
\Data\ID{1003055401}
\Updated{2020-07-15 03:07:12}
\Level{A}
\SubArea{Percent problems}
\LANGen{
\Question{
The number  \( x\) is by \( 10\% \) bigger than the number \( y \). It means that:}
{\( x-y=0.1y \)}{\( x = y +10\% \)}{\( y=\frac9{10}x \)}{\( x - 0.1x = y \)}{}{}}
\LANGpl{
\Question{
Liczba    \( x\) jest o \( 10\% \) większa od liczby  \( y \). Z tego wynika, że:}
{\( x-y=0{,}1y \)}{\( x = y +10\% \)}{\( y=\frac9{10}x \)}{\( x - 0{,}1x = y \)}{}{}}
\LANGcs{
\Question{
Číslo  \( x\) je o \( 10\,\% \) větší než číslo  \( y \). To znamená, že:}
{\( x-y=0{,}1y \)}{\( x = y +10\% \)}{\( y=\frac9{10}x \)}{\( x - 0{,}1x = y \)}{}{}}
\LANGsk{
\Question{
Číslo  \( x\) je o \( 10\% \) väčšie než číslo  \( y \). To znamená, že:}
{\( x-y=0{,}1y \)}{\( x = y +10\% \)}{\( y=\frac9{10}x \)}{\( x - 0{,}1x = y \)}{}{}}
\LANGes{
\Question{
El número   \( x\) es un \( 10\% \) más grande que el número \( y \). Esto significa que:}
{\( x-y=0.1y \)}{\( x = y +10\% \)}{\( y=\frac9{10}x \)}{\( x - 0.1x = y \)}{}{}}
\End

\Data\ID{1003055402}
\Updated{2020-08-10 07:08:25}
\Level{A}
\SubArea{Percent problems}
\LANGen{
\Question{
Mr. Nowak obtained the income of \( 20200 \) zlotys from his business activity in \( 2011 \).  He settled with the tax office, paying a tax in the amount of \( 19\% \) of his income. In \( 2011 \), Mr. Nowak's income tax amounted to:}
{\( 3838\,\text{z\l{}} \)}{\( 16362\,\text{z\l{}} \)}{\( 1063.16\,\text{z\l{}} \)}{\( 383.80\,\text{z\l{}} \)}{}{}}
\LANGpl{
\Question{
Pan Nowak uzyskał dochód  \( 20200 \) złotych z działalności gospodarczej w roku \( 2011 \).  Rozlicza się z urzędem skarbowym, płacąc podatek w wysokości  \( 19\% \) od dochodu. W \( 2011 \) roku podatek pana Nowaka z działalności gospodarczej wyniósł:}
{\( 3838\,\text{z\l{}} \)}{\( 16362\,\text{z\l{}} \)}{\( 1063{,}16\,\text{z\l{}} \)}{\( 383{,}80\,\text{z\l{}} \)}{}{}}
\LANGcs{
\Question{
Pan Novák v roce \( 2011 \) vydělal \( 20200 \) zlotých.  Podal daňové přiznání a zaplatil daň ve výši  \( 19\% \) ze svého příjmu. V roce \( 2011 \) zaplatil pan Novák na daních:}
{\( 3838\,\text{z\l{}} \)}{\( 16362\,\text{z\l{}} \)}{\( 1063{,}16\,\text{z\l{}} \)}{\( 383{,}80\,\text{z\l{}} \)}{}{}}
\LANGsk{
\Question{
Pán Novák v roku \( 2011 \) zarobil \( 20200 \) eur.  Podal daňové priznanie a zaplatil daň vo výške  \( 19\% \) zo svojho príjmu. V roku \( 2011 \) zaplatil pán Novák na daniach:}
{\( 3838\,\text{z\l{}} \)}{\( 16362\,\text{z\l{}} \)}{\( 1063{,}16\,\text{z\l{}} \)}{\( 383{,}80\,\text{z\l{}} \)}{}{}}
\LANGes{
\Question{
El señor Nowak ganó \( 20200 \) zlotys en su actividad comercial en  \( 2011 \). Presentó la declaración de la renta y pagó el impuesto sobre la renta con un porcentaje del  \( 19\% \) En \( 2011 \),  el Sr. Nowak pagó el impuesto sobre la renta por importe de:}
{\( 3838\,\text{z\l{}} \)}{\( 16362\,\text{z\l{}} \)}{\( 1063.16\,\text{z\l{}} \)}{\( 383.80\,\text{z\l{}} \)}{}{}}
\End

\Data\ID{1003055403}
\Updated{2020-07-15 03:07:12}
\Level{A}
\SubArea{Percent problems}
\LANGen{
\Question{
The price of one share of a certain firm decreased by \( 10\% \) from the initial price of \( 180 \) zlotys in one day, and the next day it increased by \( 10\% \). After these changes, the price of one share amounted to:}
{\( 178.20\,\text{z\l{}} \)}{\( 145.80\,\text{z\l{}} \)}{\( 217.80\,\text{z\l{}} \)}{\( 180\,\text{z\l{}} \)}{}{}}
\LANGpl{
\Question{
Cena jednej akcji pewnego przedsiębiorstwa w ciągu jednego dnia spadła o  \( 10\% \) z poziomu \( 180 \) złotych, a następnego dnia wzrosła o  \( 10\% \). Po tych zmianach cena jednej akcji wynosiła:}
{\( 178{,}20\,\text{z\l{}} \)}{\( 145{,}80\,\text{z\l{}} \)}{\( 217{,}80\,\text{z\l{}} \)}{\( 180\,\text{z\l{}} \)}{}{}}
\LANGcs{
\Question{
Hodnota jedné akcie jisté společnosti klesla jeden den o  \( 10\% \) z původních  \( 180 \) zlotých. Následující den vzrostla o  \( 10\% \). Jaká byla hodnota akcie po těchto změnách?}
{\( 178{,}20\,\text{z\l{}} \)}{\( 145{,}80\,\text{z\l{}} \)}{\( 217{,}80\,\text{z\l{}} \)}{\( 180\,\text{z\l{}} \)}{}{}}
\LANGsk{
\Question{
Hodnota jednej akcie istej spoločnosti klesla jeden deň o  \( 10\% \) z pôvodných  \( 180 \) zlotých. Nasledujúci deň vzrástla o  \( 10\% \). Aká bola hodnota akcie po týchto zmenách?}
{\( 178{,}20\,\text{z\l{}} \)}{\( 145{,}80\,\text{z\l{}} \)}{\( 217{,}80\,\text{z\l{}} \)}{\( 180\,\text{z\l{}} \)}{}{}}
\LANGes{
\Question{
El precio de una acción de una determinada empresa disminuyó en un  \( 10\% \) del precio inicial de  \( 180 \) zlotys en un día y al día siguiente aumentó en un  \( 10\% \). Después de estos cambios, el precio de una acción ascendió a:}
{\( 178.20\,\text{z\l{}} \)}{\( 145.80\,\text{z\l{}} \)}{\( 217.80\,\text{z\l{}} \)}{\( 180\,\text{z\l{}} \)}{}{}}
\End

\Data\ID{1003055404}
\Updated{2020-08-10 07:08:26}
\Level{A}
\SubArea{Percent problems}
\LANGen{
\Question{
\( 1100 \) people were asked to say which of the following activities (cooking, ironing, playing a musical instrument) they are bad at. \( 42\% \) of them are bad at cooking, \( 67\% \) of them are bad at ironing, \( 85\% \) of them are bad at playing a musical instrument. How many people who can cook can we expect in the group of \( 1200 \) people?}
{\( 696 \)}{\( 969 \)}{\( 504 \)}{\( 638 \)}{}{}}
\LANGpl{
\Question{
Spytano \( 1100 \) osób z jakimi umiejętnościami sobie nie radzą. Wyniki ankiety były następujące: spośród nich z gotowaniem nie radzi sobie \( 42\% \), z prasowaniem \( 67\% \), z grą na instrumencie \( 85\% \). Ilu osób potrafiących gotować należy się spodziewać w grupie \( 1200 \) osób?}
{\( 696 \)}{\( 969 \)}{\( 504 \)}{\( 638 \)}{}{}}
\LANGcs{
\Question{
\( 1100 \) lidí mělo uvést, která z následujících činností jim nejde (vaření, žehlení, hra na hudební nástroj). \( 42\% \) uvedlo vaření, \( 67\% \) uvedlo žehlení, \( 85\% \) uvedlo hru na hudební nástroj. Kolik lidí, kteří umí vařit, můžeme očekávat ve skupině  \( 1200 \) lidí?}
{\( 696 \)}{\( 969 \)}{\( 504 \)}{\( 638 \)}{}{}}
\LANGsk{
\Question{
\( 1100 \) ľudí malo uviesť, ktorá z nasledujúcich činností im nejde (varenie, žehlenie, hra na hudobný nástroj). \( 42\% \) uviedlo varenie, \( 67\% \) uviedlo žehlenie, \( 85\% \) uviedlo hru na hudobný nástroj. Koľko ľudí, ktorí vedia variť, môžeme očakávať v skupine  \( 1200 \) ľudí?}
{\( 696 \)}{\( 969 \)}{\( 504 \)}{\( 638 \)}{}{}}
\LANGes{
\Question{
\( 1100 \) personas fueron interrogadas sobre  cuál de las siguientes actividades (cocinar, planchar, tocar un instrumento musical) se les da mal. Al  \( 42\% \) de ellas se les da mal cocinar, al  \( 67\% \)de ellas se les da mal planchar, al  \( 85\% \) se les da mal tocar un instrumento musical. ¿Cuántas personas a las cuales se les da bien cocinar hay en el grupo de \( 1200 \) personas?}
{\( 696 \)}{\( 969 \)}{\( 504 \)}{\( 638 \)}{}{}}
\End

\Data\ID{1003055405}
\Updated{2020-07-15 03:07:12}
\Level{A}
\SubArea{Percent problems}
\LANGen{
\Question{
The statistical office reported that inflation amounted to \( 4.5\% \) in October and \( 4.3\% \) in September. It means that from September to October inflation has increased by:}
{\( 0.2 \) percentage points}{\( 2 \) percentage points}{\( 2\% \)}{\( 0.2\% \)}{}{}}
\LANGpl{
\Question{
Urząd statystyczny podał, że inflacja w październiku wyniosła \( 4{,}5\% \), a we wrześniu \( 4{,}3\% \). Wynika stąd, że od września do października inflacja wzrosła o:}
{\( 0{,}2 \) punktu procentowego}{\( 2 \) punktu procentowego}{\( 2\% \)}{\( 0{,}2\% \)}{}{}}
\LANGcs{
\Question{
Statistický úřad oznámil výši inflace  \( 4{,}5\,\% \) v říjnu a  \( 4{,}3\,\% \) v září. O kolik vzrostla inflace mezi zářím a říjnem?}
{o \( 0{,}2 \) procentních bodů}{o \( 2 \) procentní body}{o \( 2\% \)}{o \( 0{,}2\% \)}{}{}}
\LANGsk{
\Question{
Štatistický úrad oznámil výšku inflácie  \( 4{,}5\% \) v septembri a  \( 4{,}3\% \) v októbri. O koľko vzrástla inflácia medzi septembrom a októbrom?}
{o \( 0{,}2 \) percentných bodov}{o \( 2 \) percentné body}{o \( 2\% \)}{o \( 0{,}2\% \)}{}{}}
\LANGes{
\Question{
La oficina de estadísticas informó que la inflación ascendió a  \( 4.5\% \) en octubre y  \( 4.3\% \) en septiembre. Esto significa que de septiembre a octubre la inflación ha aumentado en:}
{\( 0.2 \) puntos porcentuales}{\( 2 \) puntos porcentuales}{\( 2\% \)}{\( 0.2\% \)}{}{}}
\End

\Data\ID{1003055406}
\Updated{2020-07-15 03:07:12}
\Level{A}
\SubArea{Percent problems}
\LANGen{
\Question{
A phone subscription in Telekomunikacja Polska is \( 35 \) zlotys. How much is the subscription fee if you add up a \( 7\% \) tax to it?}
{\( 37.45\,\text{z\l{}} \)}{about \( 38\,\text{z\l{}} \)}{\( 37\,\text{z\l{}} \)}{\( 37.50\,\text{z\l{}} \)}{}{}}
\LANGpl{
\Question{
Abonament telefoniczny w Telekomunikacji Polskiej wynosi \( 35 \) złotych. Ile wynosi opłata abonamentowa, jeśli do tej kwoty doliczony jest \( 7\% \) podatek VAT?}
{\( 37{,}45\,\text{z\l{}} \)}{około \( 38\,\text{z\l{}} \)}{\( 37\,\text{z\l{}} \)}{\( 37{,}50\,\text{z\l{}} \)}{}{}}
\LANGcs{
\Question{
Měsíční paušál za telekomunikační služby u společnosti Telekomunikacja Polska je \( 35 \) zlotých. Jaká je výše měsíčního paušálu s daní ve výši \( 7\,\% \) ?}
{\( 37{,}45\,\text{z\l{}} \)}{okolo \( 38\,\text{z\l{}} \)}{\( 37\,\text{z\l{}} \)}{\( 37{,}50\,\text{z\l{}} \)}{}{}}
\LANGsk{
\Question{
Mesačný paušál za telekomunikačné služby v spoločnosti Telekomunikacja Polska je \( 35 \) zlotých. Aká je výška mesačného paušálu s daňou vo výške \( 7\% \)?}
{\( 37{,}45\,\text{z\l{}} \)}{okolo \( 38\,\text{z\l{}} \)}{\( 37\,\text{z\l{}} \)}{\( 37{,}50\,\text{z\l{}} \)}{}{}}
\LANGes{
\Question{
Una tarifa móvil en Telekomunikacja Polska es de \( 35 \) zlotys. ¿De cuánto será la tarifa móvil si se agrega un impuesto del \( 7\% \)?}
{\( 37.45\,\text{z\l{}} \)}{aproximadamente \( 38\,\text{z\l{}} \)}{\( 37\,\text{z\l{}} \)}{\( 37.50\,\text{z\l{}} \)}{}{}}
\End

\Data\ID{1003055407}
\Updated{2022-04-23 05:04:08}
\Level{A}
\SubArea{Percent problems}
\LANGen{
\Question{
A \( 4\,\mathrm{kg} \) salt solution was prepared in the chemistry lab. The solution contains \( 250 \) g of salt. Calculate the percentage concentration of the solution.}
{\( 6.25\% \)}{\( 6\% \)}{\( 10\% \)}{\( 25\% \)}{}{}}
\LANGpl{
\Question{
W pracowni chemicznej przygotowano roztwór wodny soli o masie  \( 4\,\mathrm{kg} \). Roztwór zawiera  \( 250 \) g soli. Oblicz stężenie procentowe  roztworu.}
{\( 6{,}25\% \)}{\( 6\% \)}{\( 10\% \)}{\( 25\% \)}{}{}}
\LANGcs{
\Question{
V chemické laboratoři připravili  \( 4\,\mathrm{kg} \) solného roztoku. Roztok obsahuje  \( 250 \) g soli. O kolika procentní roztok se jedná?}
{\( 6{,}25\% \)}{\( 6\% \)}{\( 10\% \)}{\( 25\% \)}{}{}}
\LANGsk{
\Question{
V chemickom laboratóriu pripravili  \( 4\,\mathrm{kg} \) soľného roztoku. Roztok obsahuje  \( 250 \) g soli. O koľko percentný roztok sa jedná?}
{\( 6{,}25\% \)}{\( 6\% \)}{\( 10\% \)}{\( 25\% \)}{}{}}
\LANGes{
\Question{
Se preparó una solución salina de  \( 4\,\mathrm{kg} \) en el laboratorio de química. La solución contiene  \( 250 \) g  de sal. Calcula el porcentaje de concentración de la solución.}
{\( 6.25\% \)}{\( 6\% \)}{\( 10\% \)}{\( 25\% \)}{}{}}
\End

\Data\ID{1003055408}
\Updated{2020-07-15 03:07:12}
\Level{A}
\SubArea{Percent problems}
\LANGen{
\Question{
How much pure gold contains a ring weighing \( 7 \) grams if the gold fineness is \( 750 \)? (That means, the gold content is \( 75 \) percent.)}
{\( 5.25\,\mathrm{g} \)}{\( 5\,\mathrm{g} \)}{\( 6\,\mathrm{g} \)}{\( 5.5\,\mathrm{g} \)}{}{}}
\LANGpl{
\Question{
Ile jest czystego złota w obrączce próby \( 750 \) (tzn. zawartość złota to \(75\%\)), ważącej \( 7 \) g?}
{\( 5{,}25\,\mathrm{g} \)}{\( 5\,\mathrm{g} \)}{\( 6\,\mathrm{g} \)}{\( 5{,}5\,\mathrm{g} \)}{}{}}
\LANGcs{
\Question{
Kolik ryzího zlata obsahuje prsten o váze \( 7 \) gramů, pokud víme, že ryzost zlata je  \( 750 \)? (To znamená, že obsah zlata je  \( 75 \) procent.)}
{\( 5{,}25\,\mathrm{g} \)}{\( 5\,\mathrm{g} \)}{\( 6\,\mathrm{g} \)}{\( 5{,}5\,\mathrm{g} \)}{}{}}
\LANGsk{
\Question{
Koľko rýdzeho zlata obsahuje prsteň o váze \( 7 \) gramov, pokiaľ vieme, že rýdzosť zlata je  \( 750 \)? (To znamená, že obsah zlata je  \( 75 \) percent.)}
{\( 5{,}25\,\mathrm{g} \)}{\( 5\,\mathrm{g} \)}{\( 6\,\mathrm{g} \)}{\( 5{,}5\,\mathrm{g} \)}{}{}}
\LANGes{
\Question{
¿Cuánto oro puro contiene un anillo que pesa \( 7 \) gramos si la finura del oro es  \( 750 \)? (Esto significa que el contenido de oro es del  \( 75 \) por ciento).}
{\( 5.25\,\mathrm{g} \)}{\( 5\,\mathrm{g} \)}{\( 6\,\mathrm{g} \)}{\( 5.5\,\mathrm{g} \)}{}{}}
\End

\Data\ID{1003055409}
\Updated{2020-07-15 03:07:12}
\Level{A}
\SubArea{Percent problems}
\LANGen{
\Question{
In Gdańsk, a kilogram of strawberries costs \( 4 \) zlotys and in Płock it costs \( 2.50 \) zlotys. The price of strawberries in Gdansk is higher than the price of strawberries in Płock by:}
{\( 60\% \)}{\( 62.5\% \)}{\( 37.5\% \)}{\( 40\% \)}{}{}}
\LANGpl{
\Question{
W Gdańsku kilogram truskawek kosztuje \( 4 \) zł, a w Płocku  \( 2{,}50 \) zł. Cena truskawek w Gdańsku jest wyższa od ceny truskawek w Płocku o:}
{\( 60\% \)}{\( 62{,}5\% \)}{\( 37{,}5\% \)}{\( 40\% \)}{}{}}
\LANGcs{
\Question{
V Gdańsku stojí kilogram jahod \( 4 \) zloté a v Płocku pouze \( 2{,}50 \) zlotých. O kolik procent je cena jahod v Gdańsku vyšší než v Płocku?}
{\( 60\% \)}{\( 62{,}5\% \)}{\( 37{,}5\% \)}{\( 40\% \)}{}{}}
\LANGsk{
\Question{
V Košiciach stojí kilogram jahôd \( 4 \) eura a v Lučenci len \( 2{,}50 \) eura. O koľko percent je cena jahôd v Košiciach vyššia než v Lučenci?}
{\( 60\% \)}{\( 62{,}5\% \)}{\( 37{,}5\% \)}{\( 40\% \)}{}{}}
\LANGes{
\Question{
En Gdansk, un kilogramo de fresas cuesta \( 4 \) zlotys y en Płock cuesta  \( 2.50 \) zlotys. El precio de las fresas en Gdansk es más alto que el precio de las fresas en Płock en un:}
{\( 60\% \)}{\( 62.5\% \)}{\( 37.5\% \)}{\( 40\% \)}{}{}}
\End

\Data\ID{1003055410}
\Updated{2020-07-15 03:07:12}
\Level{A}
\SubArea{Percent problems}
\LANGen{
\Question{
Support for a political party X in March \( 2017 \) was \( 24\% \), and in September \( 2017 \) it was \( 30\% \). Party X noted an increase in support by:}
{\( 6 \) percentage points}{\( 20 \) percentage points}{\( 6\% \)}{\( 20\% \)}{}{}}
\LANGpl{
\Question{
Poparcie dla partii X w marcu \( 2017 \) roku wynosiło \( 24\% \), zaś we wrześniu \( 2017 \) roku było równe \( 30\% \). Partia X odnotowała wzrost poparcia o:}
{\( 6 \) punktów procentowych}{\( 20 \) punktów procentowych}{\( 6\% \)}{\( 20\% \)}{}{}}
\LANGcs{
\Question{
Politická strana X by v březnu \( 2017 \) získala \( 24\,\% \) hlasů a v září  \( 2017 \) \( 30\,\% \). Strana X tedy zaznamenala nárůst voličů o:}
{\( 6 \) procentních bodů}{\( 20 \) procentních bodů}{\( 6\,\% \)}{\( 20\,\% \)}{}{}}
\LANGsk{
\Question{
Politická strana X by v marci \( 2017 \) získala \( 24\% \) hlasov a v septembri  \( 2017 \) \( 30\% \). Strana X vtedy zaznamenala nárast voličov o:}
{\( 6 \) percentných bodov}{\( 20 \) percentných bodov}{\( 6\% \)}{\( 20\% \)}{}{}}
\LANGes{
\Question{
El apoyo a un partido político X en marzo del  \( 2017 \) fue del  \( 24\% \), y en septiembre del \( 2017 \) fue del  \( 30\% \). El aumento en el apoyo del partido X fue:}
{\( 6 \) puntos porcentuales}{\( 20 \) puntos porcentuales}{\( 6\% \)}{\( 20\% \)}{}{}}
\End

\Data\ID{1003055411}
\Updated{2020-08-10 08:08:06}
\Level{A}
\SubArea{Percent problems}
\LANGen{
\Question{
An egg weighs \( 56 \) grams. \( 55\% \) of its weight is albumen, \( 40\% \) of the egg is yolk and the rest is the shell. The weight of the shell equals:}
{\( 2.8\,\mathrm{g} \)}{\( 5.6\,\mathrm{g} \)}{\( 5\,\mathrm{g} \)}{\( 53.2\,\mathrm{g} \)}{}{}}
\LANGpl{
\Question{
Jajko waży \( 56 \) gramów. \( 55\% \) wagi jajka stanowi białko, \( 40\% \) żółtko, a resztę stanowi skorupka. Waga skorupki to:}
{\( 2{,}8\,\mathrm{g} \)}{\( 5{,}6\,\mathrm{g} \)}{\( 5\,\mathrm{g} \)}{\( 53{,}2\,\mathrm{g} \)}{}{}}
\LANGcs{
\Question{
Vajíčko váží \( 56 \) gramů. \( 55\,\% \) jeho váhy tvoří bílek, \( 40\,\% \) vajíčka tvoří žloutek a zbytek je skořápka. Skořápka tedy váží:}
{\( 2{,}8\,\mathrm{g} \)}{\( 5{,}6\,\mathrm{g} \)}{\( 5\,\mathrm{g} \)}{\( 53{,}2\,\mathrm{g} \)}{}{}}
\LANGsk{
\Question{
Vajíčko váži \( 56 \) gramov. \( 55\% \) jeho váhy tvorí bielok, \( 40\% \) vajíčka tvorí žĺtok a zbytok je škrupinka. Škrupinka teda váži:}
{\( 2{,}8\,\mathrm{g} \)}{\( 5{,}6\,\mathrm{g} \)}{\( 5\,\mathrm{g} \)}{\( 53{,}2\,\mathrm{g} \)}{}{}}
\LANGes{
\Question{
Un huevo pesa \( 56 \) gramos. El \( 55\% \) de su peso es la clara, el  \( 40\% \)  del huevo es la yema y el resto es la cáscara. El peso de la cáscara es igual a:}
{\( 2.8\,\mathrm{g} \)}{\( 5.6\,\mathrm{g} \)}{\( 5\,\mathrm{g} \)}{\( 53.2\,\mathrm{g} \)}{}{}}
\End

\Data\ID{1003055412}
\Updated{2020-07-15 03:07:12}
\Level{A}
\SubArea{Percent problems}
\LANGen{
\Question{
There are \( 40 \) students in the class, including \( 18 \) girls. What percentage of students in this class are boys?}
{\( 55\% \)}{\( 45\% \)}{\( 65\% \)}{\( 40\% \)}{}{}}
\LANGpl{
\Question{
W klasie jest  \( 40 \) uczniów, w tym \( 18 \) dziewcząt. Jaki procent uczniów tej klasy stanowią chłopcy?}
{\( 55\% \)}{\( 45\% \)}{\( 65\% \)}{\( 40\% \)}{}{}}
\LANGcs{
\Question{
Ve třídě je \( 40 \) žáků, z nichž \( 18 \) jsou dívky. Jakou část třídy tvoří chlapci?}
{\( 55\,\% \)}{\( 45\,\% \)}{\( 65\,\% \)}{\( 40\,\% \)}{}{}}
\LANGsk{
\Question{
V triede je \( 40 \) žiakov, vrátane  \( 18 \) dievčat. Akú časť triedy tvoria chlapci?}
{\( 55\% \)}{\( 45\% \)}{\( 65\% \)}{\( 40\% \)}{}{}}
\LANGes{
\Question{
Hay \( 40 \) estudiantes en una clase, de entre ellos  \( 18 \) son niñas. ¿Qué porcentaje de estudiantes en esta clase son niños?}
{\( 55\% \)}{\( 45\% \)}{\( 65\% \)}{\( 40\% \)}{}{}}
\End

\Data\ID{1003055501}
\Updated{2020-07-15 03:07:12}
\Level{A}
\SubArea{Percent problems}
\LANGen{
\Question{
The price of a certain commodity was reduced by \( 30\% \) and then by another \( 20\% \). It means that the price of the item was decreased by:}
{\( 44\% \)}{\( 50\% \)}{\( 60\% \)}{\( 56\% \)}{}{}}
\LANGpl{
\Question{
Cenę pewnego towaru obniżono o  \( 30\% \), a następnie obniżono o \( 20\% \). Zatem cenę tego towaru obniżono o:}
{\( 44\% \)}{\( 50\% \)}{\( 60\% \)}{\( 56\% \)}{}{}}
\LANGcs{
\Question{
Cena zboží klesla o  \( 30\,\% \) a následně o dalších  \( 20\,\% \). O kolik procent tedy zboží zlevnilo?}
{\( 44\,\% \)}{\( 50\,\% \)}{\( 60\,\% \)}{\( 56\,\% \)}{}{}}
\LANGsk{
\Question{
Cena tovaru klesla o  \( 30\% \) a následne o ďalších  \( 20\% \). O koľko percent bol teda tovar zľavnený?}
{\( 44\% \)}{\( 50\% \)}{\( 60\% \)}{\( 56\% \)}{}{}}
\LANGes{
\Question{
El precio de un determinado producto se redujo en un  \( 30\% \) y luego en otro  \( 20\% \). Esto significa que el precio del artículo se redujo en:}
{\( 44\% \)}{\( 50\% \)}{\( 60\% \)}{\( 56\% \)}{}{}}
\End

\Data\ID{1003055503}
\Updated{2020-07-15 03:07:12}
\Level{A}
\SubArea{Percent problems}
\LANGen{
\Question{
Tomek collected \( 320 \) coins in the piggy bank, including \( 112 \) one-zloty coins. What percentage of the coins in the  piggy bank was the number of one-zloty coins?}
{\( 35\% \)}{\( 45\% \)}{\( 40\% \)}{\( 50\% \)}{}{}}
\LANGpl{
\Question{
Tomek uzbierał w skarbonce \( 320 \) monet, w tym \( 112 \) jednozłotówek. Jaki procent liczby monet w skarbonce stanowiła liczba monet jednozłotowych?}
{\( 35\% \)}{\( 45\% \)}{\( 40\% \)}{\( 50\% \)}{}{}}
\LANGcs{
\Question{
Tomáš nastřádal do prasátka \( 320 \) v drobných, včetně \( 112 \) mincí v hodnotě jeden zlotý. Jak velkou část všech mincí v prasátku tvoří mince v hodnotě jeden zlotý?}
{\( 35\% \)}{\( 45\% \)}{\( 40\% \)}{\( 50\% \)}{}{}}
\LANGsk{
\Question{
Tomáš nastrkal do prasiatka \( 320 \) v drobných, vrátane \( 112 \) mincí v hodnote jedno euro. Akú veľkú časť všetkých mincí v prasiatku tvoria mince v hodnote jedno euro?}
{\( 35\% \)}{\( 45\% \)}{\( 40\% \)}{\( 50\% \)}{}{}}
\LANGes{
\Question{
Tomek ahorró \( 320 \) monedas en su hucha incluyendo \( 112 \) monedas de un zloty. ¿Qué porcentaje de las monedas de su hucha eran de monedas de un zloty?}
{\( 35\% \)}{\( 45\% \)}{\( 40\% \)}{\( 50\% \)}{}{}}
\End

\Data\ID{1003055504}
\Updated{2020-07-15 03:07:12}
\Level{A}
\SubArea{Percent problems}
\LANGen{
\Question{
If \( 12\% \) of the number \( a \) equals \( 15 \), then}
{\( a = 125 \).}{\( a = 75 \).}{\( a = 80 \).}{\( a = 150 \).}{}{}}
\LANGpl{
\Question{
Jeżeli \( 12\% \) liczby  \( a \) jest równa \( 15 \), to}
{\( a = 125 \).}{\( a = 75 \).}{\( a = 80 \).}{\( a = 150 \).}{}{}}
\LANGcs{
\Question{
Pokud se \( 12\,\% \) daného čísla \( a \) rovná \( 15 \), potom}
{\( a = 125 \).}{\( a = 75 \).}{\( a = 80 \).}{\( a = 150 \).}{}{}}
\LANGsk{
\Question{
Ak sa \( 12\% \) daného čísla \( a \) rovná \( 15 \), potom}
{\( a = 125 \).}{\( a = 75 \).}{\( a = 80 \).}{\( a = 150 \).}{}{}}
\LANGes{
\Question{
Si el \( 12\% \) de un número  \( a \) es igual a \( 15 \), entonces}
{\( a = 125 \).}{\( a = 75 \).}{\( a = 80 \).}{\( a = 150 \).}{}{}}
\End

\Data\ID{1003055506}
\Updated{2020-07-15 03:07:12}
\Level{A}
\SubArea{Percent problems}
\LANGen{
\Question{
A notebook and a pen cost the same amount of money. The price of the notebook was increased by \( 3\% \), and the price of the pen was increased by \( 5\% \). For the set consisting of \( 5 \) notebooks and \( 5 \) pens you need to pay more than before the rise}
{by \( 4\% \).}{by \( 8\% \).}{by \( 5\% \).}{by \( 15\% \).}{}{}}
\LANGpl{
\Question{
Zeszyt i długopis kosztowały tyle samo. Cenę zeszytu podniesiono o \( 3\% \), a cenę długopisu o \( 5\% \). Za zestaw złożony z \( 5 \) zeszytów i \( 5 \)  długopisów trzeba zapłacić więcej niż przed podwyżką}
{o \( 4\% \).}{o \( 8\% \).}{o \( 5\% \).}{o \( 15\% \).}{}{}}
\LANGcs{
\Question{
Blok a pero mají stejnou cenu. Blok zdražil o \( 3\,\% \) a cena pera vzrostla o  \( 5\,\% \). Za sadu, která se skládá z  \( 5 \) bloků a  \( 5 \) per tedy musíš zaplatit více}
{o \( 4\,\% \).}{o \( 8\,\% \).}{o \( 5\,\% \).}{o \( 15\,\% \).}{}{}}
\LANGsk{
\Question{
Blok a pero majú rovnakú cenu. Blok zdražel o \( 3\% \) a cena pera vzrástla o  \( 5\% \). Za sadu, ktorá sa skladá z  \( 5 \) blokov a  \( 5 \) pier teda musíš zaplatiť viac}
{o \( 4\% \).}{o \( 8\% \).}{o \( 5\% \).}{o \( 15\% \).}{}{}}
\LANGes{
\Question{
Un cuaderno y un bolígrafo cuestan la misma cantidad de dinero. El precio del cuaderno aumentó en un  \( 3\% \), y el precio del bolígrafo aumentó en un  \( 5\% \). Para el conjunto que consta de \( 5 \) cuadernos y  \( 5 \) bolígrafos, debes pagar más}
{en un \( 4\% \).}{en un \( 8\% \).}{en un \( 5\% \).}{en un \( 15\% \).}{}{}}
\End

\Data\ID{1003055507}
\Updated{2020-07-15 03:07:12}
\Level{A}
\SubArea{Percent problems}
\LANGen{
\Question{
The number \( a \) is by \( 30\% \) bigger than the number \( b \). It means that  \( b \) is smaller than  \( a \)}
{by \( 23\frac1{13}\% \).}{by \( 30\% \).}{by \( 25\% \).}{by \( 70\% \).}{}{}}
\LANGpl{
\Question{
Liczba  \( a \)  jest o \( 30\% \)  większa od liczby \( b \). Oznacza to, że liczba \( b \) jest mniejsza od liczby \( a \)}
{o \( 23\frac1{13}\% \).}{o \( 30\% \).}{o \( 25\% \).}{o \( 70\% \).}{}{}}
\LANGcs{
\Question{
Číslo \( a \) je o  \( 30\% \) větší než číslo  \( b \). To znamená, že  \( b \) je menší než  \( a \)}
{o \( 23\frac1{13}\% \).}{o \( 30\% \).}{o \( 25\% \).}{o \( 70\% \).}{}{}}
\LANGsk{
\Question{
Číslo \( a \) je o  \( 30\% \) väčšie než číslo  \( b \). To znamená, že  \( b \) je menšie než  \( a \)}
{o \( 23\frac1{13}\% \).}{o \( 30\% \).}{o \( 25\% \).}{o \( 70\% \).}{}{}}
\LANGes{
\Question{
El número \( a \) es un \( 30\% \) más grande que el número  \( b \). Esto significa que  \( b \) es menor que   \( a \)}
{en un \( 23\frac1{13}\% \).}{en un  \( 30\% \).}{en un \( 25\% \).}{en un \( 70\% \).}{}{}}
\End

\Data\ID{1003055508}
\Updated{2020-07-15 03:07:12}
\Level{A}
\SubArea{Percent problems}
\LANGen{
\Question{
In a chess tournament \( 70\% \) of all participants were boys. Girls were \( 20\% \) of the number of boys. Other participants of the event were coaches and judges. Find the percentage of girls with the reference to all the participants of the tournament.}
{\( 14\% \) of all participants}{\( 30\% \) of all participants}{\( 50\% \) of all participants}{Not enough information to find a solution.}{}{}}
\LANGpl{
\Question{
W turnieju szachowym chłopcy stanowili \( 70\% \)  wszystkich uczestników, a dziewczęta — \( 20\% \) liczby chłopców. Pozostali uczestnicy to trenerzy i sędziowie. Oceń, jaki procent liczby wszystkich uczestników turnieju stanowiły dziewczęta.}
{\( 14\% \) liczby uczestników}{\( 30\% \) liczby uczestników}{\( 50\% \) liczby uczestników}{Za mało danych do rozwiązania problemu.}{}{}}
\LANGcs{
\Question{
Na šachovém turnaji je mezi účastníky \( 70\,\% \) chlapců. Dívky tvoří \( 20\,\% \) z celkového počtu chlapců. Ostatními účastníky jsou trenéři a rozhodčí. Jaké procento z celkového počtu účastníků turnaje tvoří dívky?}
{\( 14\,\% \)}{\( 30\,\% \)}{\( 50\,\% \)}{Pro vyřešení úlohy není zadán dostatek informací.}{}{}}
\LANGsk{
\Question{
Na šachovom turnaji je \( 70\% \) účastníkov chlapcov. Dievčatá tvoria \( 20\% \) z celkového počtu chlapcov. Ostatnými účastníkmi sú tréneri a rozhodcovia. Aké percento z celkového počtu účastníkov turnaja tvoria dievčatá?}
{\( 14\% \) všetkých účastníkov}{\( 30\% \) všetkých účastníkov}{\( 50\% \) všetkých účastníkov}{Nie je dosť informácií pre výpočet.}{}{}}
\LANGes{
\Question{
En un torneo de ajedrez, el  \( 70\% \) de todos los participantes eran niños. Las niñas representaron el \( 20\% \) del número de niños. Otros participantes del evento eran entrenadores y jueces. Encuentra el porcentaje de chicas en relación a todos los participantes del torneo.}
{\( 14\% \) de todos los participantes}{\( 30\% \)  de todos los participantes}{\( 50\% \)  de todos los participantes}{No hay suficiente información para encontrar la solución.}{}{}}
\End

\Data\ID{1003055509}
\Updated{2020-10-17 04:10:44}
\Level{A}
\SubArea{Percent problems}
\LANGen{
\Question{
Mr. Nowak deposited \( 2000 \) zlotys into his savings account at an annual interest rate of  \( 4.5\% \), compounded annually. At the end of the year, the \( 19\% \) tax of the interest was deducted. What was the amount of the money deducted?}
{\( 17.10 \) zlotys}{\( 19 \) zlotys}{\( 38 \) zlotys}{\( 85.50 \) zlotys}{}{}}
\LANGpl{
\Question{
Pan Nowak wpłacił do banku \( 2000 \) zł na roczną lokatę oprocentowaną na  \( 4{,}5\% \) w skali roku. Po upływie roku odliczono mu \( 19\% \) podatku od odsetek. Odliczona kwota wynosi}
{\( 17{,}10 \) zł.}{\( 19 \) zł.}{\( 38 \) zł.}{\( 85{,}50 \) zł.}{}{}}
\LANGcs{
\Question{
Pan Novák uložil \( 2000 \) zlotých na spořící účet s roční úrokovou sazbou  \( 4{,}5\% \). Na konci roku mu z úroků odečetli daň ve výši  \( 19\% \). Kolik peněz zaplatil na dani z úroků?}
{\( 17{,}10 \) zlotých}{\( 19 \) zlotých}{\( 38 \) zlotých}{\( 85{,}50 \) zlotých}{}{}}
\LANGsk{
\Question{
Pán Novák uložil \( 2000 \) EUR na sporiaci účet s ročnou úrokovou sadzbou  \( 4{,}5\% \). Na konci roka mu z úrokov odčítali daň vo výške  \( 19\% \). Koľko peňazí zaplatil na dani z úrokov?}
{\( 17{,}10 \) EUR}{\( 19 \) EUR}{\( 38 \) EUR}{\( 85{,}50 \) EUR}{}{}}
\LANGes{
\Question{
El Sr. Nowak depositó  \( 2000 \) zlotys en su cuenta de ahorros a una tasa de interés anual de  \( 4.5\% \), compuesto. Al final del año, el impuesto del  \( 19\% \) de los intereses se dedujo. ¿Cuál fue la cantidad del dinero deducido?}
{\( 17.10 \) zlotys}{\( 19 \) zlotys}{\( 38 \) zlotys}{\( 85.50 \) zlotys}{}{}}
\End

\Data\ID{1003055510}
\Updated{2022-04-23 05:04:08}
\Level{A}
\SubArea{Percent problems}
\LANGen{
\Question{
A pair of trousers costs \( 126 \) zlotys after a \( 30\% \) discount. How much did the trousers cost before the discount?}
{\( 180 \) zlotys}{\( 163.80 \) zlotys}{\( 294 \) zlotys}{\( 420 \) zlotys}{}{}}
\LANGpl{
\Question{
Spodnie po obniżce ceny o \( 30\% \) kosztują \( 126 \) zł. Ile kosztowały spodnie przed obniżką?}
{\( 180 \) zł}{\( 163{,}80 \) zł}{\( 294 \) zł}{\( 420 \) zł}{}{}}
\LANGcs{
\Question{
Kalhoty stojí \( 126 \) zlotých po slevnění o \( 30\% \). Kolik stály kalhoty před slevou?}
{\( 180 \) zlotých}{\( 163{,}80 \) zlotých}{\( 294 \) zlotých}{\( 420 \) zlotých}{}{}}
\LANGsk{
\Question{
Nohavice stoja \( 126 \) EUR po odpočítaní zľavy vo výške \( 30\% \). Koľko stály nohavice pred zľavou?}
{\( 180 \) EUR}{\( 163{,}80 \) EUR}{\( 294 \) EUR}{\( 420 \) EUR}{}{}}
\LANGes{
\Question{
Unos pantalones cuestan  \( 126 \) zlotys después de un descuento del  \( 30\% \) ¿Cuánto costaban los pantalones antes del descuento?}
{\( 180 \) zlotys}{\( 163.80 \) zlotys}{\( 294 \) zlotys}{\( 420 \) zlotys}{}{}}
\End

\Data\ID{1003055511}
\Updated{2020-07-15 03:07:12}
\Level{A}
\SubArea{Percent problems}
\LANGen{
\Question{
The price of an item without the tax equals \(60\) zlotys. When we add \(22\%\) tax to the price, the item costs}
{\( 73.20 \) zlotys}{\( 49.18 \) zlotys}{\( 60.22 \) zlotys}{\( 82 \) zlotys}{}{}}
\LANGpl{
\Question{
Cena towaru bez podatku VAT jest równa 60 zł. Towar ten wraz z podatkiem VAT w wysokości $22\%$ kosztuje}
{\( 73{,}20 \) zł.}{\( 49{,}18 \) zł.}{\( 60{,}22 \) zł.}{\( 82 \) zł.}{}{}}
\LANGcs{
\Question{
Jednotková cena bez daně je \(60\) zlotých. Když k ceně přičteme daň ve výši \(22\%\), bude cena}
{\( 73{,}20 \) zlotých}{\( 49{,}18 \) zlotých}{\( 60{,}22 \) zlotých}{\( 82 \) zlotých}{}{}}
\LANGsk{
\Question{
Jednotková cena bez dane je \(60\) EUR. Keď k cene pripočítame daň vo výške \(22\%\), bude cena}
{\( 73{,}20 \) EUR}{\( 49{,}18 \) EUR}{\( 60{,}22 \) EUR}{\( 82 \) EUR}{}{}}
\LANGes{
\Question{
El precio de un artículo sin impuestos es igual a  \(60\)zlotys. Cuando agregamos  \(22\%\) de impuestos al precio, el artículo cuesta}
{\( 73.20 \) zlotys}{\( 49.18 \) zlotys}{\( 60.22 \) zlotys}{\( 82 \) zlotys}{}{}}
\End

\Data\ID{1003055512}
\Updated{2020-07-15 03:07:12}
\Level{A}
\SubArea{Percent problems}
\LANGen{
\Question{
The price of a certain commodity increased from \( 1200 \) zlotys to \( 1248 \) zlotys. By what percentage the price has risen?}
{\( 4\% \)}{\( 40\% \)}{\( 0.4\% \)}{\( 0.04\% \)}{}{}}
\LANGpl{
\Question{
Cena towaru wzrosła z \( 1200 \) zł do \( 1248 \) zł. O jaki procent wzrosła cena?}
{\( 4\% \)}{\( 40\% \)}{\( 0{,}4\% \)}{\( 0{,}04\% \)}{}{}}
\LANGcs{
\Question{
Cena zboží vzrostla z \( 1200 \) zlotých na \( 1248 \) zlotých. O kolik procent zboží zdražilo?}
{\( 4\,\% \)}{\( 40\,\% \)}{\( 0{,}4\,\% \)}{\( 0{,}04\,\% \)}{}{}}
\LANGsk{
\Question{
Cena tovaru vzrástla z \( 1200 \) zlotých na \( 1248 \) zlotých. O koľko percent tovar zdražel?}
{\( 4\% \)}{\( 40\% \)}{\( 0{,}4\% \)}{\( 0{,}04\% \)}{}{}}
\LANGes{
\Question{
El precio de un determinado producto aumentó de \( 1200 \) zlotys a  \( 1248 \) zlotys. ¿En qué porcentaje ha aumentado el precio?}
{\( 4\% \)}{\( 40\% \)}{\( 0.4\% \)}{\( 0.04\% \)}{}{}}
\End

\Data\ID{1003055513}
\Updated{2020-08-10 08:08:00}
\Level{A}
\SubArea{Percent problems}
\LANGen{
\Question{
\( 280 \) tickets were sold for the show, including \( 126 \) reduced tickets. What percentage of tickets sold were reduced tickets?}
{\( 45\% \)}{\( 33\% \)}{\( 22\% \)}{\( 63\% \)}{}{}}
\LANGpl{
\Question{
Na seans filmowy sprzedano \( 280 \) biletów, w tym  \( 126 \) ulgowych. Jaki procent sprzedanych biletów stanowiły bilety ulgowe?}
{\( 45\% \)}{\( 33\% \)}{\( 22\% \)}{\( 63\% \)}{}{}}
\LANGcs{
\Question{
Na představení se prodalo \( 280 \) lístků, včetně \( 126 \) lístků za zvýhodněnou cenu. Jaké procento prodaných lístků tvořily ty zlevněné?}
{\( 45\% \)}{\( 33\% \)}{\( 22\% \)}{\( 63\% \)}{}{}}
\LANGsk{
\Question{
Na predstavenie sa predalo \( 280 \) lístkov, vrátane \( 126 \) lístkov za zvýhodnenú cenu. Aké percento predaných lístkov tvorili tie zľavnené?}
{\( 45\% \)}{\( 33\% \)}{\( 22\% \)}{\( 63\% \)}{}{}}
\LANGes{
\Question{
Se vendieron \( 280 \) entradas para un espectáculo, incluidas \( 126 \) entradas reducidas. ¿Qué porcentaje de las entradas vendidas eran  reducidas?}
{\( 45\% \)}{\( 33\% \)}{\( 22\% \)}{\( 63\% \)}{}{}}
\End

\Data\ID{1103055502}
\Updated{2020-08-22 04:08:46}
\Level{A}
\SubArea{Percent problems}
\IMG{1103055501.pdf}
\LANGen{
\Question{
Titanic sailed to New York on April 10, 1912. Among the \( 2200 \) people on board there were Class I, Class II, Class III passengers and the crew. The pie chart shows the Titanic percent composition (to \(1\%\) accuracy). By what percentage is the number of the third-class passengers bigger than the number of the crew?}
{\(25\%\)}{\(8\%\)}{\(17\%\)}{\(125\%\)}{}{}}
\LANGpl{
\Question{
Titanic wypłynął do Nowego Jorku 10 kwietnia 1912 roku. Wśród  \( 2200 \) osób znajdujących się na pokładzie byli pasażerowie podróżujący I, II i III klasą oraz załoga. Diagram kołowy pokazuje procentowy skład osobowy Titanica (z dokładnością do \(1\%\). O ile procent liczba podróżujących III klasą była większa od liczby członków załogi? (Legenda: Class III passengers = pasażerowie klasy III, Crew = załoga, Class II passengers = pasażerowie klasy II, Class I passengers = pasażerowie klasy I)}
{\(25\%\)}{\(8\%\)}{\(17\%\)}{\(125\%\)}{}{}}
\LANGcs{
\Question{
Titanik vyplul do New Yorku 10. dubna 1912. Mezi \( 2200 \) lidmi na palubě byli pasažéři první, druhé i třetí třídy a posádka. Koláčový graf zobrazuje poměrné zastoupení jednotlivých skupin (s přesností na \(1\%\)). O kolik procent bylo na palubě více pasažérů třetí třídy než posádky? (Slovníček: Class III passengers - pasažéři 3. třídy, Crew - posádka, Class II passengers - pasažéři 2. třídy, Class I passengers - pasažéři 1. třídy)}
{\(25\%\)}{\(8\%\)}{\(17\%\)}{\(125\%\)}{}{}}
\LANGsk{
\Question{
Titanik vyplával do New Yorku 10. apríla 1912. Medzi \( 2200 \) ľuďmi na palube boli pasažieri prvej, druhej i tretej triedy a posádka. Koláčový graf zobrazuje pomerové zastúpenie jednotlivých skupín (s presnosťou na \(1\%\)). O koľko percent bolo na palube viac pasažierov tretej triedy než posádky? (Slovníček: Class III passengers - pasažieri 3. triedy, Crew - posádka, Class II passengers - pasažieri 2. triedy, Class I passengers - pasažieri 1. triedy)}
{\(25\%\)}{\(8\%\)}{\(17\%\)}{\(125\%\)}{}{}}
\LANGes{
\Question{
El Titanic salió hacia Nueva York el 10 de abril de 1912. Entre las  \( 2200 \) personas de a bordo había pasajeros de Clase I, de Clase II, de Clase III y la tripulación. El diagrama de sectores muestra la composición porcentual del Titanic (con una precisión del  \(1\%\)). ¿En qué porcentaje es mayor el número de pasajeros de tercera clase que el número de la tripulación?}
{\(25\%\)}{\(8\%\)}{\(17\%\)}{\(125\%\)}{}{}}
\End

\Data\ID{1103055505}
\Updated{2020-08-22 04:08:19}
\Level{A}
\SubArea{Percent problems}
\IMG{1103055502_0.pdf}
\LANGen{
\Question{
In the warehouse there were \( 400 \) kilograms of fruit. The diagram shows the percentage composition of fruit in the warehouse according to their weight. How many kilograms of pears were there?}
{\( 140\,\mathrm{kg} \)}{\( 150\,\mathrm{kg} \)}{\( 145\,\mathrm{kg} \)}{\( 155\,\mathrm{kg} \)}{}{}}
\LANGpl{
\Question{
W magazynie było \( 400 \) kg owoców. Na diagramie obok przedstawiono procentowy udział poszczególnych owoców według wagi. Ile było kg gruszek? (Legenda: pears = gruszki, plums = śliwki, apples = jabłka, cherries = czereśnie)}
{\( 140\,\mathrm{kg} \)}{\( 150\,\mathrm{kg} \)}{\( 145\,\mathrm{kg} \)}{\( 155\,\mathrm{kg} \)}{}{}}
\LANGcs{
\Question{
Ve skladu se nacházelo \( 400 \) kilogramů ovoce. Graf znázorňuje podíl druhů ovoce podle váhy. Kolik kilogramů hrušek bylo ve skladu? (Slovníček: plums - švestky, cherries - třešně, apples - jablka, pears - hrušky)}
{\( 140\,\mathrm{kg} \)}{\( 150\,\mathrm{kg} \)}{\( 145\,\mathrm{kg} \)}{\( 155\,\mathrm{kg} \)}{}{}}
\LANGsk{
\Question{
V sklade sa nachádzalo \( 400 \) kilogramov ovocia. Graf znázorňuje podiel druhov ovocia podľa váhy. Koľko kilogramov hrušiek bolo v sklade? (Slovníček: plums - slivky, cherries - čerešne, apples - jablká, pears - hrušky)}
{\( 140\,\mathrm{kg} \)}{\( 150\,\mathrm{kg} \)}{\( 145\,\mathrm{kg} \)}{\( 155\,\mathrm{kg} \)}{}{}}
\LANGes{
\Question{
En el almacén había \( 400 \) kilogramos de fruta. El diagrama muestra la composición porcentual de fruta en el almacén según su peso. ¿Cuántos kilogramos de peras había?}
{\( 140\,\mathrm{kg} \)}{\( 150\,\mathrm{kg} \)}{\( 145\,\mathrm{kg} \)}{\( 155\,\mathrm{kg} \)}{}{}}
\End

\Data\ID{2010001901}
\Updated{2022-11-01 09:11:29}
\Level{A}
\SubArea{Percent problems}
\LANGen{
\Question{
A \( 5\,\mathrm{kg} \) sugar solution was prepared by dissolving \( 400 \) g of sugar in water. Calculate the percentage concentration of the solution.}
{\( 8\% \)}{\( 12.5\% \)}{\( 80\% \)}{\( 16\% \)}{}{}}
\LANGpl{
\Question{
Przygotowano \( 5\,\mathrm{kg} \) roztworu cukru przez rozpuszczenie \( 400 \) g cukru w wodzie. Obliczyć procentowe stężenie roztworu.}
{\( 8\% \)}{\( 12{,}5\% \)}{\( 80\% \)}{\( 16\% \)}{}{}}
\LANGcs{
\Question{
Rozpuštěním \( 400 \) g cukru ve vodě vzniklo \( 5\,\mathrm{kg} \) cukerného roztoku. O kolika procentní roztok se jedná?}
{\( 8\% \)}{\( 12{,}5\% \)}{\( 80\% \)}{\( 16\% \)}{}{}}
\LANGsk{
\Question{
Rozpustením \( 400 \)  g cukru vo vode vzniklo \( 5\,\mathrm{kg} \) bezfarebného roztoku cukru.  Koľko percentný roztok vznikol?}
{\( 8\% \)}{\( 12{,}5\% \)}{\( 80\% \)}{\( 16\% \)}{}{}}
\LANGes{
\Question{
Se preparó una disolución de \(5\, \mathrm {kg}\) de azúcar disolviendo \(400\) g de azúcar en agua. Calcula la concentración porcentual de la disolución.}
{\( 8\% \)}{\( 12.5\% \)}{\( 80\% \)}{\( 16\% \)}{}{}}
\End

\Data\ID{2010001902}
\Updated{2022-02-20 08:02:38}
\Level{A}
\SubArea{Percent problems}
\LANGen{
\Question{
Sport shoes cost \( 112 \)  EUR after a \( 40\% \) price increase. How much did the same shoes cost before the price increase?}
{\( 80 \) EUR}{\( 28 \) EUR}{\( 78.4 \) EUR}{\( 67.2 \) EUR}{}{}}
\LANGpl{
\Question{
Buty kosztują \( 112 \)  EUR po \( 40\% \) wzroście ceny. Ile kosztowały te same buty przed wzrostem ceny?}
{\( 80 \) EUR}{\( 28 \) EUR}{\( 78{,}4 \) EUR}{\( 67{,}2 \) EUR}{}{}}
\LANGcs{
\Question{
Po zdražení o \( 40\% \) stály sportovní boty \( 112 \)  EUR. Kolik stály tytéž boty před zdražením?}
{\( 80 \) EUR}{\( 28 \) EUR}{\( 78{,}4 \) EUR}{\( 67{,}2 \) EUR}{}{}}
\LANGsk{
\Question{
Po zdražení o \( 40\% \) stáli športové topánky \( 112 \) EUR. Koľko stáli tie isté topánky pred zdražením?}
{\( 80 \) EUR}{\( 28 \) EUR}{\( 78{,}4 \) EUR}{\( 67{,}2 \) EUR}{}{}}
\LANGes{
\Question{
Unos  zapatos deportivos cuestan \(112\) EUR después de un aumento de precio del \(40\% \). ¿Cuánto costaban los mismos zapatos antes del aumento de precio?}
{\( 80 \) EUR}{\( 28 \) EUR}{\( 78.4 \) EUR}{\( 67.2 \) EUR}{}{}}
\End

\Data\ID{2010008301}
\Updated{2022-05-16 11:05:54}
\Level{A}
\SubArea{Percent problems}
\LANGen{
\Question{
Three companies made a commitment to donate a portion of their annual profit to charity (see the table). Which company did donate the biggest sum?

\[\begin{array}{|c|c|c|} \hline \text{Company} & \text{Annual profit} & \text{Ratio from annual profit} \\\hline A & 6\,000\,000 & 4.0\% \\\hline B & 12\,000\,000 & 2.0\% \\\hline C & 48\,000\,000 & 0.5\% \\\hline \end{array}\]}
{All companies donated the same sum.}{The biggest donor was the company A.}{The biggest donor was the company B.}{The biggest donor was the company C.}{}{}}
\LANGpl{
\Question{
Trzy firmy zobowiązały się przekazać część swojego rocznego zysku na cele charytatywne (patrz tabela). Która firma przekazała największą sumę?

\[\begin{array}{|c|c|c|} \hline \text{Firma} & \text{Roczny zysk} & \text{Udział w zysku} \\\hline A & 6\,000\,000 & 4{,}0\% \\\hline B & 12\,000\,000 & 2{,}0\% \\\hline C & 48\,000\,000 & 0{,}5\% \\\hline \end{array}\]}
{Wszystkie firmy przekazały taką samą sumę.}{Największą sumę przekazała firma A.}{Największą sumę przekazała firma B.}{Największą sumę przekazała firma C.}{}{}}
\LANGcs{
\Question{
Tři firmy se zavázaly věnovat daný podíl svého ročního zisku na humanitární účely (viz tabulka). Která firma věnovala největší částku?

\[\begin{array}{|c|c|c|} \hline \text{Firma} & \text{Roční obrat} & \text{Podíl ze zisku} \\\hline A & 6\,000\,000 & 4{,}0\,\% \\\hline B & 12\,000\,000 & 2{,}0\,\% \\\hline C & 48\,000\,000 & 0{,}5\,\% \\\hline \end{array}\]}
{Všechny firmy věnovaly stejnou částku.}{Největším dárcem se stala firma A.}{Největším dárcem se stala firma B.}{Největším dárcem se stala firma C.}{}{}}
\LANGsk{
\Question{
Tri firmy sa zaviazali venovať daný podiel svojho ročného zisku na humanitárnu pomoc (pozri tabuľku). Ktorá firma venovala najväčšiu sumu?

\[\begin{array}{|c|c|c|} \hline \text{Firma} & \text{Ročný obrat} & \text{Podiel zo zisku} \\\hline A & 6\,000\,000 & 4{,}0\% \\\hline B & 12\,000\,000 & 2{,}0\% \\\hline C & 48\,000\,000 & 0{,}5\% \\\hline \end{array}\]}
{Všetky firmy venovali rovnakú sumu.}{Najväčšiu sumu darovala firma A.}{Najväčšiu sumu darovala firma B.}{Najväčšiu sumu darovala firma C.}{}{}}
\LANGes{
\Question{
Tres empresas se comprometieron a donar una parte de sus ganancias anuales a la caridad (mira la tabla). ¿Qué empresa donó la mayor suma?

\[\begin{array}{|c|c|c|} \hline \text{Empresa} & \text{Ganancias anuales} & \text{Participación en las ganancias} \\\hline A & 6\,000\,000 & 4.0\% \\\hline B & 12\,000\,000 & 2.0\% \\\hline C & 48\,000\,000 & 0.5\% \\\hline \end{array}\]}
{Todas las empresas donaron la misma suma.}{El mayor donante fue la empresa A.}{El mayor donante fue la empresa B.}{El mayor donante fue la empresa C.}{}{}}
\End

\Data\ID{2010008302}
\Updated{2022-07-20 08:07:41}
\Level{A}
\SubArea{Percent problems}
\LANGen{
\Question{
During epidemic the towns reported the ratios of ill citizens (see the table). Which town had the biggest number of ill citizens on 1$^{\mathrm{st}}$ July.

\[\begin{array}{|c|c|c|} \hline \text{Town} & \text{Total population} & \text{Ratio of ill citizens on 1\(^{\mathrm{st}}\) July 2020} \\\hline A & 8\,000 & 4.0\% \\\hline B & 64\,000 & 0.5\% \\\hline C & 320\,000 & 0.1\% \\\hline \end{array}\]}
{All towns had the same number of ill citizens.}{The biggest number of ill citizens was in the town A.}{The biggest number of ill citizens was in the town B.}{The biggest number of ill citizens was in the town C.}{}{}}
\LANGpl{
\Question{
W czasie epidemii miasta podawały wskaźniki wszystkich chorych (patrz tabela). Które miasto 1 lipca miało najwięcej chorych?

\[\begin{array}{|c|c|c|} \hline \text{Miasto} & \text{Ludność} & \text{Wskaźnik chorych w dniu 1 lipca 2020} \\\hline A & 8\,000 & 4{,}0\% \\\hline B & 64\,000 & 0{,}5\% \\\hline C & 320\,000 & 0{,}1\% \\\hline \end{array}\]}
{Wszystkie miasta miały taką samą liczbę chorych.}{Najwięcej chorych było w mieście A.}{Najwięcej chorych było w miejscowości B.}{Najwięcej chorych było w mieście C.}{}{}}
\LANGcs{
\Question{
V období epidemie hlásily obce podíl nemocných obyvatel  (viz tabulka). Ve které obci byl 1. července největší počet nemocných?

\[\begin{array}{|c|c|c|} \hline \text{Obec} & \text{Celkový počet obyvatel} & \text{Podíl nemocných k 1. červenci 2020} \\\hline A & 8\,000 & 4{,}0\,\% \\\hline B & 64\,000 & 0{,}5\,\% \\\hline C & 320\,000 & 0{,}1\,\% \\\hline \end{array}\]}
{Ve všech obcích byl stejný počet nemocných.}{Nejvíce nemocných bylo v obci A.}{Nejvíce nemocných bylo v obci B.}{Nejvíce nemocných bylo v obci C.}{}{}}
\LANGsk{
\Question{
V období epidémie nahlasovali obce podiel chorých obyvateľov  (pozri tabuľku). V ktorej obci bol 1. júla najväčší počet chorých?

\[\begin{array}{|c|c|c|} \hline \text{Obec} & \text{Celkový počet obyvateľov} & \text{Podiel chorých k 1. júlu 2020} \\\hline A & 8\,000 & 4{,}0\% \\\hline B & 64\,000 & 0{,}5\% \\\hline C & 320\,000 & 0{,}1\% \\\hline \end{array}\]}
{Vo všetkých obciach bol rovnaký počet chorých.}{Najviac chorých bolo v obci A.}{Najviac chorých bolo v obci B.}{Najviac chorých bolo v obci C.}{}{}}
\LANGes{
\Question{
Durante una epidemia tres pueblos informaron sobre la proporción de ciudadanos enfermos (observa la tabla). ¿Qué ciudad tuvo el mayor número de ciudadanos enfermos el 1 de julio?

\[\begin{array}{|c|c|c|} \hline \text{Ciudad} & \text{Populación total} & \text{Proporción de los enfermos el 1 de julio 2020} \\\hline A & 8\,000 & 4.0\% \\\hline B & 64\,000 & 0.5\% \\\hline C & 320\,000 & 0.1\% \\\hline \end{array}\]}
{Todos los pueblos tenían el mismo número de ciudadanos enfermos.}{El mayor número de ciudadanos enfermos se encontraba en el pueblo A.}{El mayor número de ciudadanos enfermos se encontraba en el pueblo B.}{El mayor número de ciudadanos enfermos se encontraba en el pueblo C.}{}{}}
\End

\Data\ID{2010008303}
\Updated{2022-07-21 05:07:58}
\Level{A}
\SubArea{Percent problems}
\LANGen{
\Question{
Three merchants were selling the same product for the same price. Successively they adjusted the prices as follows:

(a) Merchant A first increased the price of the product by \(5\%\), but later reduced it by \(10\%\). 

(b) Merchant B first reduced the price of the product by \(2\%\) and later reduced it again, this time by \(3\%\). 

(c) Merchant C adjusted the price of the product only once. He cheapened the product by \(10\%\).


Which merchant offers the product for the highest price after price adjustments?}
{Merchant B}{Merchant A}{Merchant C}{All merchants will sell the product for the same price.}{}{}}
\LANGpl{
\Question{
Trzech sprzedawców sprzedawało ten sam produkt za tę samą cenę. Sukcesywnie dostosowywali ceny w następujący sposób:

(a)  Sprzedawca A najpierw podniósł cenę produktu o \(5\%\), ale później obniżył ją o \(10\%\). 

(b) Sprzedawca B najpierw obniżył cenę produktu o \(2\%\)  a później obniżył ją ponownie, tym razem o \(3\%\). 

(c) Sprzedawca C skorygował cenę produktu tylko raz. Obniżył cenę produktu o \(10\%\).

Który sprzedawca oferuje produkt za najwyższą cenę po korektach cen?}
{Sprzedawca B}{Sprzedawca A}{Sprzedawca C}{Wszyscy sprzedawcy sprzedają produkt za tę samą cenę.}{}{}}
\LANGcs{
\Question{
Tři obchodníci prodávali stejné zboží za stejnou cenu. Postupně ale všichni své ceny upravili následujícím způsobem:

(a) Obchodník A nejprve zboží o \(5\,\%\) zdražil, ale později ho zlevnil o \(10\,\%\). 

(b) Obchodník B nejprve zboží zlevnil o \(2\,\%\) a později ho znovu zlevnil, tentokrát o  \(3\,\%\). 

(c) Obchodník C cenu zboží upravil jenom jednou. Zlevnil ho o \(10\,\%\).


Který z obchodníků bude po úpravě cen prodávat toto zboží za nejvyšší cenu?}
{Obchodník B}{Obchodník A}{Obchodník C}{Zboží bude stát u všech obchodníků pořád stejně.}{}{}}
\LANGsk{
\Question{
Traja obchodníci predávali rovnaký tovar za rovnakú cenu. Postupne všetci svoje ceny upravili nasledujúcim spôsobom:

(a) Obchodník A najskôr tovar o \(5\%\) zdražil, ale neskôr ho zlacnil o \(10\%\). 

(b) Obchodník B najskôr tovar zlacnil o \(2\%\) a neskôr ho znovu zlacnil o \(3\%\). 

(c) Obchodník C cenu tovaru upravil len raz. Zlacnil ho o \(10\%\).


Ktorý obchodník bude po úprave cien predávať tento tovar za najvyššiu cenu?}
{Obchodník B}{Obchodník A}{Obchodník C}{Tovar bude stáť u všetkých troch obchodníkov stále rovnako.}{}{}}
\LANGes{
\Question{
Tres comerciantes vendían el mismo producto por el mismo precio. Sucesivamente ajustaron los precios de la siguiente manera:

(a) El comerciante A primero aumentó el precio del producto en un \(5\%\), pero luego lo redujo en un \(10\%\).

(b) El comerciante B primero redujo el precio del producto en un \(2\%\) y luego lo volvió a reducir, esta vez en un \(3\%\).

(c) El comerciante C ajustó el precio del producto solo una vez. Abarató el producto en un \(10\%\).


¿Qué comerciante ofrece el producto al precio más alto después de los ajustes de precio?}
{Comerciante B}{Comerciante A}{Comerciante C}{Todos los comerciantes venderán el producto al mismo precio.}{}{}}
\End

\Data\ID{2010008304}
\Updated{2022-07-21 05:07:19}
\Level{A}
\SubArea{Percent problems}
\LANGen{
\Question{
At the beginning of the therapy, three monitored patients were of the same weight. Their weights were changing as follows:

(a) Patient A gained \(6\%\) of his weight and then his weight did not change.

(b) Patient B lost \(2\%\) of his weight first, but later he gained \(8\%\) of his weight.

(c) Patient C gained \(2\%\) of his weight and later he gained another \(4\%\) of his weight.


Which patient has the lowest weight after the therapy?}
{Patient B}{Patient A}{Patient C}{After described weight changes everyone will weigh the same.}{}{}}
\LANGpl{
\Question{
Na początku terapii trzech monitorowanych pacjentów miało tę samą wagę. Ich wagi zmieniały się w następujący sposób:

(a) Pacjent A przybrał \(6\%\) swojej wagi a następnie jego waga się nie zmieniła.

(b) Pacjent B najpierw stracił \(2\%\) swojej wagi, ale później przybrał \(8\%\) swojej wagi.

(c) Pacjent C przybrał \(2\%\) swojej wagi a później przybrał kolejne \(4\%\).


Który pacjent ma najmniejszą wagę po terapii?}
{Pacjent B}{Pacjent A}{Pacjent C}{Po opisanych zmianach wagi każdy będzie ważył tak samo.}{}{}}
\LANGcs{
\Question{
Na začátku terapie byly hmotnosti tří sledovaných pacientů stejné. Jejich hmotnosti se měnily následujícím způsobem:

(a) Pacient A přibral \(6\,\%\) své hmotnosti a potom se jeho hmotnost už neměnila.

(b) Pacient B nejprve zhubnul o \(2\,\%\), ale později znovu přibral o \(8\,\%\) své hmotnosti.

(c) Pacient C přibral o \(2\,\%\) své hmotnosti a později opět přibral, tentokrát o \(4\,\%\) své hmotnosti.

Který pacient bude mít po ukončení terapie nejnižší hmotnost?}
{Pacient B}{Pacient A}{Pacient C}{Po daných změnách budou opět všichni vážit stejně.}{}{}}
\LANGsk{
\Question{
Na začiatku terapie boli hmotnosti troch sledovaných pacientov rovnaké. Ich hmotnosti sa menili nasledujúcim spôsobom:

(a) Pacient A pribral \(6\%\) svojej hmotnosti a potom sa jeho hmotnosť už nemenila.

(b) Pacient B najskôr schudol o \(2\%\), ale neskôr znovu pribral o \(8\%\) svojej hmotnosti.

(c) Pacient C pribral o \(2\%\) svojej hmotnosti a neskôr opäť pribral a to o \(4\%\) svojej hmotnosti.

Ktorý pacient bude mať po ukončení terapie najnižšiu hmotnosť?}
{Pacient B}{Pacient A}{Pacient C}{Po daných zmenách budú opäť všetci vážiť rovnako.}{}{}}
\LANGes{
\Question{
Al comienzo de una terapia tres pacientes monitoreados tenían el mismo peso. Sus pesos cambiaban de la siguiente manera:

(a) Paciente A aumentó su peso en un \(6\%\) y luego su peso no cambió.

(b) El paciente B primero perdió un \(2\%\) de su peso, pero luego su peso aumentó en un \(8\%\).

(c) Paciente C aumentó su peso en un \(2\%\)  y luego su peso aumentó en un \(4\%\).

¿Qué paciente tiene el peso más bajo después de la terapia?}
{Paciente B}{Paciente A}{Paciente C}{Después de estos cambios de peso, todos pesarán lo mismo.}{}{}}
\End

\Data\ID{2010018205}
\Updated{2022-11-02 09:11:06}
\Level{A}
\SubArea{Percent problems}
\LANGen{
\Question{
The number of traffic accidents in the monitored cities changed year on year. It increased by \(20\%\) in city A and decreased by \(1\%\) in city B. Decide which of the following statements is correct:
\[
\begin{array}{l}
\text{X: The same number of accidents occurred in city A as in city B the last year.}  \\
\text{Y: There were more accidents in city A than in city B the last year.} \\
\text{Z: There have been fewer accidents in city B than in city A this year.}\\
\end{array}
\]}
{The correctness of any of the above statements cannot be verified from the given data.}{Only statement X is true.}{Only statement Y is true.}{Only statement Z is true.}{}{}}
\LANGpl{
\Question{
Liczba wypadków drogowych w monitorowanych miastach zmieniała się z roku na rok. Wzrosła o \(20\%\) w mieście A i spadła o \(1\%\) w mieście B. Zdecyduj, które z poniższych stwierdzeń jest prawidłowe:
\[
\begin{array}{l}
\text{X: Ta sama liczba wypadków miała miejsce w mieście A jak i w mieście B w ubiegłym roku.}  \\
\text{Y: Więcej wypadków miało miejsce w mieście A niż w mieście B w ubiegłym roku.} \\
\text{Z: Mniej wypadków wydarzyło się w mieście B niż w mieście A w tym roku.}\\
\end{array}
\]}
{Na podstawie podanych danych nie można zweryfikować poprawności któregokolwiek z powyższych stwierdzeń.}{Tylko stwierdzenie X jest prawdziwe.}{Tylko stwierdzenie Y jest prawdziwe.}{Tylko stwierdzenie Z jest prawdziwe.}{}{}}
\LANGcs{
\Question{
Ve sledovaných městech se meziročně změnil počet dopravních nehod. Ve městě A se zvýšil o \(20\%\) a ve městě B se snížil o \(1\%\). Rozhodněte, které z následujících tvrzení je správné:
\[
\begin{array}{l}
\text{X: Ve městě A se loni stal stejný počet dopravních nehod jako ve městě B.}  \\
\text{Y: Ve městě A se loni stalo více dopravních nehod než ve městě B.} \\
\text{Z: Ve městě B se letos stalo méně dopravních nehod než ve městě A.}\\
\end{array}
\]}
{Ze zadaných údajů nelze vyhodnotit platnost žádného z uvedených tvrzení.}{Pravdivé je pouze tvrzení X.}{Pravdivé je pouze tvrzení Y.}{Pravdivé je pouze tvrzení Z.}{}{}}
\LANGsk{
\Question{
V sledovaných mestách sa medziročne zmenil počet dopravných nehôd. V meste A sa zvýšil o \(20\%\) a v meste B sa znížil o \(1\%\). Rozhodnite, ktoré z nasledujícúch tvrdení je pravdivé:
\[
\begin{array}{l}
\text{X: V meste A sa minulý rok stal rovnaký počet dopravných nehôd ako v meste B.}  \\
\text{Y: V meste A sa minulý rok stalo viac dopravních nehôd ako v meste B.} \\
\text{Z: V meste B sa stalo tento rok menej dopravných nehôd ako v meste A.}\\
\end{array}
\]}
{Zo zadaných údajov sa nedá vyhodnotiť pravdivosť žiadneho z uvedených tvrdení.}{Pravdivé je len tvrdenie X.}{Pravdivé je len tvrdenie Y.}{Pravdivé je len tvrdenie Z.}{}{}}
\LANGes{
\Question{
El número de accidentes de tráfico en dos ciudades determinadas cambió de año en año. Aumentó en \(20\%\) en la ciudad A y disminuyó en \(1\%\) en la ciudad B. Decide cuál de las siguientes afirmaciones es correcta:
\[
\begin{array}{l}
\text{X: Se produjo el mismo número de accidentes en la ciudad A que en la ciudad B el año pasado.}  \\
\text{Y: Hubo más accidentes en la ciudad A que en la ciudad B el año pasado.} \\
\text{Z: Ha habido menos accidentes en la ciudad B que en la ciudad A durante este año.}\\
\end{array}
\]}
{La veracidad de cualquiera de las afirmaciones anteriores no puede verificarse a partir de los datos proporcionados.}{Solo la afirmación X es correcta.}{Solo la afirmación Y es correcta.}{Solo la afirmación Z es correcta.}{}{}}
\End

\Data\ID{1003066901}
\Updated{2022-05-09 08:05:39}
\Level{B}
\SubArea{Percent problems}
\LANGen{
\Question{
Two sides of a rectangle were shortened by \( 20\% \). What is the percentage change in the area of the rectangle?}
{it decreased by  \( 36\% \)}{it decreased by  \( 40\% \)}{it increased by \( 40\% \)}{it decreased by  \( 20\% \)}{}{}}
\LANGpl{
\Question{
Dwa boki prostokąta skrócono o  \( 20\% \). O ile procent zmieniło się pole powierzchni tego prostokąta?}
{zmniejszyło się o \( 36\% \)}{zmniejszyło się o   \( 40\% \)}{wzrosło o \( 40\% \)}{zmniejszyło się o  \( 20\% \)}{}{}}
\LANGcs{
\Question{
Každá ze stran obdélníku se zkrátila o \( 20\,\% \). O kolik procent se změnil jeho obsah?}
{Zmenšil se o  \( 36\,\% \)}{Zmenšil se o  \( 40\,\% \)}{Zvětšil se o \( 40\,\% \)}{Zmenšil se o  \( 20\,\% \)}{}{}}
\LANGsk{
\Question{
Dve strany obdĺžnika sa skrátili o \( 20\% \). O koľko percent sa zmenil jeho obsah?}
{Zmenšil sa o  \( 36\% \)}{Zväčšil sa o \( 40\% \)}{Zmenšil sa o \( 40\% \)}{Zmenšil sa o  \( 20\% \)}{}{}}
\LANGes{
\Question{
Dos lados de un rectángulo se acortaron en un \( 20\% \). ¿Cuál es el cambio porcentual en el área del rectángulo?}
{disminuyó un  \( 36\% \)}{disminuyó un    \( 40\% \)}{aumentó en un \( 40\% \)}{disminuyó un   \( 20\% \)}{}{}}
\End

\Data\ID{1003066902}
\Updated{2020-07-15 03:07:12}
\Level{B}
\SubArea{Percent problems}
\LANGen{
\Question{
One side of a rectangle was shortened by \( 10\% \) and the other was lengthened by \( 10\% \). What was the percentage change in the area of the rectangle?}
{it decreased by  \( 1\% \)}{it decreased by  \( 20\% \)}{it increased by   \( 1\% \)}{it did not change}{}{}}
\LANGpl{
\Question{
Jeden bok prostokąta skrócono o  \( 10\% \), a drugi wydłużono o  \( 10\% \). O ile procent zmieniło się jego pole powierzchni?}
{zmniejszyło się o  \( 1\% \)}{zmniejszyło się o   \( 20\% \)}{wzrosło o   \( 1\% \)}{nie zmieniło się}{}{}}
\LANGcs{
\Question{
Jedna strana obdélníku se zkrátila o \( 10\,\% \) a druhá se prodloužila o \( 10\,\% \). Jak se změnil jeho obsah?}
{Zmenšil se o  \( 1\,\% \).}{Zmenšil se o  \( 20\,\% \).}{Zvětšil se o  \( 1\,\% \).}{Nezměnil se.}{}{}}
\LANGsk{
\Question{
Jedna strana obdĺžnika sa skrátila o \( 10\% \) a druhá sa predĺžila o \( 10\% \). Ako sa zmenil jeho obsah?}
{Zmenšil sa o  \( 1\% \)}{Zmenšil sa o  \( 20\% \)}{Zvetšil sa o  \( 1\% \)}{Nezmenil sa}{}{}}
\LANGes{
\Question{
Un lado de un rectángulo se acortó en un \( 10\% \) y el otro se alargó en un  \( 10\% \). ¿Cuál fue el cambio porcentual en el área del rectángulo?}
{disminuyó  un   \( 1\% \)}{disminuyó  un   \( 20\% \)}{aumentó en un   \( 1\% \)}{no cambió}{}{}}
\End

\Data\ID{1003066903}
\Updated{2022-04-23 05:04:26}
\Level{B}
\SubArea{Percent problems}
\LANGen{
\Question{
The number of members of a certain golf club was increasing for the last three years by \( 20\% \) a year. How many members were there in the club three years ago, assuming that  the number of them amounted to \( 216 \) a year ago?}
{\( 150 \) members}{\( 200 \) members}{\( 120 \) members}{\( 250 \) members}{}{}}
\LANGpl{
\Question{
Liczba członków pewnego klubu golfowego wzrastała w ciągu ostatnich trzech lat o \( 20\% \) rocznie. Ilu członków było w klubie trzy lata temu, jeśli w ubiegłym roku było ich \( 216 \)?}
{\( 150 \) członków}{\( 200 \) członków}{\( 120 \) członków}{\( 250 \) członków}{}{}}
\LANGcs{
\Question{
Počet členů golfového klubu v posledních třech letech rostl o \( 20\% \) ročně. Kolik členů měl klub před třemi lety, pokud víme, že loni byl počet členů \( 216 \)?}
{\( 150 \) členů}{\( 200 \) členů}{\( 120 \) členů}{\( 250 \) členů}{}{}}
\LANGsk{
\Question{
Počet členov golfového klubu v posledných troch rokoch narástol o \( 20\% \) ročne. Koľko členov mal klub pred tromi rokmi, ak vieme, že minulý rok bol počet členov \( 216 \)?}
{\( 150 \) členov}{\( 200 \) členov}{\( 120 \) členov}{\( 250 \) členov}{}{}}
\LANGes{
\Question{
El número de miembros de un determinado club de golf aumentó en los últimos tres años en un \( 20\% \) al año. ¿Cuántos miembros tuvo el club hace tres años suponiendo que el número de ellos ascendió a \( 216 \) hace un año?}
{\( 150 \) miembros}{\( 200 \) miembros}{\( 120 \) miembros}{\( 250 \) miembros}{}{}}
\End

\Data\ID{1003066904}
\Updated{2020-07-15 03:07:12}
\Level{B}
\SubArea{Percent problems}
\LANGen{
\Question{
Find the number by \( 2{\small{{}^\text{o}\mkern-5mu/\mkern-3mu_\text{oo}}}\) bigger than \( 100 \).}
{\( 100.2 \)}{\( 102 \)}{\( 120 \)}{\( 100.002 \)}{}{}}
\LANGpl{
\Question{
Znajdź liczbę o  \( 2{\small{{}^\text{o}\mkern-5mu/\mkern-3mu_\text{oo}}} \) większą od  \( 100 \).}
{\( 100{,}2 \)}{\( 102 \)}{\( 120 \)}{\( 100{,}002 \)}{}{}}
\LANGcs{
\Question{
Které číslo je o \( 2{\,\small{{}^\text{o}\mkern-5mu/\mkern-3mu_\text{oo}}} \) větší než \( 100 \)?}
{\( 100{,}2 \)}{\( 102 \)}{\( 120 \)}{\( 100{,}002 \)}{}{}}
\LANGsk{
\Question{
Ktoré číslo je o \( 2{\small{{}^\text{o}\mkern-5mu/\mkern-3mu_\text{oo}}} \) väčšie než \( 100 \)?}
{\( 100{,}2 \)}{\( 102 \)}{\( 120 \)}{\( 100{,}002 \)}{}{}}
\LANGes{
\Question{
Calcula el número en  \( 2{\small{{}^\text{o}\mkern-5mu/\mkern-3mu_\text{oo}}}\) más grande que  \( 100 \).}
{\( 100.2 \)}{\( 102 \)}{\( 120 \)}{\( 100.002 \)}{}{}}
\End

\Data\ID{1003066905}
\Updated{2022-04-23 05:04:26}
\Level{B}
\SubArea{Percent problems}
\LANGen{
\Question{
Find the number by \( 1.6{\small{{}^\text{o}\mkern-5mu/\mkern-3mu_\text{oo}}} \) bigger  than \( 500 \).}
{\( 500.8 \)}{\( 508 \)}{\( 580 \)}{\( 500.016 \)}{}{}}
\LANGpl{
\Question{
Znajdź liczbę o  \( 1{,}6{\small{{}^\text{o}\mkern-5mu/\mkern-3mu_\text{oo}}} \) większą od \( 500 \).}
{\( 500{,}8 \)}{\( 508 \)}{\( 580 \)}{\( 500{,}016 \)}{}{}}
\LANGcs{
\Question{
Které číslo je o \( 1{,}6{\small{{}^\text{o}\mkern-5mu/\mkern-3mu_\text{oo}}} \) větší než \( 500 \)?}
{\( 500{,}8 \)}{\( 508 \)}{\( 580 \)}{\( 500{,}016 \)}{}{}}
\LANGsk{
\Question{
Ktoré číslo je o \( 1{,}6{\small{{}^\text{o}\mkern-5mu/\mkern-3mu_\text{oo}}} \) väčšie než \( 500 \)?}
{\( 500{,}8 \)}{\( 508 \)}{\( 580 \)}{\( 500{,}016 \)}{}{}}
\LANGes{
\Question{
Calcula el número en  \( 1.6{\small{{}^\text{o}\mkern-5mu/\mkern-3mu_\text{oo}}} \) más grande que  \( 500 \).}
{\( 500.8 \)}{\( 508 \)}{\( 580 \)}{\( 500.016 \)}{}{}}
\End

\Data\ID{1003066906}
\Updated{2020-07-15 03:07:12}
\Level{B}
\SubArea{Percent problems}
\LANGen{
\Question{
The results of a survey considering the support for political parties \( Y \) and \( Z \) are presented in the table below. a) How many percentage points have the support for party \( Y \) increased? b) How much has the percentage of people supporting party \( Z \) decreased?
\[\begin{array}{|c|c|c|} \hline & \text{February} & \text{March} \\\hline Y & 16\% & 20\% \\\hline Z & 10\% & 8\% \\\hline \end{array}\]}
{a) by \( 4 \) percentage points, b) by \( 20\% \)}{a) by \( 25 \) percentage points, b) by \( 2\% \)}{a) by \( 25 \) percentage points, b) by \( 20\% \)}{a) by \( 4 \) percentage points, b) by \( 2\% \)}{}{}}
\LANGpl{
\Question{
Wyniki ankiety dotyczącej poparcia dla partii politycznych \( Y \) i \( Z \) przedstawiono w tabeli poniżej. a) O ile punktów procentowych wzrosło poparcie dla partii \( Y \)? b) O ile procent zmniejszyła się liczba osób popierających partię  \( Z \) ?
\[\begin{array}{|c|c|c|} \hline & \text{luty} & \text{marzec} \\\hline Y & 16\% & 20\% \\\hline Z & 10\% & 8\% \\\hline \end{array}\]}
{a) o \( 4 \) punkty procentowe, b) o \( 20\% \)}{a) o \( 25 \) punktów procentowych, b) o \( 2\% \)}{a) o \( 25 \) punktów procentowych, b) o \( 20\% \)}{a) o \( 4 \) punkty procentowe, b) o \( 2\% \)}{}{}}
\LANGcs{
\Question{
V tabulce vidíte výsledky předvolebního průzkumu stran \( Y \) a \( Z \). a) O kolik procentních bodů vzrostla obliba strany \( Y \)? b) O kolik procent klesla obliba strany \( Z \)?
\[\begin{array}{|c|c|c|} \hline & \text{Únor} & \text{Březen} \\\hline Y & 16\,\% & 20\,\% \\\hline Z & 10\,\% & 8\,\% \\\hline \end{array}\]}
{a) o \( 4 \) procentní body, b) o \( 20\,\% \)}{a) o \( 25 \) procentních bodů, b) o \( 2\,\% \)}{a) o \( 25 \) procentních bodů, b) o \( 20\,\% \)}{a) o \( 4 \) procentní body, b) o \( 2\,\% \)}{}{}}
\LANGsk{
\Question{
V tabuľke vidíte výsledky predvolebného prieskumu strán \( Y \) a \( Z \). a) O koľko percentných bodov vzrástla obľúbenosť strany \( Y \)? b) O koľko percent klesla obľúbenosť strany \( Z \)?
\[\begin{array}{|c|c|c|} \hline & \text{Február} & \text{Marec} \\\hline Y & 16\% & 20\% \\\hline Z & 10\% & 8\% \\\hline \end{array}\]}
{a) o \( 4 \) percentné body, b) o \( 20\% \)}{a) o \( 25 \) percentných bodov, b) o \( 2\% \)}{a) o \( 25 \) percentných bodov, b) o \( 20\% \)}{a) o \( 4 \) percentné body, b) o \( 2\% \)}{}{}}
\LANGes{
\Question{
Los resultados de una encuesta que considera el apoyo a los partidos políticos \( Y \) y \( Z \) están presententados en la siguiente tabla. a) ¿Cuántos puntos porcentuales ha aumentado el apoyo a la parte  \( Y \) ? b) ¿Cuántos puntos porcentuales ha disminuido el apoyo a la parte  \( Z \)?
\[\begin{array}{|c|c|c|} \hline & \text{February} & \text{March} \\\hline Y & 16\% & 20\% \\\hline Z & 10\% & 8\% \\\hline \end{array}\]}
{a) en\( 4 \) puntos porcentuales, b) en un \( 20\% \)}{a) en \( 25 \) puntos porcentuales, b) en un \( 2\% \)}{a) en \( 25 \) puntos porcentuales, b) en un \( 20\% \)}{a) en \( 4 \) puntos porcentuales, b) en un \( 2\% \)}{}{}}
\End

\Data\ID{1003066907}
\Updated{2020-08-10 08:08:33}
\Level{B}
\SubArea{Percent problems}
\LANGen{
\Question{
The results of a survey considering the support for political parties \( Y \) and \( Z \) are presented in the table below. a) How many percentage points have the support for party \( Z \) decreased?  b) How much has the percentage of people supporting party \( Y \) increased?
\[\begin{array}{|c|c|c|} \hline & \text{February} & \text{March} \\\hline Y & 16\% & 20\% \\\hline Z & 10\% & 8\% \\\hline \end{array}\]}
{a) by \( 2 \) percentage points, b) by \( 25\% \)}{a) by \( 20 \) percentage points, b) by \( 4\% \)}{a) by \( 20 \) percentage points, b) by \( 25\% \)}{a) by \( 2 \) percentage points, b) by \( 4\% \)}{}{}}
\LANGpl{
\Question{
Wyniki ankiety dotyczącej poparcia dla partii politycznych  \( Y \) i  \( Z \) przedstawiono w poniższej tabeli. a) O ile punktów procentowych zmniejszyło się poparcie dla partii  \( Z \)?  b) O ile procent zwiększyła się liczba osób popierających partię  \( Y \)?
\[\begin{array}{|c|c|c|} \hline & \text{luty} & \text{marzec} \\\hline Y & 16\% & 20\% \\\hline Z & 10\% & 8\% \\\hline \end{array}\]}
{a) o  \( 2 \) punkty procentowe, b) o \( 25\% \)}{a) o \( 20 \) punktów procentowych, b) o \( 4\% \)}{a) o \( 20 \) punktów procentowych, b) o \( 25\% \)}{a) o \( 2 \) punkty procentowe, b) o \( 4\% \)}{}{}}
\LANGcs{
\Question{
V tabulce vidíte výsledky předvolebního průzkumu stran \( Y \) a \( Z \). a) O kolik procentních bodů klesla podpora strany \( Z \)?  b) O kolik procent vzrostla podpora strany \( Y \)?
\[\begin{array}{|c|c|c|} \hline & \text{Únor} & \text{Březen} \\\hline Y & 16\,\%  & 20\,\% \\\hline Z & 10\,\%  & 8\,\% \\\hline \end{array}\]}
{a) o \( 2 \) procentní body, b) o \( 25\,\% \)}{a) o \( 20 \) procentních bodů, b) o \( 4\,\% \)}{a) o \( 20 \) procentních bodů, b) o \( 25\,\% \)}{a) o \( 2 \) procentní body, b) o \( 4\,\% \)}{}{}}
\LANGsk{
\Question{
V tabuľke vidíte výsledky predvolebného prieskumu strán \( Y \) and \( Z \). a) O koľko percentných bodov klesla podpora strany \( Z \)?  b) O koľko percent vzrástla podpora strany \( Y \)?
\[\begin{array}{|c|c|c|} \hline & \text{February} & \text{March} \\\hline Y & 16\% & 20\% \\\hline Z & 10\% & 8\% \\\hline \end{array}\]}
{a) o \( 2 \) percentné body, b) o \( 25\% \)}{a) o \( 20 \) percentných bodov, b) o \( 4\% \)}{a) o \( 20 \) percentných bodov, b) o \( 25\% \)}{a) o \( 2 \) percentné body, b) o \( 4\% \)}{}{}}
\LANGes{
\Question{
Los resultados de una encuesta que consulta el apoyo a los partidos políticos  \( Y \) y \( Z \) están representados en la siguiente tabla. a) ¿Cuántos puntos porcentuales ha disminuido el apoyo para el partido  \( Z \) ?  b)  ¿Cuánto ha aumentado el porcentaje de personas que apoyan al partido  \( Y \)?
\[\begin{array}{|c|c|c|} \hline & \text{February} & \text{March} \\\hline Y & 16\% & 20\% \\\hline Z & 10\% & 8\% \\\hline \end{array}\]}
{a) en \( 2 \) puntos porcentuales, b) en un \( 25\% \)}{a) en \( 20 \) puntos porcentuales, b) en un \( 4\% \)}{a) en \( 20 \) puntos porcentuales, b) en un  \( 25\% \)}{a) en \( 2 \) puntos porcentuales, b) en un \( 4\% \)}{}{}}
\End

\Data\ID{1003066908}
\Updated{2020-08-10 08:08:39}
\Level{B}
\SubArea{Percent problems}
\LANGen{
\Question{
The price of a wedding gown was decreased by \( 50\% \). How many percent should the new price be increased so that the gown  costs as much as before the discount?}
{by \( 100\% \)}{by \( 50\% \)}{by \( 150\% \)}{by \( 200\% \)}{}{}}
\LANGpl{
\Question{
Cena sukni ślubnej została zmniejszona o  \( 50\% \). O ile procent należy podnieść aktualną cenę sukienki tak by kosztowała ona tyle samo co przed obniżką?}
{o \( 100\% \)}{o \( 50\% \)}{o \( 150\% \)}{o \( 200\% \)}{}{}}
\LANGcs{
\Question{
Cena svatebních šatů klesla o \( 50\,\% \). O kolik procent je nyní třeba šaty zdražit, aby stály stejně jako před slevou?}
{o \( 100\,\% \)}{o \( 50\,\% \)}{o \( 150\,\% \)}{o \( 200\,\% \)}{}{}}
\LANGsk{
\Question{
Cena svadobných šiat klesla o \( 50\% \). O koľko percent treba šaty znovu zdražieť, aby stály rovnako ako pred zľavou?}
{o \( 100\% \)}{o \( 50\% \)}{o \( 150\% \)}{o \( 200\% \)}{}{}}
\LANGes{
\Question{
El precio de un vestido de novia se redujo en un  \( 50\% \). ¿Qué porcentaje se debe aumentar del nuevo precio para que el vestido cueste tanto como antes del descuento?}
{en un \( 100\% \)}{en un \( 50\% \)}{en un \( 150\% \)}{en un \( 200\% \)}{}{}}
\End

\Data\ID{1003066909}
\Updated{2022-04-23 05:04:26}
\Level{B}
\SubArea{Percent problems}
\LANGen{
\Question{
The price of a wedding gown was decreased by \( 36\% \). How many percent should the new price be increased so that the gown  costs as much as before the discount?}
{by \( 56.25\% \)}{by \( 36\% \)}{by \( 136\% \)}{by \( 64\% \)}{}{}}
\LANGpl{
\Question{
Cena sukni ślubnej została zmniejszona o  \( 36\% \). O ile procent należy podnieść aktualną cenę sukienki tak, aby  była ona taka sama jak przed obniżką?}
{o \( 56{,}25\% \)}{o  \( 36\% \)}{o \( 136\% \)}{o \( 64\% \)}{}{}}
\LANGcs{
\Question{
Cena svatebních šatů klesla o \( 36\% \). O kolik procent je třeba šaty znovu zdražit, aby stály stejně jako před slevou?}
{o \( 56{,}25\% \)}{o \( 36\% \)}{o \( 136\% \)}{o \( 64\% \)}{}{}}
\LANGsk{
\Question{
Cena svadobných šiat klesla o \( 36\% \). O koľko percent treba šaty znovu zdražiť, aby stály rovnako ako pred zľavou?}
{o \( 56{,}25\% \)}{o \( 36\% \)}{o \( 136\% \)}{o \( 64\% \)}{}{}}
\LANGes{
\Question{
El precio de un vestido de novia se redujo en un  \( 36\% \). ¿Qué porcentaje se debe aumentar del nuevo precio para que el vestido cueste tanto como antes del descuento?}
{en un  \( 56.25\% \)}{en un  \( 36\% \)}{en un \( 136\% \)}{en un  \( 64\% \)}{}{}}
\End

\Data\ID{1003082201}
\Updated{2020-08-10 08:08:26}
\Level{B}
\SubArea{Percent problems}
\LANGen{
\Question{
The price of the commodity did not change if}
{it was lowered by \( 20\% \), and the new price was increased by \( 25\% \).}{it was lowered by \( 10\% \), and the new price was increased by \( 10\% \).}{it was increased by \( 30\% \), and the new price was lowered by \( 30\% \).}{it was increased by \( 20\% \), and the new price was lowered by \( 25\% \).}{}{}}
\LANGpl{
\Question{
Cena towaru nie uległa zmianie, jeśli}
{obniżono ją o \( 20\% \), a nową cenę podniesiono o \( 25\% \).}{obniżono ją o  \( 10\% \), a nową cenę podniesiono o  \( 10\% \).}{podniesiono ją o  \( 30\% \), a nową cenę obniżono o  \( 30\% \).}{podniesiono ją o  \( 20\% \), a nową cenę obniżono o  \( 25\% \).}{}{}}
\LANGcs{
\Question{
Cena zboží se nezměnila, pokud}
{byla snížena o \( 20\,\% \) a následně znovu zvýšena o \( 25\,\% \).}{byla snížena o \( 10\,\% \) a následně znovu zvýšena o \( 10\,\% \).}{byla zvýšena o \( 30\,\% \) a následně znovu snížena o \( 30\,\% \).}{byla zvýšena o \( 20\,\% \) a následně znovu snížena o \( 25\,\% \).}{}{}}
\LANGsk{
\Question{
Cena tovaru sa nezmenila, ak}
{bola znížená o \( 20\% \) a následne znovu zvýšená o \( 25\% \).}{bola znížená o \( 10\% \) a následne znovu zvýšená o \( 10\% \).}{bola zvýšená o \( 30\% \) a následne znovu znížená o \( 30\% \).}{bola zvýšená o \( 20\% \) a následne znovu znížená o \( 25\% \).}{}{}}
\LANGes{
\Question{
El precio de un producto no cambia si}
{se redujo en un  \( 20\% \), y el nuevo precio se incrementó en un  \( 25\% \).}{se incrementó en un  \( 10\% \), y el nuevo precio se incrementó en un  \( 10\% \).}{se incrementó en un  \( 30\% \), y el nuevo precio se redujo en un  \( 30\% \).}{se incrementó en un \( 20\% \), y el nuevo precio se redujo en un  \( 25\% \).}{}{}}
\End

\Data\ID{1003082202}
\Updated{2022-04-23 05:04:26}
\Level{B}
\SubArea{Percent problems}
\LANGen{
\Question{
If \( 20\% \) of the number \( 55 \) is by \( 4 \) smaller than \( 30\% \) of the number \( x \), then:}
{\( x=50 \)}{\( x=75 \)}{\( x=45 \)}{\( x=33 \)}{}{}}
\LANGpl{
\Question{
Jeśli  \( 20\% \) liczby  \( 55 \) jest o \( 4 \) mniejsze od \( 30\% \) liczby \( x \), to:}
{\( x=50 \)}{\( x=75 \)}{\( x=45 \)}{\( x=33 \)}{}{}}
\LANGcs{
\Question{
Jestli je \( 20\% \) čísla \( 55 \) o \( 4 \) menší než \( 30\% \) čísla \( x \), pak:}
{\( x=50 \)}{\( x=75 \)}{\( x=45 \)}{\( x=33 \)}{}{}}
\LANGsk{
\Question{
Ak je \( 20\% \) čísla \( 55 \) o \( 4 \) menšie než \( 30\% \) čísla \( x \), potom:}
{\( x=50 \)}{\( x=75 \)}{\( x=45 \)}{\( x=33 \)}{}{}}
\LANGes{
\Question{
Si el  \( 20\% \) del número  \( 55 \) es en \( 4 \) unidades menor que el \( 30\% \) del número \( x \), entonces:}
{\( x=50 \)}{\( x=75 \)}{\( x=45 \)}{\( x=33 \)}{}{}}
\End

\Data\ID{1003082203}
\Updated{2020-07-15 03:07:12}
\Level{B}
\SubArea{Percent problems}
\LANGen{
\Question{
\( 30\% \) of some goods were sold with a \( 15\% \) profit and the remaining amount was sold with a \( 10\% \) profit. What was the overall profit from the sale?}
{\( 11.5\% \)}{\( 12.5\% \)}{\( 25\% \)}{\( 12\% \)}{}{}}
\LANGpl{
\Question{
\( 30\% \) pewnego towaru sprzedano z zyskiem \( 15\% \), a pozostałą część z zyskiem  \( 10\% \). Jaki był ogólny zysk ze sprzedaży?}
{\( 11{,}5\% \)}{\( 12{,}5\% \)}{\( 25\% \)}{\( 12\% \)}{}{}}
\LANGcs{
\Question{
\( 30\% \) zboží bylo prodáno s \( 15\% \) ziskem a zbývající množství se prodalo s \( 10\% \) ziskem. Jaký byl celkový zisk z prodeje?}
{\( 11{,}5\% \)}{\( 12{,}5\% \)}{\( 25\% \)}{\( 12\% \)}{}{}}
\LANGsk{
\Question{
\( 30\% \) tovaru bolo predané so \( 15\% \) ziskom a zostávajúce množstvo sa predalo s \( 10\% \) ziskom. Aký bol celkový zisk z predaja?}
{\( 11{,}5\% \)}{\( 12{,}5\% \)}{\( 25\% \)}{\( 12\% \)}{}{}}
\LANGes{
\Question{
El \( 30\% \) de algunos bienes se vendieron con un beneficio del \( 15\% \) y la cantidad restante se vendió con un beneficio del  \( 10\% \)¿Cuál fue el beneficio total de la venta?}
{\( 11.5\% \)}{\( 12.5\% \)}{\( 25\% \)}{\( 12\% \)}{}{}}
\End

\Data\ID{1003082204}
\Updated{2022-04-23 05:04:26}
\Level{B}
\SubArea{Percent problems}
\LANGen{
\Question{
By how many percent should we lengthen the edge of a cube so that its surface area is increased by \( 21\% \)?}
{by \( 10\% \)}{by \( 21\% \)}{by \( 11\% \)}{by \( 42\% \)}{}{}}
\LANGpl{
\Question{
O ile procent należy wydłużyć krawędź sześcianu, aby jego pole powierzchni wzrosło o  \( 21\% \)?}
{o \( 10\% \)}{o \( 21\% \)}{o \( 11\% \)}{o \( 42\% \)}{}{}}
\LANGcs{
\Question{
O kolik procent je třeba prodloužit délku strany krychle, aby se její povrch zvětšil o \( 21\% \)?}
{o \( 10\% \)}{o \( 21\% \)}{o \( 11\% \)}{o \( 42\% \)}{}{}}
\LANGsk{
\Question{
O koľko percent treba predĺžiť dĺžku strany kocky, aby sa jej povrch zväčšil o \( 21\% \)?}
{o \( 10\% \)}{o \( 21\% \)}{o \( 11\% \)}{o \( 42\% \)}{}{}}
\LANGes{
\Question{
¿En qué porcentaje deberíamos alargar la arista de un cubo para que su área de aumente en un \( 21\% \)?}
{en un \( 10\% \)}{en un \( 21\% \)}{en un  \( 11\% \)}{en un  \( 42\% \)}{}{}}
\End

\Data\ID{1003082205}
\Updated{2020-08-22 04:08:27}
\Level{B}
\SubArea{Percent problems}
\LANGen{
\Question{
One side of a rectangle was increased by \( 10\% \). By how many percent has the area of the rectangle changed?}
{it has increased by \( 10\% \)}{it has decreased by  \( 10\% \)}{it has increased by  \( 20\% \)}{it has decreased by  \( 20\% \)}{}{}}
\LANGpl{
\Question{
Jeden bok prostokąta zwiększono o \( 10\% \). O ile procent zmieniło się pole prostokąta?}
{zwiększyło się o \( 10\% \)}{zmniejszyło się o  \( 10\% \)}{zwiększyło się o  \( 20\% \)}{zmniejszyło się o  \( 20\% \)}{}{}}
\LANGcs{
\Question{
Jedna strana obdélníku se zvětšila o \( 10\,\% \). O kolik procent se změnil jeho obsah?}
{zvětšil se o \( 10\,\% \)}{zmenšil se o  \( 10\,\% \)}{zvětšil se o  \( 20\,\% \)}{zmenšil se o  \( 20\,\% \)}{}{}}
\LANGsk{
\Question{
Strana obdĺžnika sa zväčšila o \( 10\% \). O koľko percent sa zmenil jeho obsah?}
{zväčšil sa o \( 10\% \)}{zmenšil sa o  \( 10\% \)}{zväčšil sa o  \( 20\% \)}{zmenšil sa o  \( 20\% \)}{}{}}
\LANGes{
\Question{
Un lado de un rectángulo se incrementó en un \( 10\% \). ¿En qué porcentaje ha cambiado el área del rectángulo?}
{ha aumentado en un  \( 10\% \)}{ha disminuido en un  \( 10\% \)}{ha aumentado en un  \( 20\% \)}{ha disminuido en un   \( 20\% \)}{}{}}
\End

\Data\ID{1003082206}
\Updated{2020-08-22 04:08:56}
\Level{B}
\SubArea{Percent problems}
\LANGen{
\Question{
The price of the book was increased by \( 4.5 \) złotys which is \( 15\% \) of the initial price.  Then it was increased by another \( 1.5 \) złotys. By how many percent did the overall price of the book  increase?}
{by \( 20\% \)}{by \( 6 \) złotys}{by \( 25\% \)}{by \( 16\% \)}{}{}}
\LANGpl{
\Question{
Cenę książki podwyższono o \( 4{,}5 \) zł, co stanowi \( 15\% \) ceny początkowej. Następnie jeszcze podwyższono cenę o \( 1{,}5 \) zł. O ile procent łącznie podwyższono cenę tej książki?}
{o \( 20\% \)}{o \( 6 \) zł}{o \( 25\% \)}{o \( 16\% \)}{}{}}
\LANGcs{
\Question{
Kniha zdražila o \( 4{,}5 \) zlotých, což je o \( 15\% \) z původní ceny. Následně zdražila o dalších \( 1{,}5 \) zlotých. Jak celkem narostla cena knihy?}
{o \( 20\% \)}{o \( 6 \) zlotých}{o \( 25\% \)}{o \( 16\% \)}{}{}}
\LANGsk{
\Question{
Kniha zdražela o \( 4{,}5 \) EUR, čo je o \( 15\% \) z pôvodnej ceny. Následne zdražela o ďalších \( 1{,}5 \) EUR. Ako celkom narástla cena knihy?}
{o \( 20\% \)}{o \( 6 \) EUR}{o \( 25\% \)}{o \( 16\% \)}{}{}}
\LANGes{
\Question{
El precio de un libro se incrementó en  \( 4.5 \) złotys que es el \( 15\% \) del precio inicial. Luego se incrementó en otros  \( 1.5 \) złotys. ¿En qué porcentaje aumentó el precio total del libro?}
{en un  \( 20\% \)}{en \( 6 \) złotys}{en un \( 25\% \)}{en un \( 16\% \)}{}{}}
\End

\Data\ID{1003082207}
\Updated{2022-04-23 05:04:26}
\Level{B}
\SubArea{Percent problems}
\LANGen{
\Question{
The angle \( \beta \)  is by \( 25\% \) less  than the angle \( \gamma \) and by \( 50\% \) greater than the angle \( \alpha \).  The angle \( \gamma \) is}
{by \( 100\% \) greater than \( \alpha \).}{by \( 50\% \) greater than \( \alpha \).}{by \( 75\% \) greater than \( \alpha \).}{by \( 25\% \) greater than \( \alpha \).}{}{}}
\LANGpl{
\Question{
Kąt  \( \beta \) ma miarę o \( 25\% \) mniejszą niż kąt  \( \gamma \)  i o \( 50\% \) większą niż kąt \( \alpha \).  Miara kąta  \( \gamma \) jest}
{o \( 100\% \) większa niż  \( \alpha \).}{o \( 50\% \) większa niż  \( \alpha \).}{o \( 75\% \) większa niż  \( \alpha \).}{o \( 25\% \) większa niż  \( \alpha \).}{}{}}
\LANGcs{
\Question{
Úhel \( \beta \)  je o \( 25\% \) menší než úhel \( \gamma \) a o  \( 50\% \) větší než úhel \( \alpha \).  Úhel \( \gamma \) je}
{o \( 100\% \) větší než \( \alpha \).}{o \( 50\% \) větší než \( \alpha \).}{o \( 75\% \) větší než \( \alpha \).}{o \( 25\% \) větší než \( \alpha \).}{}{}}
\LANGsk{
\Question{
Uhol \( \beta \)  je o \( 25\% \) menší než uhol \( \gamma \) a o  \( 50\% \) väčší než uhol \( \alpha \).  Uhol \( \gamma \) je}
{o \( 100\% \) väčší než \( \alpha \).}{o \( 50\% \) väčší než \( \alpha \).}{o \( 75\% \) väčší než \( \alpha \).}{o \( 25\% \) väčší než \( \alpha \).}{}{}}
\LANGes{
\Question{
El ángulo  \( \beta \) es un \( 25\% \) menor que el ángulo \( \gamma \) y un \( 50\% \) mayor que el ángulo  \( \alpha \).  El ángulo  \( \gamma \) es}
{en un \( 100\% \) mayor que  \( \alpha \).}{en un \( 50\% \) mayor que  \( \alpha \).}{en un \( 75\% \) mayor que \( \alpha \).}{en un \( 25\% \) mayor que  \( \alpha \).}{}{}}
\End

\Data\ID{1003082208}
\Updated{2020-07-15 03:07:12}
\Level{B}
\SubArea{Percent problems}
\LANGen{
\Question{
Find a positive number whose square is by \( 500\% \) greater than this number.}
{\( 6 \)}{\( 5 \)}{\( 7 \)}{\( 8 \)}{}{}}
\LANGpl{
\Question{
Liczbą dodatnią, której kwadrat jest większy od niej o  \( 500\% \) jest:}
{\( 6 \)}{\( 5 \)}{\( 7 \)}{\( 8 \)}{}{}}
\LANGcs{
\Question{
Najdi kladné číslo, jehož druhá mocnina je o \( 500\,\% \) větší než toto číslo.}
{\( 6 \)}{\( 5 \)}{\( 7 \)}{\( 8 \)}{}{}}
\LANGsk{
\Question{
Nájdi kladné číslo, ktorého druhá mocnina je o \( 500\% \) väčšia než toto číslo.}
{\( 6 \)}{\( 5 \)}{\( 7 \)}{\( 8 \)}{}{}}
\LANGes{
\Question{
Encuentra un número positivo cuyo cuadrado sea un  \( 500\% \) mayor que este número.}
{\( 6 \)}{\( 5 \)}{\( 7 \)}{\( 8 \)}{}{}}
\End

\Data\ID{1003082209}
\Updated{2020-07-15 03:07:12}
\Level{B}
\SubArea{Percent problems}
\LANGen{
\Question{
Two circles are given. The radius of the first circle is by \( 30\% \) longer than the radius of the other circle. It means that the area of the first circle is bigger than the area of the second one}
{by more than \( 60\% \).}{by less than \( 50\% \), but more than \( 40\% \).}{by less than \( 60\% \), but more than \( 50\% \).}{by exactly \( 60\% \).}{}{}}
\LANGpl{
\Question{
Dane są dwa koła. Promień pierwszego koła jest większy od promienia drugiego koła o \( 30\% \). Wynika stąd, że pole pierwszego koła jest większe od pola drugiego koła}
{o więcej niż  \( 60\% \).}{o mniej niż \( 50\% \), ale więcej niż  \( 40\% \).}{o  mniej niż  \( 60\% \), ale więcej niż \( 50\% \).}{dokładnie o \( 60\% \).}{}{}}
\LANGcs{
\Question{
Jsou dány dva kruhy. Poloměr prvního kruhu je o \( 30\% \) větší než poloměr druhého kruhu. To znamená, že obsah prvního kruhu je větší než obsah druhého o}
{více než \( 60\% \).}{méně než \( 50\% \), ale více než \( 40\% \).}{méně než \( 60\% \), ale více než \( 50\% \).}{přesně \( 60\% \).}{}{}}
\LANGsk{
\Question{
Sú dané dva kruhy. Polomer prvého kruhu je o \( 30\% \) väčší než polomer druhého kruhu. To znamená, že obsah prvého kruhu je väčší než obsah druhého kruhu o}
{viac než \( 60\% \).}{menej než \( 50\% \), ale viac než \( 40\% \).}{menej než \( 60\% \), ale viac než \( 50\% \).}{presne \( 60\% \).}{}{}}
\LANGes{
\Question{
Tenemos dos círculos. El radio del primer círculo es un  \( 30\% \) más largo que el radio del otro círculo. Esto significa que el área del primer círculo es más grande que el área del segundo}
{en más del  \( 60\% \).}{en menos del  \( 50\% \), pero más del  \( 40\% \).}{en menos del\( 60\% \), pero más del  \( 50\% \).}{en exactamente el  \( 60\% \).}{}{}}
\End

\Data\ID{1103066910}
\Updated{2020-08-22 04:08:50}
\Level{B}
\SubArea{Percent problems}
\IMG{1103066901.pdf}
\LANGen{
\Question{
The diagram shows the number of individual grades that the brothers Bolek, Olek and Alek received at the end of the first semester. How many percent of all grades received by the brothers are better than E?}
{\( 75\% \)}{\( 50\% \)}{\( 15\% \)}{\( 25\% \)}{}{}}
\LANGpl{
\Question{
Poniższy diagram przedstawia liczbę poszczególnych ocen, które bracia Bolek, Olek i Alek otrzymali na koniec pierwszego półrocza.  Jaki procent wszystkich ocen uzyskanych przez braci stanowią oceny wyższe od  E?}
{\( 75\% \)}{\( 50\% \)}{\( 15\% \)}{\( 25\% \)}{}{}}
\LANGcs{
\Question{
Graf znázorňuje počet známek (hodnoceno na škále A - E), které bratři Bolek, Olek a Alek dostali na konci prvního pololetí. Kolik procent známek všech tří bratrů je lepších než E?}
{\( 75\,\% \)}{\( 50\,\% \)}{\( 15\,\% \)}{\( 25\,\% \)}{}{}}
\LANGsk{
\Question{
Graf znázorňuje počet známok (hodnotenie na škále A - E), ktoré bratia Bolek, Olek a Alek dostali na konci prvého polroka. Koľko percent známok všetkých troch bratov je lepších ako stupeň E?}
{\( 75\% \)}{\( 50\% \)}{\( 15\% \)}{\( 25\% \)}{}{}}
\LANGes{
\Question{
El diagrama muestra la cantidad de calificaciones individuales que recibieron los hermanos Bolek, Olek y Alek al final del primer semestre. ¿Qué pocentaje de todas las calificaciones recibidas por los hermanos son mejores que E?}
{\( 75\% \)}{\( 50\% \)}{\( 15\% \)}{\( 25\% \)}{}{}}
\End

\Data\ID{2010006201}
\Updated{2022-08-25 10:08:59}
\Level{B}
\SubArea{Percent problems}
\LANGen{
\Question{
If \( 40\% \) of the number \( 60 \) is by \( 18 \) smaller than \( 60\% \) of the number \( x \), then:}
{\( x=70 \)}{\( x=20 \)}{\( x=126 \)}{\( x=193 \)}{}{}}
\LANGpl{
\Question{
Jeżeli \( 40\% \) liczby \( 60 \) jest o \( 18 \) mniejsze niż \( 60\% \) liczby  \( x \), to:}
{\( x=70 \)}{\( x=20 \)}{\( x=126 \)}{\( x=193 \)}{}{}}
\LANGcs{
\Question{
Jestli je \( 40\,\% \) čísla \( 60 \) o \( 18 \) menší než \( 60\,\% \) čísla \( x \), pak:}
{\( x=70 \)}{\( x=20 \)}{\( x=126 \)}{\( x=193 \)}{}{}}
\LANGsk{
\Question{
Ak \( 40\% \) z čísla \( 60 \) je o \( 18 \) menšie než \( 60\% \) z čísla \( x \), tak:}
{\( x=70 \)}{\( x=20 \)}{\( x=126 \)}{\( x=193 \)}{}{}}
\LANGes{
\Question{
Si el \( 40\% \) del número \( 60 \) es en \( 18 \) unidades menor que el \( 60\% \) del número \( x \) entonces:}
{\( x=70 \)}{\( x=20 \)}{\( x=126 \)}{\( x=193 \)}{}{}}
\End

\Data\ID{2010006202}
\Updated{2022-08-25 10:08:38}
\Level{B}
\SubArea{Percent problems}
\LANGen{
\Question{
By how many percent should we shorten the edge of a cube so that its surface area is decreased by \( 19\% \)?}
{by \( 10\% \)}{by \( 19\% \)}{by \( 9\% \)}{by \( 90\% \)}{}{}}
\LANGpl{
\Question{
O ile procent powinniśmy skrócić krawędź sześcianu, aby zmniejszyć jego powierzchnię o \( 19\% \)?}
{o \( 10\% \)}{o \( 19\% \)}{o \( 9\% \)}{o \( 90\% \)}{}{}}
\LANGcs{
\Question{
O kolik procent je třeba zkrátit délku strany krychle, aby se její povrch zmenšil o \( 19\,\% \)?}
{o \( 10\,\% \)}{o \( 19\,\% \)}{o \( 9\,\% \)}{o \( 90\,\% \)}{}{}}
\LANGsk{
\Question{
O koľko percent treba skrátiť dĺžku strany kocky, aby sa jej povrch zmenšil o \( 19\% \)?}
{o \( 10\% \)}{o \( 19\% \)}{o \( 9\% \)}{o \( 90\% \)}{}{}}
\LANGes{
\Question{
¿En qué porcentaje deberíamos alargar la arista de un cubo para que su área de aumente en un \( 19\% \)?}
{en \( 10\% \)}{en \( 19\% \)}{en \( 9\% \)}{en \( 90\% \)}{}{}}
\End

\Data\ID{2010006203}
\Updated{2022-08-25 10:08:38}
\Level{B}
\SubArea{Percent problems}
\LANGen{
\Question{
Two sides of a rectangle were increased by \( 10\% \). What is the percentage change of the area of the rectangle?}
{it increased by  \( 21\% \)}{it increased by  \( 20\% \)}{it decreased by  \( 40\% \)}{it decreased by  \( 20\% \)}{}{}}
\LANGpl{
\Question{
Dwa boki prostokąta zostały powiększone o \( 10\% \). Jak zmieniło się pole prostokąta w procentach?}
{zwiększyło się o \( 21\% \)}{zwiększyło się o  \( 20\% \)}{zmniejszyło się o  \( 40\% \)}{zmniejszyło się o  \( 20\% \)}{}{}}
\LANGcs{
\Question{
Každá ze stran obdélníku se zvětšila o  \( 10\,\% \). O kolik procent se změnil jeho obsah?}
{Zvětšil se o \( 21\,\% \).}{Zvětšil se o \( 20\,\% \).}{Zmenšil se o \( 40\,\% \).}{Zmenšil se o \( 20\,\% \).}{}{}}
\LANGsk{
\Question{
Dve strany obdĺžnika sa zväčšili o \( 10\% \). O koľko percent sa zmenil jeho obsah?}
{Zväčšil sa o \( 21\% \).}{Zväčšil sa o \( 20\% \).}{Zmenšil sa o  \( 40\% \).}{Zmenšil sa o \( 20\% \).}{}{}}
\LANGes{
\Question{
Dos lados de un rectángulo se incrementaron en \( 10\% \). ¿Cuál es el cambio porcentual en el área del rectángulo?}
{aumentó en un   \( 21\% \)}{aumentó en un  \( 20\% \)}{disminuyó un \( 40\% \)}{disminuyó un \( 20\% \)}{}{}}
\End

\Data\ID{2010006204}
\Updated{2022-08-25 10:08:38}
\Level{B}
\SubArea{Percent problems}
\LANGen{
\Question{
The number of permanent residents of a city has decreased by \( 10\% \) per year in the last three years. How many residents did the city have three years ago, if we know that the last year the population was \(972\)?}
{\( 1\,200 \)}{\( 1\,080 \)}{\( 900 \)}{\( 1\,070 \)}{}{}}
\LANGpl{
\Question{
Liczba stałych mieszkańców miasta spadała w ciągu ostatnich trzech lat o \( 10\% \) rocznie. Ilu mieszkańców miało miasto trzy lata temu, jeśli wiemy, że w zeszłym roku było \(972\) mieszkańców?}
{\( 1\,200 \)}{\( 1\,080 \)}{\( 900 \)}{\( 1\,070 \)}{}{}}
\LANGcs{
\Question{
Počet stálých obyvatel obce v posledních třech letech klesal o \( 10\,\% \) ročně. Kolik obyvatel měla obec před třemi lety, pokud víme, že loni byl počet obyvatel \(972\)?}
{\( 1\,200 \)}{\( 1\,080 \)}{\( 900 \)}{\( 1\,070 \)}{}{}}
\LANGsk{
\Question{
Počet stálých obyvateľov obce za posledné tri roky klesal  o \( 10\% \) ročne. Koľko obyvateľov mala obec pred troma rokmi, ak vieme, že minulý rok bol počet obyvateľov \(972\)?}
{\( 1\,200 \)}{\( 1\,080 \)}{\( 900 \)}{\( 1\,070 \)}{}{}}
\LANGes{
\Question{
El número de residentes permanentes de una ciudad va disminuyendo un \( 10\% \) cada año en los últimos tres años. ¿Cuántos residentes vivieron en la ciudad hace tres años si sabemos que el último año la población fue de \(972\)?}
{\( 1\,200 \)}{\( 1\,080 \)}{\( 900 \)}{\( 1\,070 \)}{}{}}
\End

\Data\ID{2010006205}
\Updated{2022-08-25 10:08:38}
\Level{B}
\SubArea{Percent problems}
\LANGen{
\Question{
The price of oil rose by \(60\%\). By how many percent it is necessary to reduce the price of oil so that it costs the same as before?}
{by \( 37.5\% \)}{by \( 60\% \)}{by \( 40\% \)}{by \( 160\% \)}{}{}}
\LANGpl{
\Question{
Cena ropy wzrosła o \(60\%\). O ile procent trzeba obniżyć cenę ropy, aby kosztowała tyle samo, co poprzednio.}
{o \( 37{,}5\% \)}{o \( 60\% \)}{o \( 40\% \)}{o \( 160\% \)}{}{}}
\LANGcs{
\Question{
Cena ropy vzrostla o  \(60\,\%\). O kolik procent je třeba ropu zlevnit, aby stála stejně jako původně?}
{o \( 37{,}5\,\% \)}{o \( 60\,\% \)}{o \( 40\,\% \)}{o \( 160\,\% \)}{}{}}
\LANGsk{
\Question{
Cena ropy vzrástla o \(60\%\). O koľko percent treba ropu znovu zlacniť, aby stála rovnako ako pred zľavou?}
{o \( 37{,}5\% \)}{o \( 60\% \)}{o \( 40\% \)}{o \( 160\% \)}{}{}}
\LANGes{
\Question{
El precio del petróleo subió un \(60\%\). ¿En cuánto por ciento es necesario reducir el precio del petróleo para que cueste lo mismo que antes?}
{un \( 37.5\% \)}{un \( 60\% \)}{un \( 40\% \)}{un \( 160\% \)}{}{}}
\End

\Data\ID{2010006206}
\Updated{2022-08-25 10:08:38}
\Level{B}
\SubArea{Percent problems}
\LANGen{
\Question{
Find the number by \( 2.5{\small{{}^\text{o}\mkern-5mu/\mkern-3mu_\text{oo}}} \) greater  than \( 800 \).}
{\(802\)}{\(820\)}{\(800.2\)}{\(800.025\)}{}{}}
\LANGpl{
\Question{
Znajdź liczbę o \( 2{,}5{\small{{}^\text{o}\mkern-5mu/\mkern-3mu_\text{oo}}} \) większą niż \( 800 \).}
{\(802\)}{\(820\)}{\(800{,}2\)}{\(800{,}025\)}{}{}}
\LANGcs{
\Question{
Které číslo je o \( 2{,}5\,{\small{{}^\text{o}\mkern-5mu/\mkern-3mu_\text{oo}}} \) větší než \( 800 \)?}
{\(802\)}{\(820\)}{\(800{,}2\)}{\(800{,}025\)}{}{}}
\LANGsk{
\Question{
Ktoré číslo je o \( 2{,}5{\small{{}^\text{o}\mkern-5mu/\mkern-3mu_\text{oo}}} \) väčšie ako \( 800 \)?}
{\(802\)}{\(820\)}{\(800{,}2\)}{\(800{,}025\)}{}{}}
\LANGes{
\Question{
Calcula el número en \( 2.5{\small{{}^\text{o}\mkern-5mu/\mkern-3mu_\text{oo}}} \) más grande que \( 800 \).}
{\(802\)}{\(820\)}{\(800.2\)}{\(800.025\)}{}{}}
\End

\Data\ID{2010006207}
\Updated{2022-08-25 10:08:38}
\Level{B}
\SubArea{Percent problems}
\LANGen{
\Question{
If we increase an unknown number \(x\) 
by \(8\, \%\), we
get \(432\).
Find \(x\).}
{\(400\)}{\(466.56\)}{\(397.44\)}{\(380\)}{}{}}
\LANGpl{
\Question{
Jeśli zwiększymy nieznaną liczbę \(x\) 
o \(8\, \%\), otrzymamy \(432\).
Znajdź \(x\).}
{\(400\)}{\(466{,}56\)}{\(397{,}44\)}{\(380\)}{}{}}
\LANGcs{
\Question{
Zvětšíme-li neznámé číslo \(x\) 
o \(8\, \%\), dostaneme číslo \(432\).
Určete neznámé číslo \(x\).}
{\(400\)}{\(466{,}56\)}{\(397{,}44\)}{\(380\)}{}{}}
\LANGsk{
\Question{
Ak zväčšíme neznáme číslo o \(8\, \%\), dostaneme číslo \(432\). Určte neznáme číslo \(x\).}
{\(400\)}{\(466{,}56\)}{\(397{,}44\)}{\(380\)}{}{}}
\LANGes{
\Question{
Si aumentamos un número desconocido \(x\) 
en un  \(8\, \%\), obtenemos \(432\).
Determina \(x\).}
{\(400\)}{\(466.56\)}{\(397.44\)}{\(380\)}{}{}}
\End

\Data\ID{2010006208}
\Updated{2022-08-25 10:08:38}
\Level{B}
\SubArea{Percent problems}
\LANGen{
\Question{
There are \(14\) right-handed and \(18\) left-handed children among the children attending the ceramic group. Find the difference between percentages  of right-handers and left-handers in the group. (Round the result to one decimal place.)}
{differ by \(27.3\, \%\)}{differ by \(63.6\, \%\)}{differ by \(36.4\, \%\)}{differ by \(57.1\, \%\)}{}{}}
\LANGpl{
\Question{
Wśród dzieci uczęszczających do grupy ceramicznej jest \(14\) dzieci praworęcznych i  \(18\) leworęcznych.  Znajdź różnicę między odsetkiem osób praworęcznych i leworęcznych w grupie. (Zaokrąglij wynik do jednego miejsca po przecinku.)}
{różnica to \(27{,}3\, \%\)}{różnica to \(63{,}6\, \%\)}{różnica to \(36{,}4\, \%\)}{różnica \(57{,}1\, \%\)}{}{}}
\LANGcs{
\Question{
Mezi dětmi, které navštěvují keramický kroužek je \(14\) praváků a \(18\) leváků. O kolik se liší procentuální zastoupení praváků a leváků v kroužku? (Výsledek zaokrouhlete na jedno desetinné místo.)}
{o \(27{,}3\, \%\)}{o \(63{,}6\, \%\)}{o \(36{,}4\, \%\)}{o \(57{,}1\, \%\)}{}{}}
\LANGsk{
\Question{
Medzi deťmi na keramickom krúžku je \(14\) pravákov a \(18\) ľavákov. O koľko sa líši percentuálne zastúpenie pravákov a ľavákov na krúžku? (Výsledok zaokrúhlite na jedno desatinné miesto.)}
{o \(27{,}3\, \%\)}{o \(63{,}6\, \%\)}{o \(36{,}4\, \%\)}{o \(57{,}1\, \%\)}{}{}}
\LANGes{
\Question{
Hay \(14\) niños diestros y \(18\) zurdos entre los niños que asisten al grupo de cerámica. Calcula la diferencia entre los porcentajes de diestros y zurdos en el grupo. (Redondea el resultado a un decimal.)}
{difieren en \(27.3\, \%\)}{difieren en \(63.6\, \%\)}{difieren en \(36.4\, \%\)}{difieren en \(57.1\, \%\)}{}{}}
\End

\Data\ID{2010006209}
\Updated{2022-11-02 09:11:40}
\Level{B}
\SubArea{Percent problems}
\LANGen{
\Question{
The angle \( \beta \)  is by \( 20\% \) smaller  than the angle \( \gamma \) and by \( 100\% \) greater than the angle \( \alpha \).  The angle \( \alpha \) is}
{by \( 60\% \) smaller than \( \gamma \).}{by \( 40\% \) smaller than \( \gamma \).}{by \( 20\% \) smaller than \( \gamma \).}{by \( 80\% \) smaller than \( \gamma \).}{}{}}
\LANGpl{
\Question{
Kąt \( \beta \)  jest o \( 20\% \) mniejszy niż kąt \( \gamma \) i o \( 100\% \) większy od kąta \( \alpha \). Kąt \( \alpha \) jest}
{o \( 60\% \) mniejszy niż \( \gamma \).}{o \( 40\% \) mniejszy niż \( \gamma \).}{o \( 20\% \) mniejszy niż \( \gamma \).}{o \( 80\% \) mniejszy niż \( \gamma \).}{}{}}
\LANGcs{
\Question{
Úhel \( \beta \)  je o \( 20\,\% \) menší než úhel \( \gamma \) a o \( 100\,\% \) větší než úhel \( \alpha \).  Úhel \( \alpha \) je}
{o \( 60\,\% \) menší než \( \gamma \).}{o \( 40\,\% \) menší než \( \gamma \).}{o \( 20\,\% \) menší než \( \gamma \).}{o \( 80\,\% \) menší než \( \gamma \).}{}{}}
\LANGsk{
\Question{
Uhol \( \beta \)  je o \( 20\,\% \) menší ako uhol \( \gamma \) a o \( 100\,\% \) väčší ako uhol \( \alpha \).  Uhol \( \alpha \) je}
{o \( 60\,\% \) menší než \( \gamma \).}{o \( 40\,\% \) menší ako \( \gamma \).}{o \( 20\,\% \) menší ako \( \gamma \).}{o \( 80\,\% \) menší ako \( \gamma \).}{}{}}
\LANGes{
\Question{
El ángulo \( \beta \)  es un \( 20\% \) menor que el ángulo \( \gamma \) y un \( 100\% \) mayor que el ángulo \( \alpha \).  El ángulo \( \alpha \) es}
{en un  \( 60\% \) menor que \( \gamma \).}{en un \( 40\% \) menor que \( \gamma \).}{en un \( 20\% \) menor que \( \gamma \).}{en un \( 80\% \) menor que \( \gamma \).}{}{}}
\End

\Data\ID{2010018201}
\Updated{2022-11-02 09:11:06}
\Level{B}
\SubArea{Percent problems}
\LANGen{
\Question{
We measured the length of a rod to be \(2\,\mathrm{m}\). What is the possible maximum length of the rod, if the percentage error of the measurement was \(5\%\).}
{\(210\,\mathrm{cm}\)}{\(205\,\mathrm{cm}\)}{\(200\,\mathrm{cm}\)}{\(195\,\mathrm{cm}\)}{}{}}
\LANGpl{
\Question{
Długość pręta to \(2\,\mathrm{m}\). Jaka jest możliwa maksymalna długość pręta, jeśli błąd procentowy pomiaru wynosił \(5\%\)?}
{\(210\,\mathrm{cm}\)}{\(205\,\mathrm{cm}\)}{\(200\,\mathrm{cm}\)}{\(195\,\mathrm{cm}\)}{}{}}
\LANGcs{
\Question{
Měřením jsme zjistili délku tyče \(2\,\mathrm{m}\). Jakou maximální délku tyče musíme připustit, jestliže měření délky proběhlo s relativní chybou měření \(5\%\)?}
{\(210\,\mathrm{cm}\)}{\(205\,\mathrm{cm}\)}{\(200\,\mathrm{cm}\)}{\(195\,\mathrm{cm}\)}{}{}}
\LANGsk{
\Question{
Meraním sme zistili dĺžku tyče \(2\,\mathrm{m}\). Aká je maximálna možná dĺžka tyče, ak bola percentuálna chyba merania \(5\%\)?}
{\(210\,\mathrm{cm}\)}{\(205\,\mathrm{cm}\)}{\(200\,\mathrm{cm}\)}{\(195\,\mathrm{cm}\)}{}{}}
\LANGes{
\Question{
Hemos medido la longitud de una varilla obteniendo \(2\,\mathrm{m}\). ¿Cuál es la posible longitud máxima de la varilla, si el porcentaje de error de la medición fue del \(5\%\)?}
{\(210\,\mathrm{cm}\)}{\(205\,\mathrm{cm}\)}{\(200\,\mathrm{cm}\)}{\(195\,\mathrm{cm}\)}{}{}}
\End

\Data\ID{2010018202}
\Updated{2022-11-02 09:11:06}
\Level{B}
\SubArea{Percent problems}
\LANGen{
\Question{
A rectangle is drawn in a scale drawing with the scale factor of \(1:5\). By what percentage is the area of the drawn rectangle smaller than the area of the real unscaled rectangle?}
{by \(96\%\)}{by \(4\%\)}{by \(20\%\)}{by \(80\%\)}{}{}}
\LANGpl{
\Question{
Prostokąt jest narysowany na rysunku w skali \(1:5\). O jaki procent powierzchnia narysowanego prostokąta jest mniejsza niż powierzchnia rzeczywistego prostokąta nieskalowanego?}
{o \(96\%\)}{o \(4\%\)}{o \(20\%\)}{o \(80\%\)}{}{}}
\LANGcs{
\Question{
Na plánu s měřítkem \(1:5\) je vyznačen obdélník. O kolik procent je obsah obdélníku na plánu menší než obsah skutečného obdélníku?}
{o \(96\,\%\)}{o \(4\,\%\)}{o \(20\,\%\)}{o \(80\,\%\)}{}{}}
\LANGsk{
\Question{
Na pláne s mierkou \(1:5\) je vyznačený obdĺžnik. O koľko percent je obsah obdĺžnika na pláne menší ako obsah skutočného obdĺžnika?}
{o \(96\%\)}{o \(4\%\)}{o \(20\%\)}{o \(80\%\)}{}{}}
\LANGes{
\Question{
Se dibuja un rectángulo a escala \(1:5\). ¿En qué porcentaje el área del rectángulo dibujado es menor que el área del rectángulo real?}
{en un \(96\%\)}{en un \(4\%\)}{en un \(20\%\)}{en un \(80\%\)}{}{}}
\End

\Data\ID{2010018203}
\Updated{2022-11-02 09:11:06}
\Level{B}
\SubArea{Percent problems}
\LANGen{
\Question{
A protective layer of thickness \(d\) reduces the level of harmful radiation by \(10\%\). Determine what is the percentage decrease in the original level of harmful radiation after passing through the layer of thickness \(3d\). Round the result to the whole percent.}
{\(73\%\)}{\(70\%\)}{\(30\%\)}{\(27\%\)}{}{}}
\LANGpl{
\Question{
Gruba warstwa ochronna \(d\) zmniejsza poziom szkodliwego promieniowania o \(10\%\). Określ, jaki jest procentowy spadek pierwotnego poziomu szkodliwego promieniowania po przejściu przez warstwę grubości \(3d\). Zaokrąglij wynik do pełnego procentu.}
{\(73\%\)}{\(70\%\)}{\(30\%\)}{\(27\%\)}{}{}}
\LANGcs{
\Question{
Ochranná vrstva o tloušťce \(d\) sníží úroveň škodlivého záření o \(10\%\). Určete, na kolik procent z původní hodnoty klesne úroveň škodlivého záření po průchodu vrstvou o tloušťce \(3d\). Výsledek zaokrouhlete na celá procenta.}
{\(73\%\)}{\(70\%\)}{\(30\%\)}{\(27\%\)}{}{}}
\LANGsk{
\Question{
Ochranná vrstva hrúbky \(d\) zníží úroveň škodlivého žiarenia o \(10\%\). Určte, na koľko percent z pôvodnej hodnoty klesne úroveň škodlivého žiarenia po prechode vrstvou hrúbky \(3d\). Výsledok zaokrúhlite na celé percentá.}
{\(73\%\)}{\(70\%\)}{\(30\%\)}{\(27\%\)}{}{}}
\LANGes{
\Question{
Una capa protectora de espesor \(d\) reduce el nivel de radiación nociva en un \(10\%\). Determina cuál es el porcentaje de disminución del nivel original de radiación nociva que llega después de atravesar una capa de espesor \(3d\). Redondea el resultado al porcentaje entero.}
{\(73\%\)}{\(70\%\)}{\(30\%\)}{\(27\%\)}{}{}}
\End

\Data\ID{2010018204}
\Updated{2024-02-02 09:02:37}
\Level{B}
\SubArea{Percent problems}
\LANGen{
\Question{
The aluminium rod and brass rod have the same length at a given temperature. The material constants of the rods are: \(\alpha_{\mathrm{aluminium}}=24\cdot 10^{-6}\,\mathrm{K}^{-1}\) and \(\alpha_{\mathrm{brass}}=18\cdot 10^{-6}\,\mathrm{K}^{-1}\). Suppose both rods are heated to the same higher temperature. Find the true statement about the extensions of both rods. Round the percentage difference in the rods’ extensions to the whole percent.
\[~\]
Hint: Solid materials expand upon heating. The rod with the initial length \(l_0\) is upon the temperature increasing by \(\Delta t\) extended by the value \(\Delta l = l_0 \cdot \alpha \cdot \Delta t\), where \(\alpha\) is the material constant (coefficient of linear thermal expansion) that is indicative of the extent to which a material expands upon heating.}
{The extension of the aluminium rod is by \(33\%\) greater than the extension of the brass rod.}{The extension of the aluminium rod is by \(67\%\) greater than the extension of the brass rod.}{The extension of the aluminium rod is by \(133\%\) greater than the extension of the brass rod.}{The extension of the aluminium rod is by \(33\%\) less than the extension of the brass rod.}{}{}}
\LANGpl{
\Question{
Pręt aluminiowy i pręt mosiężny mają tę samą długość w danej temperaturze. Stałe materiałowe prętów to: \(\alpha_{\mathrm{aluminium}}=24\cdot 10^{-6}\,\mathrm{K}^{-1}\) i  \(\alpha_{\mathrm{brass}}=18\cdot 10^{-6}\,\mathrm{K}^{-1}\). Załóżmy, że oba pręty są podgrzewane do tej samej wyższej temperatury. Znajdź prawdziwe stwierdzenie dotyczące przedłużeń obu prętów. Zaokrąglij procentową różnicę w wysunięciach prętów do pełnego procentu.
\[~\]
Wskazówka: Materiały stałe rozszerzają się po podgrzaniu. Pręt o długości początkowej \(l_0\) jest pod wpływem wzrostu temperatury o \(\Delta t\) przedłużonej o wartość \(\Delta l = l_0 \cdot \alpha \cdot \Delta t\), gdzie \(\alpha\) to stała materiałowa (współczynnik liniowej rozszerzalności cieplnej), która wskazuje na stopień, w jakim materiał rozszerza się po podgrzaniu.}
{Wydłużenie pręta aluminiowego jest o \(33\%\) większe niż przedłużenie pręta mosiężnego.}{Wydłużenie pręta aluminiowego jest o \(67\%\) większe niż przedłużenie pręta mosiężnego.}{Wydłużenie  pręta aluminiowego jest o \(133\%\) większe niż przedłużenie pręta mosiężnego.}{Wydłużenie pręta aluminiowego jest o \(33\%\) mniejsze niż przedłużenie pręta mosiężnego.}{}{}}
\LANGcs{
\Question{
Hliníková a mosazná tyč mají při dané teplotě stejnou délku. Známe materiálové konstanty obou tyčí: \(\alpha_{\text{hliník}}=24\cdot 10^{-6}\,\mathrm{K}^{-1}\) a \(\alpha_{\text{mosaz}}=18\cdot 10^{-6}\,\mathrm{K}^{-1}\). Rozhodněte, co platí o prodloužení obou tyčí, když je zahřejeme o stejnou teplotu. Procentuální rozdíl v prodloužení tyčí zaokrouhlete na celá procenta.
\[~\]
Nápověda: Při zahřívání se tělesa prodlužují. Při počáteční délce tyče \(l_0\) a při zahřátí o teplotu \(\Delta t\) se tyč prodlouží o hodnotu \(\Delta l = l_0 \cdot \alpha \cdot \Delta t\), kde \(\alpha\) je materiálová konstanta (součinitel teplotní délkové roztažnosti).}
{Prodloužení hliníkové tyče bude o \(33\%\) větší, než prodloužení mosazné tyče.}{Prodloužení hliníkové tyče bude o \(67\%\) větší, než prodloužení mosazné tyče.}{Prodloužení hliníkové tyče bude o \(133\%\) větší, než prodloužení mosazné tyče.}{Prodloužení hliníkové tyče bude o \(33\%\) menší, než prodloužení mosazné tyče}{}{}}
\LANGsk{
\Question{
Hliníková a mosadzná tyč majú pri danej teplote rovnakú dĺžku. Materiálové konštanty obidvoch tyčí sú: \(\alpha_{\text{hliník}}=24\cdot 10^{-6}\,\mathrm{K}^{-1}\) a \(\alpha_{\mathrm{mosaz}}=18\cdot 10^{-6}\,\mathrm{K}^{-1}\). Rozhodnite, čo platí o predĺžení oboch tyčí, ak ich zohrejeme o rovnakú teplotu. Percentuálny rozdiel v predĺžení tyčí zaokrúhlite na celé percentá.
\[~\]
Pomôcka: Pri zohrievaní sa telesá predlžujú. Pri začiatočnej dĺžke tyče \(l_0\) a pri zohrievaní o teplotu \(\Delta t\) sa tyč predĺži o hodnotu \(\Delta l = l_0 \cdot \alpha \cdot \Delta t\), kde \(\alpha\) je materiálová konštanta (súčiniteľ teplotnej dĺžkovej rozťažnosti).}
{Predĺženie hliníkovej tyče bude o \(33\%\) väčšie, ako predĺženie mosadznej tyče.}{Predĺženie hliníkovej tyče bude o \(67\%\) väčšie, ako predĺženie mosadznej tyče.}{Predĺženie hliníkovej tyče bude o \(133\%\) väčšie, ako predĺženie mosadznej tyče.}{Predĺženie hliníkovej tyče bude o \(33\%\) menšie, ako predĺženie mosadznej tyče.}{}{}}
\LANGes{
\Question{
Una varilla de aluminio y una de latón tienen la misma longitud para una determinada temperatura. Las constantes de los materiales son: \(\alpha_{\mathrm{aluminium}}=24\cdot 10^{-6}\,\mathrm{K}^{-1}\) y \(\alpha_{\mathrm{brass}}=18\cdot 10^{-6}\,\mathrm{K}^{-1}\). Supongamos que ambas varillas se calientan hasta la misma temperatura. Elige la afirmación correcta en cuanto a la dilatación de ambas varillas se refiere. Redondea la diferencia de porcentaje de dilatación en las varillas a las unidades.
\[~\]
Pista: Los materiales sólidos se dilatan con el calor. La dilatación de una varilla, \(\Delta l \) viene dada por la expresión \(\Delta l = l_0 \cdot \alpha \cdot \Delta t\), siendo \(l_0\) la longitud inicial de la varilla, \(\Delta t\) la variación de temperatura y \(\alpha\) la constante del material (constante de dilatación térmica lineal), que es un indicativo del alcance que tiene la dilatación en un material debido al calor.}
{La dilatación de la varilla de aluminio es un \(33\%\) mayor que la de latón.}{La dilatación de la varilla de aluminio es un \(67\%\) mayor que la de latón.}{La dilatación de la varilla de aluminio es un \(133\%\) mayor que la de latón.}{La dilatación de la varilla de aluminio es un \(33\%\) menor que la de latón.}{}{}}
\End

\Data\ID{9000078901}
\Updated{2022-04-23 05:04:26}
\Level{B}
\SubArea{Percent problems}
\LANGen{
\Question{
If we increase an unknown number \(x\) 
by \(5\, \%\), we
get \(378\).
Find \(x\).}
{\(360\)}{\(359.1\)}{\(396.9\)}{\(350\)}{}{}}
\LANGpl{
\Question{
Jeśli zwiększymy niewiadomą  \(x\) 
o \(5\, \%\), otrzymamy 
\(378\).
Znajdź \(x\).}
{\(360\)}{\(359.1\)}{\(396.9\)}{\(350\)}{}{}}
\LANGcs{
\Question{
Zvětšíme-li neznámé číslo o
\(5\, \%\), dostaneme
číslo \(378\).
Určete neznámé číslo.}
{\(360\)}{\(359{,}1\)}{\(396{,}9\)}{\(350\)}{}{}}
\LANGsk{
\Question{
Ak zväčšíme neznáme číslo o
\(5\, \%\), dostaneme
číslo \(378\).
Určte neznáme číslo.}
{\(360\)}{\(359{,}1\)}{\(396{,}9\)}{\(350\)}{}{}}
\LANGes{
\Question{
Si aumentamos un número desconocido  \(x\) 
en un  \(5\, \%\), obtenemos \(378\).
Calcula \(x\).}
{\(360\)}{\(359.1\)}{\(396.9\)}{\(350\)}{}{}}
\End

\Data\ID{9000078902}
\Updated{2020-07-15 03:07:37}
\Level{B}
\SubArea{Percent problems}
\LANGen{
\Question{
If we decrease an unknown number \(x\) 
by \(14\, \%\), we
get \(602\).
Find \(x\).}
{\(700\)}{\(686.28\)}{\(517.72\)}{\(680\)}{}{}}
\LANGpl{
\Question{
Jeśli zmniejszymy niewiadomą  \(x\) 
o \(14\, \%\), otrzymamy
\(602\).
Znajdź \(x\).}
{\(700\)}{\(686.28\)}{\(517.72\)}{\(680\)}{}{}}
\LANGcs{
\Question{
Zmenšíme-li neznámé číslo o \(14\, \%\),
dostaneme číslo \(602\).
Určete neznámé číslo.}
{\(700\)}{\(686{,}28\)}{\(517{,}72\)}{\(680\)}{}{}}
\LANGsk{
\Question{
Ak zmenšíme neznáme číslo o \(14\, \%\),
dostaneme číslo \(602\).
Určte neznáme číslo.}
{\(700\)}{\(686{,}28\)}{\(517{,}72\)}{\(680\)}{}{}}
\LANGes{
\Question{
Si disminuimos un número desconocido  \(x\) 
en un \(14\, \%\), obtenemos \(602\).
Calcula \(x\).}
{\(700\)}{\(686.28\)}{\(517.72\)}{\(680\)}{}{}}
\End

\Data\ID{9000078903}
\Updated{2020-07-15 03:07:37}
\Level{B}
\SubArea{Percent problems}
\LANGen{
\Question{
The number \(234\) 
is by \(20\, \%\) bigger
than \(x\).
Find \(x\).}
{\(195\)}{\(187.2\)}{\(280.8\)}{\(205\)}{}{}}
\LANGpl{
\Question{
Liczba \(234\) 
jest o \(20\, \%\) większa
niż \(x\).
Znajdź \(x\).}
{\(195\)}{\(187.2\)}{\(280.8\)}{\(205\)}{}{}}
\LANGcs{
\Question{
Číslo \(234\) 
je o \(20\, \%\) 
větší než neznámé číslo. Určete neznámé číslo.}
{\(195\)}{\(187{,}2\)}{\(280{,}8\)}{\(205\)}{}{}}
\LANGsk{
\Question{
Číslo \(234\) 
je o \(20\, \%\) 
väčšie než neznáme číslo. Určte neznáme číslo.}
{\(195\)}{\(187{,}2\)}{\(280{,}8\)}{\(205\)}{}{}}
\LANGes{
\Question{
El número \(234\) 
es un \(20\, \%\) más grande que \(x\).
Calcula \(x\).}
{\(195\)}{\(187.2\)}{\(280.8\)}{\(205\)}{}{}}
\End

\Data\ID{9000078904}
\Updated{2020-07-15 03:07:37}
\Level{B}
\SubArea{Percent problems}
\LANGen{
\Question{
The number \(87.5\) 
is a \(35\, \%\) from an
unknown number \(x\).
Find \(x\).}
{\(250\)}{\(240\)}{\(260\)}{\(270\)}{}{}}
\LANGpl{
\Question{
Liczba  \(87.5\) 
stanowi \(35\, \%\) niewiadomej \(x\).
Znajdź \(x\).}
{\(250\)}{\(240\)}{\(260\)}{\(270\)}{}{}}
\LANGcs{
\Question{
\(35\, \%\) z neznámého
čísla je \(87{,}5\).
Určete neznámé číslo.}
{\(250\)}{\(240\)}{\(260\)}{\(270\)}{}{}}
\LANGsk{
\Question{
\(35\, \%\) z neznámeho
čísla je \(87{,}5\).
Určte neznáme číslo.}
{\(250\)}{\(240\)}{\(260\)}{\(270\)}{}{}}
\LANGes{
\Question{
El número  \(87.5\) 
es un \(35\, \%\) de un número desconocido \(x\).
Calcula  \(x\).}
{\(250\)}{\(240\)}{\(260\)}{\(270\)}{}{}}
\End

\Data\ID{9000078905}
\Updated{2020-07-15 03:07:37}
\Level{B}
\SubArea{Percent problems}
\LANGen{
\Question{
A new laptop has been introduced to the marked with the original price
\(\$1\: 090\).
Soon after the introduction the price has been increased temporarily by
\(20\, \%\) and after several
months decreased by \(30\, \%\).
Find the final price for the laptop. Round your result to the nearest integer.}
{\(\$916\)}{\(\$981\)}{\(\$952\)}{\(\$930\)}{}{}}
\LANGpl{
\Question{
Nowy laptop został wprowadzony na rynek w początkowej cenie
\(\$1\: 090\).
Niedługo potem początkowa cena została chwilowo zmniejszona o \(20\, \%\), a po kilku miesiącach zmniejszona o 
 \(30\, \%\).
Jaka była końcowa cena laptopa? Zaokrąglij do tysięcy.}
{\(\$916\)}{\(\$981\)}{\(\$952\)}{\(\$930\)}{}{}}
\LANGcs{
\Question{
Původní cena šatů byla \(1\: 090\) 
Kč. Potom je zdražili o \(20\, \%\) 
a za měsíc zlevnili o \(30\, \%\).
Určete konečnou cenu šatů. (Výsledek zaokrouhlete na koruny.)}
{\(916\) Kč}{\(981\) Kč}{\(952\) Kč}{\(930\) Kč}{}{}}
\LANGsk{
\Question{
Pôvodná cena šiat bola \(1\: 090\) 
Kč. Potom zdraželi o \(20\, \%\) 
a za mesiac zlacneli o \(30\, \%\).
Určte konečnú cenu šiat. (Výsledok zaokrúhlite na koruny.)}
{\(916\) Kč}{\(981\) Kč}{\(952\) Kč}{\(930\) Kč}{}{}}
\LANGes{
\Question{
Se ha presentado un nuevo ordenador portátil marcado con el precio original de 
\(\$1\: 090\).
Poco después de la presentación, el precio aumentó temporalmente en un 
\(20\, \%\) y, después de varios meses, disminuyó  un \(30\, \%\).
Calcula el precio final del ordenador portátil. Redondea tu resultado al entero más cercano.}
{\(\$916\)}{\(\$981\)}{\(\$952\)}{\(\$930\)}{}{}}
\End

\Data\ID{9000078906}
\Updated{2020-07-15 03:07:37}
\Level{B}
\SubArea{Percent problems}
\LANGen{
\Question{
The price of a ship decreased in autumn by
\(16\, \%\) and in spring
increased by \(4\, \%\). This
final price was \(\$367\: 000\).
Find the original price of the ship. (Round to the nearest thousand.)}
{\(\$420\: 000\)}{\(\$417\: 000\)}{\(\$409\: 000\)}{\(\$415\: 000\)}{}{}}
\LANGpl{
\Question{
Cena statku zmniejszyła się jesienią o
\(16\, \%\), a wiosną wzrosła o
 \(4\, \%\). Ostateczna cena wyniosła
 \(\$367\: 000\).
Znajdź początkową cenę statku. Zaokrąglij do tysięcy.}
{\(\$420\: 000\)}{\(\$417\: 000\)}{\(\$409\: 000\)}{\(\$415\: 000\)}{}{}}
\LANGcs{
\Question{
Původní cena automobilu byla snížena o
\(16\, \%\) a později
zvýšena o \(4\, \%\) 
na \(367\: 000\) 
Kč. Určete původní cenu automobilu. (Výsledek zaokrouhlete na
tisícikoruny.)}
{\(420\: 000\) Kč}{\(417\: 000\) Kč}{\(409\: 000\) Kč}{\(415\: 000\) Kč}{}{}}
\LANGsk{
\Question{
Pôvodná cena automobilu bola znížená o
\(16\, \%\) a neskôr
zvýšená o \(4\, \%\) 
na \(367\: 000\) 
Kč. Určte pôvodnú cenu automobilu. (Výsledok zaokrúhlite na
tisíckoruny.)}
{\(420\: 000\) Kč}{\(417\: 000\) Kč}{\(409\: 000\) Kč}{\(415\: 000\) Kč}{}{}}
\LANGes{
\Question{
El precio de un barco disminuyó en otoño un 
\(16\, \%\) y en primavera aumentó en un \(4\, \%\). Este precio final fue de \(\$367\: 000\).
Calcula el precio original del barco. (Redondea a los millares más cercanos)}
{\(\$420\: 000\)}{\(\$417\: 000\)}{\(\$409\: 000\)}{\(\$415\: 000\)}{}{}}
\End

\Data\ID{9000078907}
\Updated{2020-07-15 03:07:37}
\Level{B}
\SubArea{Percent problems}
\LANGen{
\Question{
The unknown number \(x\) 
has been increased by \(20\, \%\) 
and then decreased by \(5\, \%\).
The final value was \(513\).
Find \(x\).
(Round to the closest integer.)}
{\(450\)}{\(446\)}{\(430\)}{\(436\)}{}{}}
\LANGpl{
\Question{
Niewiadoma \(x\) 
została zwiększona o  \(20\, \%\),
a następnie zmniejszona o \(5\, \%\).
Ostateczna wartość wyniosła \(513\).
Znajdź \(x\).
(Zaokrąglij do liczby całkowitej.)}
{\(450\)}{\(446\)}{\(430\)}{\(436\)}{}{}}
\LANGcs{
\Question{
Zvětšíme-li neznámé číslo o
\(20\, \%\) a potom ho
zmenšíme o \(5\, \%\),
dostaneme \(513\).
Určete neznámé číslo. (Neznámé číslo zaokrouhlete na jednotky.)}
{\(450\)}{\(446\)}{\(430\)}{\(436\)}{}{}}
\LANGsk{
\Question{
Ak zväčšíme neznáme číslo o
\(20\, \%\) a potom ho
zmenšíme o \(5\, \%\),
dostaneme \(513\).
Určte neznáme číslo. (Neznáme číslo zaokrúhlite na jednotky.)}
{\(450\)}{\(446\)}{\(430\)}{\(436\)}{}{}}
\LANGes{
\Question{
El número desconocido \(x\) 
aumentó en un \(20\, \%\) 
y luego disminuyó un  \(5\, \%\).
El valor final fue  \(513\).
Calcula \(x\).
(Redondea al entero más cercano).}
{\(450\)}{\(446\)}{\(430\)}{\(436\)}{}{}}
\End

\Data\ID{9000078908}
\Updated{2020-07-15 03:07:37}
\Level{B}
\SubArea{Percent problems}
\LANGen{
\Question{
A house price in the March was \(\$259\: 000\).
The house has not been sold for several months and the price dropped to
\(\$234\: 000\) in
August. Write the discount as percents of the original price. (Round to the tenth.)}
{\(9.7\, \%\)}{\(10.7\, \%\)}{\(10.5\, \%\)}{\(9.5\, \%\)}{}{}}
\LANGpl{
\Question{
Cena za dom w marcu wyniosła \(\$259\: 000\).
Dom nie został sprzedany przez kilka miesięcy i cena spadła do
\(\$234\: 000\) w sierpniu.
Napisz, o ile obniżono początkową cenę w procentach. (Zaokrąglij do części tysięcznych.)}
{\(9.7\, \%\)}{\(10.7\, \%\)}{\(10.5\, \%\)}{\(9.5\, \%\)}{}{}}
\LANGcs{
\Question{
Na začátku roku se jistý automobil prodával za
\(259\: 000\) Kč. Na konci roku jeho
cena klesla na \(234\: 000\) Kč.
O kolik procent klesla cena automobilu? (Výsledek zaokrouhlete na desetiny
procenta.)}
{o \(9{,}7\, \%\)}{o \(10{,}7\, \%\)}{o \(10{,}5\, \%\)}{o \(9{,}5\, \%\)}{}{}}
\LANGsk{
\Question{
Na začiatku roku sa jeden automobil predával za
\(259\: 000\) Kč. Na konci roku jeho
cena klesla na \(234\: 000\) Kč.
O koľko percent klesla cena automobilu? (Výsledok zaokrúhlite na desatiny
procenta.)}
{o \(9{,}7\, \%\)}{o \(10{,}7\, \%\)}{o \(10{,}5\, \%\)}{o \(9{,}5\, \%\)}{}{}}
\LANGes{
\Question{
El precio de una casa en marzo fue de  \(\$259\: 000\).
La casa no se ha vendido durante varios meses y el precio bajó a 
\(\$234\: 000\) en agosto. ¿Cuál es el descenso porcentual del precio principal? (Redondea a la décima.)}
{\(9.7\, \%\)}{\(10.7\, \%\)}{\(10.5\, \%\)}{\(9.5\, \%\)}{}{}}
\End

\Data\ID{9000078909}
\Updated{2020-07-15 03:07:37}
\Level{B}
\SubArea{Percent problems}
\LANGen{
\Question{
The salary of an employee is decreased by the tax at a flat rate
\(14.5\, \%\). An employee
paid the tax \(\$3\: 683\).
Find his/her take home pay (i.e. the pre-tax income decreased by the tax).}
{\(\$21\: 717\)}{\(\$25\: 400\)}{\(\$21\: 971\)}{\(\$23\: 352\)}{}{}}
\LANGpl{
\Question{
Pensja pracownika została pomniejszona o podatek od pensji stałej, który wyniósł
\(14.5\, \%\). Pracownik zapłacił podatek w wysokości 
 \(\$3\: 683\).
Znajdź wysokość jego pensji stałej (przed odprowadzeniem podatku).}
{\(\$21\: 717\)}{\(\$25\: 400\)}{\(\$21\: 971\)}{\(\$23\: 352\)}{}{}}
\LANGcs{
\Question{
Z hrubé mzdy byly pracovníkovi odečteny zákonné odvody
\(3\: 683\) Kč, což
představovalo \(14{,}5\, \%\) 
jeho hrubé mzdy. Jak velká částka (čistá mzda) mu byla vyplacena?
(Poznámka: Zaměstnancům se vyplácí tzv. čistá mzda, která je rozdílem
hrubé mzdy a zákonných odvodů -- daň z příjmu, sociální a zdravotní
pojištění.)}
{\(21\: 717\) Kč}{\(25\: 400\) Kč}{\(21\: 971\) Kč}{\(23\: 352\) Kč}{}{}}
\LANGsk{
\Question{
Z hrubej mzdy boli pracovníkovi odrátané zákonné odvody
\(3\: 683\) Kč, čo
predstavovalo \(14{,}5\, \%\) 
jeho hrubej mzdy. Aká veľká čiastka (čistá mzda) mu bola vyplatená?
(Poznámka: Zamestnancom sa vypláca tzv. čistá mzda, ktorá je rozdielom
hrubej mzdy a zákonných odvodov -- daň z príjmu, sociálne a zdravotné
poistenie.)}
{\(21\: 717\) Kč}{\(25\: 400\) Kč}{\(21\: 971\) Kč}{\(23\: 352\) Kč}{}{}}
\LANGes{
\Question{
El salario de un empleado fue reducido por el impuesto de la tasa fija 
\(14.5\, \%\). El empleado pagó el impuesto de \(\$3\: 683\).
Calcula su salario neto (es decir, el ingreso antes de impuestos reducido por el impuesto).}
{\(\$21\: 717\)}{\(\$25\: 400\)}{\(\$21\: 971\)}{\(\$23\: 352\)}{}{}}
\End

\Data\ID{9000078910}
\Updated{2022-04-23 05:04:26}
\Level{B}
\SubArea{Percent problems}
\LANGen{
\Question{
There is \(19\) 
girls and \(12\) 
boys in the class. Find the difference between the percentage of the girls and boys in
this class.}
{\(22.6\, \%\)}{\(15.1\, \%\)}{\(18.5\, \%\)}{\(23.5\, \%\)}{}{}}
\LANGpl{
\Question{
W klasie jest \(19\) 
dziewcząt i \(12\) 
chłopców. Znajdź różnicę w procentach pomiędzy ilością dziewcząt i chłopców w tej klasie.}
{\(22.6\, \%\)}{\(15.1\, \%\)}{\(18.5\, \%\)}{\(23.5\, \%\)}{}{}}
\LANGcs{
\Question{
Ve třídě je \(19\) 
dívek a \(12\) 
chlapců. O kolik se liší procentuální zastoupení dívek a chlapců ve
třídě?}
{o \(22{,}6\, \%\)}{o \(15{,}1\, \%\)}{o \(18{,}5\, \%\)}{o \(23{,}5\, \%\)}{}{}}
\LANGsk{
\Question{
V triede je \(19\) 
dievčat a \(12\) 
chlapcov. O koľko sa líši percentuálne zastúpenie dievčat a chlapcov v
triede?}
{o \(22{,}6\, \%\)}{o \(15{,}1\, \%\)}{o \(18{,}5\, \%\)}{o \(23{,}5\, \%\)}{}{}}
\LANGes{
\Question{
Hay \(19\) 
niñas y  \(12\) 
niños en la clase. Calcula la diferencia entre el porcentaje de niñas y niños en esta clase.}
{\(22.6\, \%\)}{\(15.1\, \%\)}{\(18.5\, \%\)}{\(23.5\, \%\)}{}{}}
\End

\Data\ID{9000085301}
\Updated{2020-08-22 04:08:42}
\Level{B}
\SubArea{Percent problems}
\LANGen{
\Question{
A bike was introduced to the market with the initial price
\(\$10\: 000\) in
April. As a consequence of bad sales values, the price has been decreased by
\(20\, \%\) in June and
by another \(20\, \%\) 
in September. Find the price in September.}
{\(\$6\: 400\)}{\(\$6\: 000\)}{\(\$6\: 125\)}{\(\$6\: 500\)}{}{}}
\LANGpl{
\Question{
Motor został wprowadzony na rynek w kwietniu w początkowej cenie \(\$10\: 000\). Ze względu na niski wskaźnik sprzedaży jego cena została zmniejszona o 
\(20\, \%\) w czerwcu i o kolejne \(20\, \%\) we wrześniu. Jaka była cena motoru we wrześniu?}
{\(\$6\: 400\)}{\(\$6\: 000\)}{\(\$6\: 125\)}{\(\$6\: 500\)}{}{}}
\LANGcs{
\Question{
Televizor stál \(10\: 000\, \text{Kč}\).
Postupně byl dvakrát zlevněn a to vždy o
\(20\, \%\), aby
se lépe prodával. Jaká je jeho konečná prodejní cena?}
{\(6\: 400\, \text{Kč}\)}{\(6\: 000\, \text{Kč}\)}{\(6\: 125\, \text{Kč}\)}{\(6\: 500\, \text{Kč}\)}{}{}}
\LANGsk{
\Question{
Televízor stál \(10\: 000\, \text{Kč}\).
Postupne bol dvakrát zlacnený a to vždy o
\(20\, \%\), aby
sa lepšie predával. Aká je jeho konečná predajná cena?}
{\(6\: 400\, \text{Kč}\)}{\(6\: 000\, \text{Kč}\)}{\(6\: 125\, \text{Kč}\)}{\(6\: 500\, \text{Kč}\)}{}{}}
\LANGes{
\Question{
Una bicicleta fue presentada en el mercado con el precio inicial de 
\(\$10\: 000\) en abril. Su precio disminuyó un
\(20\, \%\)  en junio y otro 
 \(20\, \%\) 
en septiembre para que se vendiera más. Calcula el precio en septiembre.}
{\(\$6\: 400\)}{\(\$6\: 000\)}{\(\$6\: 125\)}{\(\$6\: 500\)}{}{}}
\End

\Data\ID{9000085302}
\Updated{2020-07-15 03:07:37}
\Level{B}
\SubArea{Percent problems}
\LANGen{
\Question{
There are \(32\) students
in the class (\(20\) 
boys and \(12\) 
girls). One quarter of the boys and one quarter of the girls are excellent students.
Suppose that one boy and one girl, both excellent students, move to another school.
Find the difference in the ratio of excellent students in the class before and after the
change in the class.}
{\(5\, \%\)}{\(6.25\, \%\)}{\(7.5\, \%\)}{\(8.25\, \%\)}{}{}}
\LANGpl{
\Question{
W klasie jest \(32\) uczniów (\(20\) 
chłopców i  \(12\) dziewcząt). Jedna czwarta chłopców i jedna czwarta dziewcząt są świetnymi uczniami. Załóżmy, że 1 chłopiec i 1 dziewczyna, obydwoje świetni uczniowie, przeniosą się do innej szkoły. Jaka będzie różnica we wskaźniku wspaniałych uczniów przed zmianą i po zmianie w tej klasie?}
{\(5\, \%\)}{\(6.25\, \%\)}{\(7.5\, \%\)}{\(8.25\, \%\)}{}{}}
\LANGcs{
\Question{
Ve třídě je \(32\) 
žáků - \(20\) 
chlapců a \(12\) 
dívek. Čtvrtina všech chlapců a čtvrtina všech dívek má vyznamenání.
O kolik procent klesne počet všech vyznamenaných ve třídě, jestliže jeden
chlapec a jedna dívka se samými jedničkami přestoupí na jinou školu?}
{\(5\, \%\)}{\(6{,}25\, \%\)}{\(7{,}5\, \%\)}{\(8{,}25\, \%\)}{}{}}
\LANGsk{
\Question{
V triede je \(32\) 
žiakov - \(20\) 
chlapcov a \(12\) 
dievčat. Štvrtina všetkých chlapcov a štvrtina všetkých dievčat má vyznamenanie.
O koľko percent klesne počet všetkých vyznamenaných v triede, keď jeden
chlapec a jedno dievča so samými jednotkami prestúpi na inú školu?}
{\(5\, \%\)}{\(6{,}25\, \%\)}{\(7{,}5\, \%\)}{\(8{,}25\, \%\)}{}{}}
\LANGes{
\Question{
Hay \(32\) estudiantes en la clase  (\(20\) 
niños y  \(12\) 
niñas). Un cuarto de todos los niños y un cuarto de todas las niñas son excelentes estudiantes. Supongamos que un niño y una niña, ambos excelentes estudiantes, se van a otra escuela. Calcula la diferencia en la proporción de los estudiantes excelentes en la clase antes y después del cambio en la clase.}
{\(5\, \%\)}{\(6.25\, \%\)}{\(7.5\, \%\)}{\(8.25\, \%\)}{}{}}
\End

\Data\ID{9000085303}
\Updated{2020-07-15 03:07:37}
\Level{B}
\SubArea{Percent problems}
\LANGen{
\Question{
The shop received \(30\) 
pieces of TV from the manufacturer NoName. Five TVs in this delivery
were defect. The shop also received TVs from the manufacturer
GoodName. All these TVs were perfect. The final inspection revealed that
\(10\, \%\) of
TVs were defect in total. Use these information to find out how many TVs have been
supplied from the producer GoodName.}
{\(20\)}{\(25\)}{\(18\)}{\(16\)}{}{}}
\LANGpl{
\Question{
Sklep otrzymał \(30\) sztuk telewizorów od producenta NoName. W przesyłce znalazło się 5 wadliwych telewizorów. Sklep otrzymał również telewizory od producenta GoodName. Wszystkie telewizory były w idealnym stanie. Ostateczna inspekcja dowiodła, że  \(10\, \%\) wszystkich telewizorów stanowiły te wadliwe. Na podstawie podanych informacji powiedz ile telewizorów dostarczył producent GoodName.}
{\(20\)}{\(25\)}{\(18\)}{\(16\)}{}{}}
\LANGcs{
\Question{
Do obchodu bylo dodáno \(30\) 
kusů výrobků od výrobce \(A\),
přičemž \(5\) 
z nich nefungovalo, a určité množství výrobků od výrobce
\(B\),
které fungovaly všechny. Kolik výrobků dodal výrobce
\(B\), jestliže
\(10\, \%\) ze
všech výrobků bylo nefunkčních?}
{\(20\)}{\(25\)}{\(18\)}{\(16\)}{}{}}
\LANGsk{
\Question{
Do obchodu bolo dodaných \(30\) 
kusov výrobkov od výrobcu \(A\),
pričom \(5\) 
z nich nefungovalo, a určité množstvo výrobkov od výrobcu
\(B\),
ktoré fungovali všetky. Koľko výrobkov dodal výrobca
\(B\), ak
\(10\, \%\) zo
všetkých výrobkov bolo nefunkčných?}
{\(20\)}{\(25\)}{\(18\)}{\(16\)}{}{}}
\LANGes{
\Question{
Una tienda recibió  \(30\) 
productos de fabricante A, pero cinco televisores de esta entrega estaban defectuosos. La tienda también recibió televisores del fabricante B cuyos productos eran perfectos. La inspección final reveló que el \(10\, \%\) de los productos tenían defectos en total. ¿Cuántos productos se han suministrado del fabricante B?}
{\(20\)}{\(25\)}{\(18\)}{\(16\)}{}{}}
\End

\Data\ID{9000085304}
\Updated{2020-07-15 03:07:37}
\Level{B}
\SubArea{Percent problems}
\LANGen{
\Question{
The final score of the hockey match between teams
\(A\) and
\(B\) was
\(2 : 2\). The goaltender
of \(A\) had
success rate \(90\, \%\) 
(he stopped \(90\, \%\) 
of all shots on goal) and the goaltender of the team
\(B\) had fail
rate \(20\, \%\) 
(\(20\, \%\) of all
shots on goal went into his net). How many shots on goal has been in the match in
total?}
{\(30\)}{\(25\)}{\(35\)}{\(40\)}{}{}}
\LANGpl{
\Question{
Ostateczny wynik meczu hokeja pomiędzy drużynami
\(A\) i
\(B\) wyniósł
\(2 : 2\). Skuteczność bramkarza drużyny 
 \(A\) wyniosła \(90\, \%\) 
(obronił \(90\, \%\) 
wszystkich strzałów na bramkę), a nieskuteczność bramkarza drużyny
\(B\) wyniosła
 \(20\, \%\) 
(\(20\, \%\) wszystkich strzałów trafiło do bramki). Ile strzałów na bramkę w sumie oddano?}
{\(30\)}{\(25\)}{\(35\)}{\(40\)}{}{}}
\LANGcs{
\Question{
Hokejové utkání mezi mužstvy \(A\) 
a \(B\) skončilo
nerozhodně \(2 : 2\).
Brankář mužstva \(A\) 
chytil \(90\, \%\) 
všech střel vystřelených na jeho branku, brankář mužstva
\(B\) nechytil
\(20\, \%\) 
všech střel vystřelených na jeho branku. Kolik střel celkem bylo během
zápasu vystřeleno na obě branky?}
{\(30\)}{\(25\)}{\(35\)}{\(40\)}{}{}}
\LANGsk{
\Question{
Hokejové stretnutie medzi mužstvami \(A\) 
a \(B\) skončilo
nerozhodne \(2 : 2\).
Brankár mužstva \(A\) 
chytil \(90\, \%\) 
všetkých striel vystrelených na jeho bránku, brankár mužstva
\(B\) nechytil
\(20\, \%\) 
všetkých striel vystrelených na jeho bránku. Koľko striel celkom bolo behom
zápasu vystrelených na obe bránky?}
{\(30\)}{\(25\)}{\(35\)}{\(40\)}{}{}}
\LANGes{
\Question{
La puntuación final del partido de hockey entre los equipos 
\(A\) y
\(B\) fue
\(2 : 2\). El portero de 
 \(A\) tuvo una tasa de éxito del  \(90\, \%\) 
(detuvo el \(90\, \%\) 
de todos los tiros a puerta) y el portero del equipo 
\(B\) tuvo una tasa de fracaso del  \(20\, \%\) 
(\(20\, \%\) de todos los tiros a puerta se metió en su red). ¿Cuántos tiros a portería hubo en total en el partido?}
{\(30\)}{\(25\)}{\(35\)}{\(40\)}{}{}}
\End

\Data\ID{9000085305}
\Updated{2020-08-22 04:08:30}
\Level{B}
\SubArea{Percent problems}
\LANGen{
\Question{
The price for the first issue of the book was
\(\$100\), the price for the
second issue \(\$125\).
Find the rate of sale of the second issue if the target price for the second issue is the
same as for the first issue.}
{\(20\, \%\)}{\(22.5\, \%\)}{\(17.5\, \%\)}{\(25\, \%\)}{}{}}
\LANGpl{
\Question{
Cena za pierwsze wydanie książki wyniosła
\(\$100\), a za drugie wydanie \(\$125\). O ile procent należy obniżyć cenę drugiego wydania, aby otrzymać cenę pierwszego wydania?}
{\(20\, \%\)}{\(22{,}5\, \%\)}{\(17{,}5\, \%\)}{\(25\, \%\)}{}{}}
\LANGcs{
\Question{
První vydání učebnice stálo \(100\, \text{Kč}\),
druhé vydání téže učebnice \(125\, \text{Kč}\).
O kolik procent je potřeba zlevnit druhé vydání, aby stálo tolik, co první?}
{\(20\, \%\)}{\(22{,}5\, \%\)}{\(17{,}5\, \%\)}{\(25\, \%\)}{}{}}
\LANGsk{
\Question{
Prvé vydanie učebnice stálo \(100\, \text{Kč}\),
druhé vydanie tej istej učebnice \(125\, \text{Kč}\).
O koľko percent je treba zlacniť druhé vydanie, aby stálo toľko, ako prvé?}
{\(20\, \%\)}{\(22{,}5\, \%\)}{\(17{,}5\, \%\)}{\(25\, \%\)}{}{}}
\LANGes{
\Question{
El precio de la primera edición de un libro fue de 
\(\$100\), y el precio de la segunda edición fue de  \(\$125\).
¿En qué porcentaje es necesario abaratar la segunda edición para que cueste lo mismo que la primera?}
{\(20\, \%\)}{\(22.5\, \%\)}{\(17.5\, \%\)}{\(25\, \%\)}{}{}}
\End

\Data\ID{9000085306}
\Updated{2020-08-22 04:08:03}
\Level{B}
\SubArea{Percent problems}
\LANGen{
\Question{
The ski has been sold for a reduced price in a sale after the season. The discount
\(18\, \%\) 
of the original price resumed in a price which was cheaper by
\(\$36\). Find
the original price of the ski.}
{\(\$200\)}{\(\$250\)}{\(\$450\)}{\(\$500\)}{}{}}
\LANGpl{
\Question{
Narty zostały sprzedane w obniżonej cenie podczas posezonowej wyprzedaży. Po rabacie w wysokości
\(18\, \%\) początkowej ceny były one tańsze o 
\(\$36\). Jaka była początkowa cena nart?}
{\(\$200\)}{\(\$250\)}{\(\$450\)}{\(\$500\)}{}{}}
\LANGcs{
\Question{
Sjezdové lyže byly po sezoně zlevněny o
\(18\, \%\) původní ceny
a prodávaly se o \(360\, \text{Kč}\) 
levněji. Jaká byla cena lyží před slevou?}
{\(2\: 000\, \text{Kč}\)}{\(2\: 500\, \text{Kč}\)}{\(4\: 500\, \text{Kč}\)}{\(5\: 000\, \text{Kč}\)}{}{}}
\LANGsk{
\Question{
Zjazdové lyže boli po sezóne zlacnené o
\(18\, \%\) pôvodnej ceny
a predávali sa o \(360\, \text{Kč}\) 
lacnejšie. Aká bola cena lyží pred zľavou?}
{\(2\: 000\, \text{Kč}\)}{\(2\: 500\, \text{Kč}\)}{\(4\: 500\, \text{Kč}\)}{\(5\: 000\, \text{Kč}\)}{}{}}
\LANGes{
\Question{
Después de la temporada, un esquí se ha vendido por un precio reducido. Un descuento del 
\(18\, \%\) 
del precio inicial resultó 
\(\$36\) de descuento en el precio. Calcula el precio inicial del esquí.}
{\(\$200\)}{\(\$250\)}{\(\$450\)}{\(\$500\)}{}{}}
\End

\Data\ID{9000085307}
\Updated{2020-07-15 03:07:37}
\Level{B}
\SubArea{Percent problems}
\LANGen{
\Question{
A machine fills \(2\: 000\) 
bottles per hour. As a consequence of technical problems, the efficiency of the machine
decreased by \(10\, \%\).
How many bottles can be filled by the machine with the mentioned technical problems during a
time period of \(8\) 
hours?}
{\(14\: 400\)}{\(13\: 800\)}{\(14\: 500\)}{\(15\: 200\)}{}{}}
\LANGpl{
\Question{
Maszyna napełnia \(2\: 000\) butelek
na godzinę. W wyniku problemów technicznych wydajność maszyny zmniejszyła się o 
\(10\, \%\).
Ile butelek zostanie napełnionych przez tę maszynę w czasie \(8\) 
godzin?}
{\(14\: 400\)}{\(13\: 800\)}{\(14\: 500\)}{\(15\: 200\)}{}{}}
\LANGcs{
\Question{
Automat na plnění lahví naplní obvykle
\(2\: 000\) 
lahví za hodinu. V důsledku technické závady klesl jeho výkon o
\(10\, \%\). Kolik lahví
naplní automat za \(8\) 
hodin při tomto sníženém výkonu?}
{\(14\: 400\)}{\(13\: 800\)}{\(14\: 500\)}{\(15\: 200\)}{}{}}
\LANGsk{
\Question{
Automat na plnenie fliaš naplní obvykle
\(2\: 000\) 
fliaš za hodinu. V dôsledku technickej závady klesol jeho výkon o
\(10\, \%\). Koľko fliaš
naplní automat za \(8\) 
hodín pri tomto zníženom výkone?}
{\(14\: 400\)}{\(13\: 800\)}{\(14\: 500\)}{\(15\: 200\)}{}{}}
\LANGes{
\Question{
Una máquina rellena  \(2\: 000\) 
botellas por hora. Como consecuencia de problemas técnicos, la eficiencia de la máquina disminuyó un \(10\, \%\).
¿Cuántas botellas puede rellenar la máquina con los problemas técnicos mencionados durante un período de tiempo de  \(8\) 
horas?}
{\(14\: 400\)}{\(13\: 800\)}{\(14\: 500\)}{\(15\: 200\)}{}{}}
\End

\Data\ID{9000085308}
\Updated{2020-07-15 03:07:37}
\Level{B}
\SubArea{Percent problems}
\LANGen{
\Question{
The number \(\frac{1}
{2} + \frac{1} 
{3}\) 
is bigger than \(\frac{1}
{2} -\frac{1} 
{3}\).
Express the difference between these numbers in percents of the smaller quantity.}
{\(400\, \%\)}{\(500\, \%\)}{\(300\, \%\)}{\(600\, \%\)}{}{}}
\LANGpl{
\Question{
Liczba  \(\frac{1}
{2} + \frac{1} 
{3}\) 
jest większa niż \(\frac{1}
{2} -\frac{1} 
{3}\).
Wyraź różnicę między tymi liczbami w procentach.}
{\(400\, \%\)}{\(500\, \%\)}{\(300\, \%\)}{\(600\, \%\)}{}{}}
\LANGcs{
\Question{
O kolik procent je \(\frac{1}
{2} + \frac{1} 
{3}\) 
větší než \(\frac{1}
{2} -\frac{1} 
{3}\)?}
{\(400\, \%\)}{\(500\, \%\)}{\(300\, \%\)}{\(600\, \%\)}{}{}}
\LANGsk{
\Question{
O koľko percent je \(\frac{1}
{2} + \frac{1} 
{3}\) 
väčšie než \(\frac{1}
{2} -\frac{1} 
{3}\)?}
{\(400\, \%\)}{\(500\, \%\)}{\(300\, \%\)}{\(600\, \%\)}{}{}}
\LANGes{
\Question{
El número \(\frac{1}
{2} + \frac{1} 
{3}\) 
es mayor que  \(\frac{1}
{2} -\frac{1} 
{3}\).
Expresa la diferencia entre estos números en porcentaje respecto a la cantidad menor.}
{\(400\, \%\)}{\(500\, \%\)}{\(300\, \%\)}{\(600\, \%\)}{}{}}
\End

\Data\ID{9000085309}
\Updated{2020-08-22 04:08:56}
\Level{B}
\SubArea{Percent problems}
\LANGen{
\Question{
The edge of a cube is increased by \(100\, \%\).
Find the increase of the volume of the cube?}
{\(700\, \%\)}{\(400\, \%\)}{\(200\, \%\)}{\(100\, \%\)}{}{}}
\LANGpl{
\Question{
Krawędź sześcianu zostaje zwiększona o \(100\, \%\).
O ile procent wzrośnie objętość tego sześcianu?}
{\(700\, \%\)}{\(400\, \%\)}{\(200\, \%\)}{\(100\, \%\)}{}{}}
\LANGcs{
\Question{
Hranu krychle zvětšíme o \(100\, \%\).
O kolik procent se zvětší objem této krychle?}
{\(700\, \%\)}{\(400\, \%\)}{\(200\, \%\)}{\(100\, \%\)}{}{}}
\LANGsk{
\Question{
Hranu kocky zväčšíme o \(100\, \%\).
O koľko percent sa zväčší objem tejto kocky?}
{\(700\, \%\)}{\(400\, \%\)}{\(200\, \%\)}{\(100\, \%\)}{}{}}
\LANGes{
\Question{
La arista de un cubo aumenta en un  \(100\, \%\).
Calcula el aumento del volumen del cubo?}
{\(700\, \%\)}{\(400\, \%\)}{\(200\, \%\)}{\(100\, \%\)}{}{}}
\End

\Data\ID{9000085310}
\Updated{2020-08-22 04:08:47}
\Level{B}
\SubArea{Percent problems}
\IMG{000853_10-figure0.pdf}
\LANGen{
\Question{
The pad in the form of (not necessarily regular) octagon
is manufactured from a square. The side of the square is
\(4\, \mathrm{cm}\). An isosceles right
triangle of the hypotenuse \(1\, \mathrm{cm}\) 
is removed from each corner of the square which results in the desired octagon. How
much of the material of the initial square goes to the waste?}
{\(12.5\, \%\)}{\(10\, \%\)}{\(15\, \%\)}{\(20\, \%\)}{}{}}
\LANGpl{
\Question{
Podkładka w kształcie ośmiokąta (niekoniecznie regularna) jest produkowana z kwadratu. Bok kwadratu ma długość
\(4\, \mathrm{cm}\). Aby otrzymać pożądany ośmiokąt, z każdego rogu kwadratu usuwany jest trójkąt równoramienny o przeciwprostokątnej długości \(1\, \mathrm{cm}\). Ile materiału wykorzystanego na kwadrat zostaje zmarnowane?}
{\(12.5\, \%\)}{\(10\, \%\)}{\(15\, \%\)}{\(20\, \%\)}{}{}}
\LANGcs{
\Question{
Podložka tvaru osmiúhelníku se lisuje ze čtverce o straně
\(4\, \mathrm{cm}\). Při její
výrobě se ze všech jeho rohů odlomí pravoúhlý trojúhelník s odvěsnou
délky \(1\, \mathrm{cm}\).
Kolik procent plochy původního čtverce tvoří odpad?}
{\(12{,}5\, \%\)}{\(10\, \%\)}{\(15\, \%\)}{\(20\, \%\)}{}{}}
\LANGsk{
\Question{
Podložka tvaru osemuholníka sa lisuje zo štvorce o strane
\(4\, \mathrm{cm}\). Pri jej
výrobe sa zo všetkých jeho rohov odlomí pravouhlý trojuholník s odvesnou
dĺžky \(1\, \mathrm{cm}\).
Koľko percent plochy pôvodného štvorca tvorí odpad?}
{\(12{,}5\, \%\)}{\(10\, \%\)}{\(15\, \%\)}{\(20\, \%\)}{}{}}
\LANGes{
\Question{
Una colchoneta en forma de octágono (no necesariamente regular) se fabrica a partir de un cuadrado. El lado del cuadrado es  \(4\, \mathrm{cm}\). Un triángulo rectángulo isósceles con hipotenusa de  \(1\, \mathrm{cm}\) 
se forma a partir de cada esquina del cuadrado, lo que resulta en el octágono deseado. ¿Qué parte del material del cuadrado inicial va a la basura?}
{\(12.5\, \%\)}{\(10\, \%\)}{\(15\, \%\)}{\(20\, \%\)}{}{}}
\End

\Data\ID{1003083301}
\Updated{2020-08-22 04:08:02}
\Level{C}
\SubArea{Percent problems}
\LANGen{
\Question{
If  a number \( 3b \) is by \( 20\% \) bigger than half of  the number \( 2a + b \), then  the number \( a \) is bigger than \( b \) by}
{\( 100\% \).}{\( 80\% \).}{\( 50\% \).}{\( 200\% \).}{}{}}
\LANGpl{
\Question{
Jeżeli liczba \( 3b \) jest o \( 20\% \) większa od połowy liczby \( 2a + b \), to liczba \( a \) jest większa od \( b \) o}
{\( 100\% \).}{\( 80\% \).}{\( 50\% \).}{\( 200\% \).}{}{}}
\LANGcs{
\Question{
Jestliže je číslo \( 3b \) o \( 20\% \) větší než polovina čísla \( 2a + b \), pak je číslo \( a \) větší než \( b \) o}
{\( 100\% \).}{\( 80\% \).}{\( 50\% \).}{\( 200\% \).}{}{}}
\LANGsk{
\Question{
Ak je číslo \( 3b \) o \( 20\% \) väčšie než polovica čísla \( 2a + b \), potom je číslo \( a \) väčšie než \( b \) o}
{\( 100\% \).}{\( 80\% \).}{\( 50\% \).}{\( 200\% \).}{}{}}
\LANGes{
\Question{
Si un número \( 3b \) es un \( 20\% \) mayor que la mitad del número  \( 2a + b \), entonces el número  \( a \) es mayor que  \( b \) en un}
{\( 100\% \).}{\( 80\% \).}{\( 50\% \).}{\( 200\% \).}{}{}}
\End

\Data\ID{1003083302}
\Updated{2020-07-15 03:07:12}
\Level{C}
\SubArea{Percent problems}
\LANGen{
\Question{
Let \( a \) and \( c \) be positive numbers. A number \( b \) equals \( 96\% \) of the number \( 2a + b \) and \( 64\% \) of the number \( 5a + c \). It means that}
{\( c=70a \).}{\( c=1.5a \).}{\( c=14a \).}{\( c=48a \).}{}{}}
\LANGpl{
\Question{
Liczby \( a \) i \( c \) są dodatnie. Liczba  \( b \) stanowi \( 96\% \) liczby \( 2a + b \) oraz \( 64\% \) liczby \( 5a + c \). Wynika stąd, że}
{\( c=70a \).}{\( c=1{,}5a \).}{\( c=14a \).}{\( c=48a \).}{}{}}
\LANGcs{
\Question{
Jsou dána kladná čísla \( a \) a \( c \). Číslo \( b \) se rovná \( 96\% \) čísla \( 2a + b \) a \( 64\% \) čísla \( 5a + c \). Platí, že}
{\( c=70a \).}{\( c=1{,}5a \).}{\( c=14a \).}{\( c=48a \).}{}{}}
\LANGsk{
\Question{
Sú dané kladné čísla \( a \) a \( c \). Číslo \( b \) sa rovná \( 96\% \) čísla \( 2a + b \) a \( 64\% \) čísla \( 5a + c \). Platí, že}
{\( c=70a \).}{\( c=1{,}5a \).}{\( c=14a \).}{\( c=48a \).}{}{}}
\LANGes{
\Question{
Sean \( a \) y \( c \) números positivos. Un número \( b \) es igual al  \( 96\% \) del número  \( 2a + b \) y el \( 64\% \) del número \( 5a + c \). Esto significa que}
{\( c=70a \).}{\( c=1.5a \).}{\( c=14a \).}{\( c=48a \).}{}{}}
\End

\Data\ID{1003083303}
\Updated{2022-04-23 05:04:26}
\Level{C}
\SubArea{Percent problems}
\LANGen{
\Question{
One half of a number \( a \) is by \( 25\% \) smaller than one third of a number \( b \). Then the number \( b \) is}
{by \( 100\% \) bigger than \( a \).}{by \( 200\% \) bigger than \( a \).}{by \( 50\% \) bigger than \( a \).}{by \( 150\% \) bigger than \( a \).}{}{}}
\LANGpl{
\Question{
Połowa liczby \( a \) jest o \( 25\% \) mniejsza od trzeciej części liczby  \( b \). Wtedy liczba \( b \) jest}
{o \( 100\% \) większa od  \( a \).}{o \( 200\% \) większa od  \( a \).}{o \( 50\% \) większa od  \( a \).}{o \( 150\% \) większa od  \( a \).}{}{}}
\LANGcs{
\Question{
Polovina čísla \( a \) je o \( 25\, \% \) menší než třetina čísla \( b \). Číslo \( b \) je tedy}
{o \( 100\,\% \) větší než \( a \).}{o \( 200\,\% \) větší než \( a \).}{o \( 50\,\% \) větší než \( a \).}{o \( 150\,\% \) větší než \( a \).}{}{}}
\LANGsk{
\Question{
Polovica čísla \( a \) je o \( 25\% \) menšia než tretina čísla \( b \). Číslo \( b \) je teda}
{o \( 100\% \) väčšie než \( a \).}{o \( 200\% \) väčšie než \( a \).}{o \( 50\% \) väčšie než \( a \).}{o \( 150\% \) väčšie než \( a \).}{}{}}
\LANGes{
\Question{
La mitad de un número \( a \) es un \( 25\% \) menor que un tercio de un número  \( b \). Entonces el número  \( b \) es}
{en un \( 100\% \) mayor que  \( a \).}{en un \( 200\% \) mayor que \( a \).}{en un  \( 50\% \) mayor que \( a \).}{en un \( 150\% \) mayor que  \( a \).}{}{}}
\End

\Data\ID{1003083304}
\Updated{2020-08-22 04:08:06}
\Level{C}
\SubArea{Percent problems}
\LANGen{
\Question{
The initial price of  ice bags was \( 24.20 \) zlotys. Now, you can buy \( 22 \) ice bags (which is \( 10\% \) more than before) for the same price. What is the percentage decrease in the price of one ice bag?}
{by app. \( 9\% \)}{by app. \( 8\% \)}{by \( 10\% \)}{by \( 2.2\% \)}{}{}}
\LANGpl{
\Question{
Początkowa cena woreczków do lodu była  \( 24{,}20 \) złote. Teraz możesz kupić \( 22 \) woreczki do lodów (czyli o \( 10\% \) więcej niż poprzednio) za tę samą cenę. O ile procent obniżyła się cena jednego woreczka do robienia kulek lodu?}
{o ok. \( 9\% \)}{o ok. \( 8\% \)}{o \( 10\% \)}{o \( 2{,}2\% \)}{}{}}
\LANGcs{
\Question{
Původní cena sáčků na led byla \( 24{,}20 \) zlotých. Nyní se dá za stejnou cenu koupit \( 22 \) sáčků (což je o \( 10\% \) více než dříve). O kolik procent klesla cena jednoho sáčku?}
{přibližně o \( 9\% \)}{přibližně o \( 8\% \)}{o \( 10\% \)}{o \( 2{,}2\% \)}{}{}}
\LANGsk{
\Question{
Pôvodná cena sáčkov na ľad bola \( 24{,}20 \) zlotých. Teraz sa dá za rovnakú cenu kúpiť \( 22 \) sáčkov (čo je o \( 10\% \) viac než predtým). O koľko percent klesla cena jedného sáčku?}
{približne o \( 9\% \)}{približne o \( 8\% \)}{o \( 10\% \)}{o \( 2{,}2\% \)}{}{}}
\LANGes{
\Question{
El precio inicial de unas bolsas de hielo fue de \( 24.20 \) zlotys. Ahora se pueden comprar \( 22 \) bolsas de hielo (que es un \( 10\% \) más que antes) por el mismo precio. ¿ En qué porcentaje disminuyó el precio de una bolsa de hielo?}
{aproximadamente \( 9\% \)}{aproximadamente \( 8\% \)}{en un \( 10\% \)}{en un  \( 2.2\% \)}{}{}}
\End

\Data\ID{1003083305}
\Updated{2023-06-28 12:06:50}
\Level{C}
\SubArea{Percent problems}
\LANGen{
\Question{
Economists predict that in a certain country \( 10\% \) annual inflation will last for two years. What will be the percentage increase in prices in this country in two years?}
{by \( 21\% \)}{by \( 20\% \)}{by \( 10\% \)}{by \( 1.21\% \)}{}{}}
\LANGpl{
\Question{
Ekonomiści przewidują, że w pewnym kraju inflacja  \( 10\% \) rocznie utrzyma się przez dwa lata. O ile procent wzrosną ceny w tym kraju w ciągu dwóch lat?}
{o \( 21\% \)}{o \( 20\% \)}{o \( 10\% \)}{o \( 1{,}21\% \)}{}{}}
\LANGcs{
\Question{
Ekonomové předpovídají, že v jisté zemi dojde k roční inflaci \( 10\% \), a to po dobu dvou let. O kolik procent po těchto dvou letech vzrostou v zemi ceny?}
{o \( 21\% \)}{o \( 20\% \)}{o \( 10\% \)}{o \( 1{,}21\% \)}{}{}}
\LANGsk{
\Question{
Ekonómovia predpovedali, že v istej krajine dôjde k ročnej inflácii \( 10\% \), a to po dobu dvoch rokov. O koľko percent po týchto dvoch rokoch vzrastú v zemi ceny?}
{o \( 21\% \)}{o \( 20\% \)}{o \( 10\% \)}{o \( 1{,}21\% \)}{}{}}
\LANGes{
\Question{
Los economistas predicen que en cierto país la inflación anual del  \( 10\% \) durará dos años. ¿Cuál será el aumento porcentual de los precios en este país en dos años?}
{de un  \( 21\% \)}{de un  \( 20\% \)}{de un  \( 10\% \)}{de un \( 1.21\% \)}{}{}}
\End

\Data\ID{1003083306}
\Updated{2020-08-22 04:08:38}
\Level{C}
\SubArea{Percent problems}
\LANGen{
\Question{
In a certain school, grades are given according to the rules listed below:
\[\begin{array}{|c|c|} \hline \text{The percentage of the points gained} & \text{grade} \\\hline 0\%\text{ to } 49\% & 1 \\\hline 50\% \text{ to } 59\%  & 2 \\\hline 60\% \text{ to } 79\%  & 3 \\\hline 80\% \text{ to } 89\%  & 4 \\\hline  90\% \text{ to } 94\%  & 5 \\\hline 95\% \text{ to } 100\% & 6 \\\hline \end{array}\]
Krysia gained \( 39 \) points from the history test. There were only \( 2 \) percentage points missing for her to get grade \( 4 \). What was the total amount of points one  could get from the test?}
{\( 50 \) points}{\( 41 \) points}{\( 45 \) points}{\( 55 \) points}{}{}}
\LANGpl{
\Question{
W pewnej szkole oceny ustalane są według zasad określonych w poniższej tabeli:
\[\begin{array}{|c|c|} \hline \text{Procent zdobytych punktów} & \text{Ocena} \\\hline 0\%\text{ do } 49\% & 1 \\\hline 50\% \text{ do } 59\%  & 2 \\\hline 60\% \text{ do } 79\%  & 3 \\\hline 80\% \text{ do } 89\%  & 4 \\\hline  90\% \text{ do } 94\%  & 5 \\\hline 95\% \text{ do } 100\% & 6 \\\hline \end{array}\]
Krysia zdobyła na sprawdzianie z historii  \( 39 \) Tylko dwóch punktów procentowych brakowało jej do czwórki. Ile maksymalnie punktów można było zdobyć na tym sprawdzianie?}
{\( 50 \) punktów}{\( 41 \) punktów}{\( 45 \) punktów}{\( 55 \) punktów}{}{}}
\LANGcs{
\Question{
V jedné škole se známkuje podle těchto pravidel:
\[\begin{array}{|c|c|} \hline \text{Procento získaných bodů} & \text{známka} \\\hline 0\%\text{ až } 49\% & 1 \\\hline 50\% \text{ až } 59\%  & 2 \\\hline 60\% \text{ až } 79\%  & 3 \\\hline 80\% \text{ až } 89\%  & 4 \\\hline  90\% \text{ až } 94\%  & 5 \\\hline 95\% \text{ až } 100\% & 6 \\\hline \end{array}\]
Katka získala v testu z dějepisu \( 39 \) bodů. Chyběly jí jen \( 2 \) procentní body, aby dostala známku \( 4 \). Jaký  maximální počet bodů mohli žáci z testu získat?}
{\( 50 \) bodů}{\( 41 \) bodů}{\( 45 \) bodů}{\( 55 \) bodů}{}{}}
\LANGsk{
\Question{
V jednej škole sa známkuje podľa týchto pravidiel:
\[\begin{array}{|c|c|} \hline \text{Percento získaných bodov} & \text{známka} \\\hline 0\%\text{ až } 49\% & 1 \\\hline 50\% \text{ až } 59\%  & 2 \\\hline 60\% \text{ až } 79\%  & 3 \\\hline 80\% \text{ až } 89\%  & 4 \\\hline  90\% \text{ až } 94\%  & 5 \\\hline 95\% \text{ až } 100\% & 6 \\\hline \end{array}\]
Katka získala v teste z dejepisu \( 39 \) bodov. Chýbali jej len \( 2 \) percentné body, aby dostala známku \( 4 \). Aký  maximálny počet bodov mohli žiaci z testu získať?}
{\( 50 \) bodov}{\( 41 \) bodov}{\( 45 \) bodov}{\( 55 \) bodov}{}{}}
\LANGes{
\Question{
En cierta escuela se dan las notas de acuerdo con las siguientes reglas:
\[\begin{array}{|c|c|} \hline \text{El porcentaje de los puntos obtenidos } & \text{nota} \\\hline 0\%\text{ a } 49\% & 1 \\\hline 50\% \text{ a } 59\%  & 2 \\\hline 60\% \text{ a } 79\%  & 3 \\\hline 80\% \text{ a } 89\%  & 4 \\\hline  90\% \text{ a } 94\%  & 5 \\\hline 95\% \text{ a } 100\% & 6 \\\hline \end{array}\]
Krysia obtuvo  \( 39 \) puntos del examen de historia. Solo le faltaban  \( 2 \) puntos porcentuales para obtener la nota  \( 4 \). ¿Cuál era la cantidad total de puntos que uno podría obtener en el examen?}
{\( 50 \) puntos}{\( 41 \) puntos}{\( 45 \) puntos}{\( 55 \) puntos}{}{}}
\End

\Data\ID{1003083307}
\Updated{2023-06-28 12:06:52}
\Level{C}
\SubArea{Percent problems}
\LANGen{
\Question{
Amount of \( 3000 \) zlotys was paid to bank \( A \) for interest-bearing investment of \( p\% \) per year. Amount \( 4000 \) zlotys was paid to bank \( B \), where the interest rate is by \( 2 \) percentage points higher than in bank \( A \). After one year, the interest earned in both banks amounted to \( 360 \) zlotys. What was the interest rate \( p \)?}
{\( 4\% \)}{\( 4.5\% \)}{\( 5\% \)}{\( 2\% \)}{}{}}
\LANGpl{
\Question{
Kwotę  \( 3000 \) zł wpłacono do banku \( A \) na lokatę oprocentowaną   \( p\% \) w skali roku. Kwotę \( 4000 \) zł wpłacono do banku  \( B \), w którym oprocentowanie jest o  \( 2 \) punkty procentowe wyższe niż w banku \( A \). Po roku odsetki uzyskane w obu bankach wyniosły \( 360 \) zł. Jakie było oprocentowanie   \( p \)?}
{\( 4\% \)}{\( 4{,}5\% \)}{\( 5\% \)}{\( 2\% \)}{}{}}
\LANGcs{
\Question{
Banka \( A \) obdržela \( 3000 \) zlotých za investice s roční úrokovou sazbou \( p\% \). Banka \( B \) získala \( 4000 \) zlotých za investice s úrokovou sazbou, která byla o \( 2 \) procentní body lepší než u banky \( A \). Po jednom roce byl celkový zisk obou bank \( 360 \) zlotých. Jaká byla sazba \( p \)?}
{\( 4\% \)}{\( 4{,}5\% \)}{\( 5\% \)}{\( 2\% \)}{}{}}
\LANGsk{
\Question{
Banka \( A \) obdržala \( 3000 \) zlotých za investície s ročnou sadzbou \( p\% \). Banka \( B \) získala \( 4000 \) zlotých za úrokovú sadzbu, ktorá bola o \( 2 \) percentné body lepšia než banka \( A \). Po jednom roku bol celkový zisk oboch bánk \( 360 \) zlotých. Aká bola sadzba \( p \)?}
{\( 4\% \)}{\( 4{,}5\% \)}{\( 5\% \)}{\( 2\% \)}{}{}}
\LANGes{
\Question{
Se pagó una cantidad de  \( 3000 \) zlotys al banco  \( A \) por una inversión con una tasa de interés anual del  \( p\% \). Se pagó  \( 4000 \) zlotys al banco  \( B \), donde la tasa de interés era  \( 2 \) puntos porcentuales más alta que en el banco \( A \). Después de un año, los intereses ganados en ambos bancos ascendieron a \( 360 \) zlotys. ¿Cuál era la tasa de interés \( p \)?}
{\( 4\% \)}{\( 4.5\% \)}{\( 5\% \)}{\( 2\% \)}{}{}}
\End

\Data\ID{1003083308}
\Updated{2023-06-28 12:06:27}
\Level{C}
\SubArea{Percent problems}
\LANGen{
\Question{
Company \( A \)  had  a turnover by \( 50\% \) higher in the second quarter  than in the first quarter. In the third quarter its turnover decreased by \( 20\% \) and in the fourth quarter the turnover decreased by another \( 20\% \). What was the turnover of company \( A \) in the second quarter, if in the fourth quarter it amounted to \( 1440000 \) zlotys?}
{\( 2250000\) zlotys}{\( 1500000 \) zlotys}{\( 1800000\) zlotys}{\( 2000000 \) zlotys}{}{}}
\LANGpl{
\Question{
Firma \( A \)  w II kwartale miała obroty o  \( 50\% \) większe niż w I kwartale. W III kwartale jej obroty spadły o  \( 20\% \), zaś w IV spadły o kolejne  \( 20\% \). Jakie były obroty firmy  \( A \) w II kwartale, jeśli w IV kwartale wyniosły  \( 1440000 \) zł?}
{\( 2250000\) zł}{\( 1500000 \) zł}{\( 1800000\) zł}{\( 2000000 \) zł}{}{}}
\LANGcs{
\Question{
Společnost \( A \)  měla ve druhém kvartále obrat o \( 50\% \) vyšší než v prvním. Ve třetím kvartále její obrat klesl o \( 20\% \) a ve čtvrtém klesl o dalších \( 20\% \). Jaký byl obrat společnosti \( A \) ve druhém kvartále, pokud víme, že ve čtvrtém kvartále byl obrat \( 1440000 \) zlotých?}
{\( 2250000\) zlotých}{\( 1500000 \) zlotých}{\( 1800000\) zlotých}{\( 2000000 \) zlotých}{}{}}
\LANGsk{
\Question{
Spoločnosť \( A \)  mala v druhom kvartáli obrat o \( 50\% \) vyšší než v prvom. V treťom kvartále jej obrat klesol o \( 20\% \) a vo štvrtom klesol o ďalších \( 20\% \). Aký bol obrat spoločnosti \( A \) v druhom kvartále, ak vieme, že vo štvrtom kvartále bol obrat vyčíslený na \( 1440000 \) zlotých?}
{\( 2250000\) zlotých}{\( 1500000 \) zlotých}{\( 1800000\) zlotých}{\( 2000000 \) zlotých}{}{}}
\LANGes{
\Question{
La empresa \( A \)  tuvo una facturación un  \( 50\% \) mayor en el segundo trimestre que en el primer trimestre. En el tercer trimestre, su facturación disminuyó  un  \( 20\% \) y en el cuarto trimestre la facturación disminuyó  otro  \( 20\% \). ¿Cuál fue la facturación de la empresa  \( A \) en el segundo trimestre, si en el cuarto trimestre ascendió a \( 1440000 \) zlotys?}
{\( 2250000\) zlotys}{\( 1500000 \) zlotys}{\( 1800000\) zlotys}{\( 2000000 \) zlotys}{}{}}
\End

\Data\ID{1003090501}
\Updated{2022-04-23 05:04:26}
\Level{C}
\SubArea{Percent problems}
\LANGen{
\Question{
Amongst \( 400 \) people surveyed, \( 65\% \) of them speak fluent English, \( 47\% \) speak French and \( 24\% \) speak both languages. How many of them do not know any of these languages?}
{\( 48 \) people}{\( 50 \) people}{\( 88 \) people}{\( 58 \) people}{}{}}
\LANGpl{
\Question{
Wśród \( 400 \) badanych osób, \( 65\% \) biegle zna język angielski, \( 47\% \) zna biegle język francuski, a \( 24\% \) zna oba te języki. Ile spośród badanych nie zna żadnego z tych języków?}
{\( 48 \) osób}{\( 50 \) osób}{\( 88 \) osób}{\( 58 \) osób}{}{}}
\LANGcs{
\Question{
Ze \( 400 \) dotázaných lidí mluvilo \( 65\, \% \) plynnou angličtinou, \( 47\, \% \) mluvilo francouzsky, \( 24\, \% \) mluvilo oběma jazyky. Kolik lidí nemluvilo žádným z uvedených jazyků?}
{\( 48 \) lidí}{\( 50 \) lidí}{\( 88 \) lidí}{\( 58 \) lidí}{}{}}
\LANGsk{
\Question{
Z \( 400 \) opýtaných ľudí hovorilo \( 65\% \) plynulou angličtinou, \( 47\% \) hovorilo francúzsky, \( 24\% \) hovorilo oboma jazykmi. Koľko ľudí nehovorilo žiadnym z uvedených jazykov?}
{\( 48 \) ľudí}{\( 50 \) ľudí}{\( 88 \) ľudí}{\( 58 \) ľudí}{}{}}
\LANGes{
\Question{
Entre las \( 400 \) personas encuestadas, el \( 65\% \) de ellas habla inglés con fluidez, el \( 47\% \) habla francés y el  \( 24\% \) habla ambos idiomas. ¿Cuántos de ellos no conocen ninguno de estos idiomas?}
{\( 48 \) personas}{\( 50 \) personas}{\( 88 \) personas}{\( 58 \) personas}{}{}}
\End

\Data\ID{1003090502}
\Updated{2022-04-23 05:04:26}
\Level{C}
\SubArea{Percent problems}
\LANGen{
\Question{
Find the number \( x \) if \( 60\% \) of its double is \( 30\% \) of the number \( x +5\) .}
{\( 1\frac23 \)}{\( 1.67 \)}{\( 1.6 \)}{\( 1\frac13 \)}{}{}}
\LANGpl{
\Question{
Jaka to liczba \( x \), której  \( 60\% \) dwukrotności liczby   \( x \)  jest równe \( 30\% \) liczby   \( x \) powiększonej o  \( 5 \).}
{\( 1\frac23 \)}{\( 1{,}67 \)}{\( 1{,}6 \)}{\( 1\frac13 \)}{}{}}
\LANGcs{
\Question{
Najdi číslo \( x \), pro které platí, že \( 60\% \) jeho dvojnásobku se rovná \( 30\% \) čísla \( x \) zvětšeného o \( 5 \).}
{\( 1\frac23 \)}{\( 1{,}67 \)}{\( 1{,}6 \)}{\( 1\frac13 \)}{}{}}
\LANGsk{
\Question{
Nájdi číslo \( x \), pre ktoré platí, že \( 60\% \) jeho dvojnásobku sa rovná \( 30\% \) čísla \( x \) zväčšeného o \( 5 \).}
{\( 1\frac23 \)}{\( 1{,}67 \)}{\( 1{,}6 \)}{\( 1\frac13 \)}{}{}}
\LANGes{
\Question{
Encuentra el número  \( x \) si el  \( 60\% \) de su doble es el  \( 30\% \) del número \( x \) más \( 5 \).}
{\( 1\frac23 \)}{\( 1.67 \)}{\( 1.6 \)}{\( 1\frac13 \)}{}{}}
\End

\Data\ID{1003090503}
\Updated{2020-08-22 04:08:51}
\Level{C}
\SubArea{Percent problems}
\LANGen{
\Question{
The total weight of two alloys is \( 800\,\mathrm{kg} \). In the alloy containing \( 4.5\% \) of  lead (Pb) there is by \( 19\,\mathrm{kg} \) of lead more than in the other, containing \( 4\% \) of lead. How many kilograms of lead is in each of these alloys?}
{\( 8 \) kilograms (\( 4\% \) Pb) and \( 27 \) kilograms (\( 4.5\% \) Pb).}{\( 18 \) kilograms (\( 4\% \) Pb) and \( 37 \) kilograms (\( 4.5\% \) Pb).}{\( 36 \) kilograms (\( 4\% \) Pb) and \( 55 \) kilograms (\( 4.5\% \) Pb).}{\( 10 \) kilograms (\( 4\% \) Pb) and \( 29 \) kilograms (\( 4.5\% \) Pb).}{}{}}
\LANGpl{
\Question{
Łączna masa dwóch stopów jest równa  \( 800\,\mathrm{kg} \). W stopie zawierającym  \( 4{,}5\% \) ołowiu (Pb) jest o  \( 19\,\mathrm{kg} \) więcej niż w drugim, zawierającym  \( 4\% \) ołowiu. Ile kilogramów ołowiu jest w każdym z tych stopów?}
{\( 8 \) kg(\( 4\% \) Pb) i \( 27 \) kg(\( 4{,}5\% \) Pb).}{\( 18 \) kg(\( 4\% \) Pb) i \( 37 \) kg(\( 4{,}5\% \) Pb).}{\( 36 \) kg(\( 4\% \) Pb) i \( 55 \) kg(\( 4{,}5\% \) Pb).}{\( 10 \) kg(\( 4\% \) Pb) i \( 29 \) kg(\( 4{,}5\% \) Pb).}{}{}}
\LANGcs{
\Question{
Celková váha dvou kusů slitiny je \( 800\,\mathrm{kg} \). Ve slitině, která obsahuje \( 4{,}5\% \) olova (Pb), je o \( 19\,\mathrm{kg} \) více olova než v druhé, která obsahuje \( 4\% \) olova. Kolik kilogramů olova je v každém kusu?}
{\( 8 \) kilogramů (\( 4\% \) Pb) a \( 27 \) kilogramů (\( 4{,}5\% \) Pb).}{\( 18 \) kilogramů (\( 4\% \) Pb) a \( 37 \) kilogramů (\( 4{,}5\% \) Pb).}{\( 36 \) kilogramů (\( 4\% \) Pb) a \( 55 \) kilogramů (\( 4{,}5\% \) Pb).}{\( 10 \) kilogramů (\( 4\% \) Pb) a \( 29 \) kilogramů (\( 4{,}5\% \) Pb).}{}{}}
\LANGsk{
\Question{
Celková váha dvoch kusov zliatiny je \( 800\,\mathrm{kg} \). V zliatine, ktorá obsahuje \( 4{,}5\% \) olova (Pb) je o \( 19\,\mathrm{kg} \) viac olova než v druhej, ktorá obsahuje \( 4\% \) olova. Koľko kilogramov olova je v každom kuse?}
{\( 8 \) kilogramov (\( 4\% \) Pb) a \( 27 \) kilogramov (\( 4{,}5\% \) Pb).}{\( 18 \) kilogramov (\( 4\% \) Pb) a \( 37 \) kilogramov (\( 4{,}5\% \) Pb).}{\( 36 \) kilogramov (\( 4\% \) Pb) a \( 55 \) kilogramov (\( 4{,}5\% \) Pb).}{\( 10 \) kilogramov (\( 4\% \) Pb) a \( 29 \) kilogramov (\( 4{,}5\% \) Pb).}{}{}}
\LANGes{
\Question{
El peso total de dos aleaciones es de  \( 800\,\mathrm{kg} \). En la aleación que contiene  \( 4.5\% \) de plomo  (Pb) hay  \( 19\,\mathrm{kg} \) de plomo más que en la otra, que contiene  \( 4\% \) de plomo. ¿Cuántos kilogramos de plomo hay en cada una de estas aleaciones?}
{\( 8 \) kilogramos (\( 4\% \) Pb) y \( 27 \) kilogramos (\( 4.5\% \) Pb).}{\( 18 \) kilogramos (\( 4\% \) Pb) y \( 37 \) kilogramos (\( 4.5\% \) Pb).}{\( 36 \) kilogramos (\( 4\% \) Pb) y \( 55 \) kilogramos (\( 4.5\% \) Pb).}{\( 10 \) kilogramos (\( 4\% \) Pb) y \( 29 \) kilogramos (\( 4.5\% \) Pb).}{}{}}
\End

\Data\ID{1003090504}
\Updated{2020-11-30 04:11:18}
\Level{C}
\SubArea{Percent problems}
\LANGen{
\Question{
Interest on two construction loans of the total value of \( 100000 \) zlotys amounts to \( 3.150 \) zlotys a year with the interest rate on one of the loans being \( 3\% \) and the other \( 3.5\% \). Calculate the amount of each of these loans.}
{\( 70000 \) zlotys of \( 3\% \)  interest rate loan  and \( 30000 \) zlotys  of \( 3.5\% \) interest rate loan}{\( 30000 \) zlotys of \( 3\% \)  interest rate loan  and  \( 70000 \) zlotys  of \( 3.5\% \) interest rate loan}{\( 50000 \) zlotys of \( 3\% \)  interest rate loan  and  \( 50000 \) zlotys  of \( 3.5\% \) interest rate loan}{\( 80000 \) zlotys of \( 3\% \)  interest rate loan  and  \( 20000 \) zlotys  of \( 3.5\% \) interest rate loan}{}{}}
\LANGpl{
\Question{
Odsetki od dwóch kredytów budowlanych o łącznej wartości  \( 100000 \) zł wynoszą rocznie  \( 3{,}150 \) zł, przy czym stopa procentowa w skali jednego z kredytów jest  \( 3\% \), a drugiego  \( 3{,}5\% \). Oblicz kwotę każdego z tych kredytów.}
{\( 70000 \) zł kredytu oprocentowanego \( 3\% \)  i  \( 30000 \) zł kredytu oprocentowanego  \( 3{,}5\% \)}{\( 30000 \) zł kredytu oprocentowanego  \( 3\% \)  i  \( 70000 \) zł kredytu oprocentowanego  \( 3{,}5\% \)}{\( 50000 \) zł kredytu oprocentowanego  \( 3\% \)  i  \( 50000 \) zł kredytu oprocentowanego  \( 3{,}5\% \)}{\( 80000 \) zł kredytu oprocentowanego  \( 3\% \)  i  \( 20000 \) zł kredytu oprocentowanego  \( 3{,}5\% \)}{}{}}
\LANGcs{
\Question{
Roční výše zaplacených úroků dvou hypotečních úvěrů v celkové výši \( 100000 \) zlotých je \( 3{,}150 \) zlotých, přičemž úroková sazba jednoho úvěru je \( 3\% \) a druhého \( 3{,}5\% \). Vypočítejte výši jednotlivých úvěrů.}
{\( 70000 \) zlotých na \( 3\% \) úrok a \( 30000 \) zlotých na \( 3{,}5\% \) úrok}{\( 30000 \) zlotých na \( 3\% \) úrok a \( 70000 \) zlotých na \( 3{,}5\% \) úrok}{\( 50000 \) zlotých na \( 3\% \) úrok a \( 50000 \) zlotých na \( 3{,}5\% \) úrok}{\( 80000 \) zlotých na \( 3\% \) úrok a \( 20000 \) zlotých na \( 3{,}5\% \) úrok}{}{}}
\LANGsk{
\Question{
Ročná výška zaplatených úrokov v dvoch hypotekárnych úveroch v celkovej výške \( 100000 \) zlotých je \( 3{,}150 \) zlotých, pričom úroková sadzba jedného úveru je \( 3\% \) a druhého \( 3{,}5\% \). Vypočítajte výšku jednotlivých úverov.}
{\( 70000 \) zlotých na \( 3\% \) úrok a \( 30000 \) zlotých na \( 3{,}5\% \) úrok}{\( 30000 \) zlotých na \( 3\% \) úrok a \( 70000 \) zlotých na \( 3{,}5\% \) úrok}{\( 50000 \) zlotých na \( 3\% \) úrok a \( 50000 \) zlotých na \( 3{,}5\% \) úrok}{\( 80000 \) zlotých na \( 3\% \) úrok a \( 20000 \) zlotých na \( 3{,}5\% \) úrok}{}{}}
\LANGes{
\Question{
El interés en dos créditos hipotecarios por un valor total de  \( 100000 \) zlotys asciende a  \( 3.150 \) zlotys al año, con una tasa de interés en uno de los préstamos del \( 3\% \) y en el otro \( 3.5\% \). Calcula la cantidad de cada uno de estos créditos.}
{\( 70000 \) zlotys de crédito con una tasa de interés del \( 3\% \) y \( 30000 \) zlotys  de crédito con una tasa de interés del \( 3.5\% \)}{\( 30000 \) zlotys de crédito con una tasa de interés del \( 3\% \) y  \( 70000 \) zlotys de crédito con una tasa de interés del \( 3.5\% \)}{\( 50000 \) zlotys de crédito con una tasa de interés del \( 3\% \)  y  \( 50000 \) zlotys de crédito con una tasa de interés del \( 3.5\% \)}{\( 80000 \) zlotys de crédito con una tasa de interés del  \( 3\% \)  y \( 20000 \) zlotys de crédito con una tasa de interés del \( 3.5\% \)}{}{}}
\End

\Data\ID{1003090505}
\Updated{2020-07-15 03:07:12}
\Level{C}
\SubArea{Percent problems}
\LANGen{
\Question{
The interest rate on deposits  in a certain bank initially equaled \( 5\% \) per year but then it was increased  by \( 1.2 \) percentage point. What was the percentage increase of the interest rate?}
{\( 24\% \)}{\( 6\% \)}{\( 6.2\% \)}{\( 20\% \)}{}{}}
\LANGpl{
\Question{
Oprocentowanie lokat w pewnym banku, równe początkowo  \( 5\% \) w skali roku, wzrosło o  \( 1{,}2 \) punktu procentowego. O ile procent wzrosło to oprocentowanie?}
{\( 24\% \)}{\( 6\% \)}{\( 6{,}2\% \)}{\( 20\% \)}{}{}}
\LANGcs{
\Question{
Úroková sazba na vklad v nejmenované bance byla původně stanovena na \( 5\% \) ročně, ovšem potom se zvýšila o \( 1{,}2 \) procentního bodu. O kolik procent se úroková sazba zvýšila?}
{\( 24\% \)}{\( 6\% \)}{\( 6{,}2\% \)}{\( 20\% \)}{}{}}
\LANGsk{
\Question{
Úroková sadzba na vklad v nemenovanej banke bola pôvodne stanovená na \( 5\% \) ročne, avšak potom sa zvýšila o \( 1{,}2 \) percentného bodu. O koľko percent sa úroková sadzba zvýšila?}
{\( 24\% \)}{\( 6\% \)}{\( 6{,}2\% \)}{\( 20\% \)}{}{}}
\LANGes{
\Question{
La tasa de interés de los depósitos en un determinado banco fue inicialmente del  \( 5\% \) por año, pero luego se incrementó en \( 1.2 \) puntos porcentuales. ¿Cuál fue el aumento porcentual de la tasa de interés?}
{\( 24\% \)}{\( 6\% \)}{\( 6.2\% \)}{\( 20\% \)}{}{}}
\End

\Data\ID{1003090506}
\Updated{2020-07-15 03:07:12}
\Level{C}
\SubArea{Percent problems}
\LANGen{
\Question{
The interest rate on deposits in  a certain bank initially equaled \( 5\% \) per year but then it was decreased by \( 20\% \). By how many percentage points has the interest rate decreased?}
{by \( 1\) percentage point}{by \( 1.2 \) percentage point}{by \( 2 \) percentage points}{by \( 2.1 \) percentage point}{}{}}
\LANGpl{
\Question{
Oprocentowanie lokat w pewnym banku, równe początkowo  \( 5\% \) w skali roku, zmniejszyło się o  \( 20\% \). O ile punktów procentowych zmniejszyło się to oprocentowanie?}
{o \( 1\) punkt procentowy}{o \( 1{,}2 \) punktu procentowego}{o \( 2 \) punkty procentowe}{o \( 2{,}1 \) punktu procentowego}{}{}}
\LANGcs{
\Question{
Úroková sazba v nejmenované bance byla stanovena na \( 5\% \) ročně, ovšem potom klesla o \( 20\% \). O kolik procentních bodů se úroková sazba snížila?}
{o \( 1\) procentní bod}{o \( 1{,}2 \) procentního bodu}{o \( 2 \) procentní body}{o \( 2{,}1 \) procentního bodu}{}{}}
\LANGsk{
\Question{
Úroková sadzba v nemenovanej banke bola stanovená na \( 5\% \) ročne, avšak potom klesla o \( 20\% \). O koľko percentných bodov sa úroková sadzba znížila?}
{o \( 1\) percentný bod}{o \( 1{,}2 \) percentného bodu}{o \( 2 \) percentné body}{o \( 2{,}1 \) percentného bodu}{}{}}
\LANGes{
\Question{
La tasa de interés de los depósitos en un determinado banco fue inicialmente del \( 5\% \) por año, pero luego se redujo en un  \( 20\% \). ¿En cuántos puntos porcentuales disminuyó la tasa de interés?}
{en \( 1\) puntos porcentuales}{en \( 1.2 \) puntos porcentuales}{en \( 2 \) puntos porcentuales}{en \( 2.1 \) puntos porcentuales}{}{}}
\End

\Data\ID{1003090507}
\Updated{2020-07-15 03:07:12}
\Level{C}
\SubArea{Percent problems}
\LANGen{
\Question{
A certain amount of money was put into a bank account for an annual deposit of \( 3.5\% \) per year. After a year the bank paid a tax in the amount of \( 14 \) zlotys from the gained interests. What amount was placed on the deposit if the interest tax amounted to \( 20\% \)?}
{\( 2000 \) zlotys}{\( 2500 \) zlotys}{\( 3000 \) zlotys}{\( 1400 \) zlotys}{}{}}
\LANGpl{
\Question{
Pewną kwotę wpłacono do banku na lokatę roczną oprocentowaną  \( 3{,}5\% \) skali roku. Od dopisanych po roku odsetek bank odprowadził podatek w wysokości  \( 14 \) zł. Jaką kwotę umieszczono na lokacie, jeżeli podatek od odsetek wynosił \( 20\% \)?}
{\( 2000 \) zł}{\( 2500 \) zł}{\( 3000 \) zł}{\( 1400 \) zł}{}{}}
\LANGcs{
\Question{
Na účet s ročním úročením \( 3{,}5\% \) byla uložena jistá suma peněz. Po roce banka zaplatila ze získaných úroků daň ve výši \( 14 \) zlotých. Jaká suma byla uložena na účet, když víme, že úroky jsou zdaněny \( 20\% \)?}
{\( 2000 \) zlotých}{\( 2500 \) zlotých}{\( 3000 \) zlotých}{\( 1400 \) zlotých}{}{}}
\LANGsk{
\Question{
Na účet s ročným úročením \( 3{,}5\% \) bola uložená istá suma peňazí. Po roku banka zaplatila zo získaných úrokov daň vo výške \( 14 \) zlotých. Aká suma bola uložená na účet, keď vieme, že úroky sú zdanené \( 20\% \)?}
{\( 2000 \) zlotých}{\( 2500 \) zlotých}{\( 3000 \) zlotých}{\( 1400 \) zlotých}{}{}}
\LANGes{
\Question{
Una cierta cantidad de dinero fue depositada en una cuenta bancaria para un depósito anual de  \( 3.5\% \) por año. Después de un año, el banco pagó un impuesto por la cantidad de  \( 14 \) zlotys de los intereses ganados. ¿Qué cantidad fue colocada en el depósito si el impuesto de interés ascendió al \( 20\% \)?}
{\( 2000 \) zlotys}{\( 2500 \) zlotys}{\( 3000 \) zlotys}{\( 1400 \) zlotys}{}{}}
\End

\Data\ID{1003090508}
\Updated{2020-07-15 03:07:12}
\Level{C}
\SubArea{Percent problems}
\LANGen{
\Question{
A  factory is supplied with raw materials by rail. Due to the renovation of the track, the time of passing the goods from the warehouse to the factory increased by \( 25\% \). How much did the average speed of carriage  on the delivery route decrease?}
{by \( 20\% \)}{by \( 25\% \)}{by \( 4\% \)}{by \( 75\% \)}{}{}}
\LANGpl{
\Question{
Fabryka zaopatrywana jest w surowce drogą kolejową. Z powodu remontu torów czas przejazdu składu towarowego do fabryki wydłużył się o \( 25\% \). O ile procent zmniejszyła się średnia prędkość składu na trasie dostaw?}
{o \( 20\% \)}{o \( 25\% \)}{o \( 4\% \)}{o \( 75\% \)}{}{}}
\LANGcs{
\Question{
Do továrny se vozí suroviny po železnici. Kvůli práci na trati se doba dopravy zboží ze skladu do továrny prodloužila o \( 25\% \). O kolik se snížila rychlost zásobovacího vlaku?}
{o \( 20\% \)}{o \( 25\% \)}{o \( 4\% \)}{o \( 75\% \)}{}{}}
\LANGsk{
\Question{
Do továrne sa vozia suroviny po železnici. Kvôli práci na trati sa doba dopravy tovaru zo skladu do továrne predĺžila o \( 25\% \). O koľko sa znížila rýchlosť zásobovacieho vlaku?}
{o \( 20\% \)}{o \( 25\% \)}{o \( 4\% \)}{o \( 75\% \)}{}{}}
\LANGes{
\Question{
Una fábrica se abastece de materias primas por ferrocarril. Debido a la renovación de la vía, el tiempo de paso de los productos desde el almacén a la fábrica aumentó en un \( 25\% \). ¿Cuánto disminuyó la velocidad media del tren de abastecimiento?}
{en un \( 20\% \)}{en un \( 25\% \)}{en un \( 4\% \)}{en un \( 75\% \)}{}{}}
\End

\Data\ID{2000016301}
\Updated{2022-07-21 05:07:26}
\Level{C}
\SubArea{Percent problems}
\LANGen{
\Question{
Investors Thomas and Paul invested the same amount. After the first year, the value of Thomas's investments decreased by \(5{\small{{}^\text{o}\mkern-5mu/\mkern-3mu_\text{o}}}\), but after the following year, their value increased by \(5{\small{{}^\text{o}\mkern-5mu/\mkern-3mu_\text{o}}}\). Paul's investments were more stable. After the first year, the value of his investments increased by \(2{\small{{}^\text{o}\mkern-5mu/\mkern-3mu_\text{oo}}}\), but after the second year, it decreased again by \(2{\small{{}^\text{o}\mkern-5mu/\mkern-3mu_\text{oo}}}\). Determine the true statement about the value of Thomas and Paul's investments two years after the investment.}
{Paul's investments will have a higher value.}{Thomas's investments will have a higher value.}{The values of both investments will again be the same.}{It is not possible to determine the ratio of the values of Paul and Thomas's investments from the given data.}{}{}}
\LANGpl{
\Question{
Inwestorzy Thomas i Paul zainwestowali tę samą kwotę. Po pierwszym roku wartość inwestycji Thomasa spadła o \(5{\small{{}^\text{o}\mkern-5mu/\mkern-3mu_\text{o}}}\), ale po kolejnym roku ich wartość wzrosła o \(5{\small{{}^\text{o}\mkern-5mu/\mkern-3mu_\text{o}}}\). Inwestycje Paula były bardziej stabilne. Po pierwszym roku wartość jego inwestycji wzrosła o \(2{\small{{}^\text{o}\mkern-5mu/\mkern-3mu_\text{oo}}}\),  ale po drugim roku ponownie spadła o \(2{\small{{}^\text{o}\mkern-5mu/\mkern-3mu_\text{oo}}}\). Określ prawdziwe stwierdzenie dotyczące wartości inwestycji Tomasza i Pawła dwa lata po inwestycji.}
{Inwestycje Paula będą miały większą wartość.}{Inwestycje Thomasa będą miały większą wartość.}{Wartości obu inwestycji znów będą takie same.}{Z podanych danych nie jest możliwe określenie stosunku wartości inwestycji Paula i Thomasa.}{}{}}
\LANGcs{
\Question{
Investoři Tomáš a Pavel investovali stejnou částku. Tomášovi po prvním roce klesla hodnota jeho investic o \(5\,{\small{{}^\text{o}\mkern-5mu/\mkern-3mu_\text{o}}}\), ale po dalším roce se jejich hodnota o \(5\,{\small{{}^\text{o}\mkern-5mu/\mkern-3mu_\text{o}}}\) zvýšila. Pavlovy investice  byly stabilnější. Po prvním roce se zvýšila hodnota jeho investic  o \(2\,{\small{{}^\text{o}\mkern-5mu/\mkern-3mu_\text{oo}}}\), ale po druhém roce se zase o \(2\,{\small{{}^\text{o}\mkern-5mu/\mkern-3mu_\text{oo}}}\) snížila. Určete pravdivé tvrzení ohledně hodnoty investic Tomáše a Pavla po dvou letech od investice.}
{Vyšší hodnotu budou mít investice Pavla.}{Vyšší hodnotu budou mít investice Tomáše.}{Hodnoty obou investic budou opět stejné.}{Z daných údajů není možno poměr hodnot investic Pavla a Tomáše určit.}{}{}}
\LANGsk{
\Question{
Investori Tomáš a Peter investovali rovnakú sumu. Tomášovi po prvom roku klesla hodnota jeho investícií o \(5{\small{{}^\text{o}\mkern-5mu/\mkern-3mu_\text{o}}}\), ale po ďalšom roku sa jeho hodnota o \(5{\small{{}^\text{o}\mkern-5mu/\mkern-3mu_\text{o}}}\) zvýšila. Petrove investície  boli stabilnejšie. Po prvom roku sa zvýšila hodnota jeho investícií  o \(2{\small{{}^\text{o}\mkern-5mu/\mkern-3mu_\text{oo}}}\), ale po druhom roku sa zase o \(2{\small{{}^\text{o}\mkern-5mu/\mkern-3mu_\text{oo}}}\) znížila. Určte pravdivé tvrdenie o hodnotách investícií Tomáša a Petra po dvoch rokoch od investície.}
{Vyššiu hodnotu budú mať  Petrove investície.}{Vyššiu hodnotu budú mať Tomášove investície.}{Investície obidvoch budú mať rovnakú hodnotu.}{Z daných údajov nie je možné pomer hodnôt investícií Petra a Tomáša určiť.}{}{}}
\LANGes{
\Question{
Dos inversores, Tomás y Paulo, invirtieron la misma cantidad. Después del primer año, el valor de las inversiones de Tomás disminuyó en un \(5{\small{{}^\text{o}\mkern-5mu/\mkern-3mu_\text{o}}}\), pero después del siguiente año, su valor aumentó en un \(5{\small{{}^\text{o}\mkern-5mu/\mkern-3mu_\text{o}}}\). Las inversiones de Paulo eran más estables. Después del primer año, el valor de sus inversiones aumentó en un \(2{\small{{}^\text{o}\mkern-5mu/\mkern-3mu_\text{oo}}}\), pero después del segundo año, volvió a disminuir en un \(2{\small{{}^\text{o}\mkern-5mu/\mkern-3mu_\text{oo}}}\). Determina la proposición verdadera sobre el valor de las inversiones de Tomás y Paulo dos años después de la inversión.}
{Las inversiones de Paulo tendrán un valor más alto.}{Las inversiones de Tomás tendrán un valor más alto.}{Los valores de ambas inversiones volverán a ser los mismos.}{No es posible determinar la proporción de los valores de las inversiones de Paulo y Tomás a partir de los datos proporcionados.}{}{}}
\End

\Data\ID{2010006401}
\Updated{2022-05-16 11:05:29}
\Level{C}
\SubArea{Percent problems}
\LANGen{
\Question{
Of the school's \(300\) pupils, \(67\%\) attend the sports club and \(42\%\) attend the music club. \(18\%\) attend both clubs. How many children do not attend any club?}
{\(27\) children}{\(81\) children}{\(87\) children}{\(30\) children}{}{}}
\LANGpl{
\Question{
Spośród \(300\) uczniów szkoły,  \(67\%\) uczęszcza do klubu sportowego a \(42\%\) do klubu muzycznego. \(18\%\) uczęszcza do obu klubów. Ile dzieci nie uczęszcza do żadnego klubu?}
{\(27\) dzieci}{\(81\) dzieci}{\(87\) dzieci}{\(30\) dzieci}{}{}}
\LANGcs{
\Question{
Ze \(300\) žáků školy navštěvuje sportovní kroužek \(67\,\%\) a hudební kroužek \(42\,\%\) dětí. Oba kroužky navštěvuje \(18\,\%\). Kolik dětí nechodí ani do jednoho kroužku?}
{\(27\) dětí}{\(81\) dětí}{\(87\) dětí}{\(30\) dětí}{}{}}
\LANGsk{
\Question{
V škole je \(300\) žiakov. Športový krúžok navštevuje \(67\%\) žiakov a hudobný krúžok navštevuje \(42\%\) žiakov. Obidva krúžky navštevuje \(18\%\) žiakov. Koľko žiakov školy nenavštevuje ani jeden z krúžkov?}
{\(27\) žiakov}{\(81\) žiakov}{\(87\) žiakov}{\(30\) žiakov}{}{}}
\LANGes{
\Question{
De los \(300\) alumnos de la escuela \(67 \%\) asiste al club deportivo, \(42 \%\) asiste al club de música y \(18 \%\) asiste a ambos clubes. ¿Cuántos niños no asisten a ningún club?}
{\(27\) niños}{\(81\) niños}{\(87\) niños}{\(30\) niños}{}{}}
\End

\Data\ID{2010006402}
\Updated{2022-05-16 11:05:29}
\Level{C}
\SubArea{Percent problems}
\LANGen{
\Question{
One third of a number \( a \) is by \( 75\% \) greater than one seventh of a number \( b \). Choose the true statement.}
{The number \(a\) is by \(25\%\) smaller than \(b\).}{The number \(a\) is by \(75\%\) smaller than \(b\).}{The number \(a\) is by \(400\%\) greater than \(b\).}{The number \(a\) is by \(75\%\) greater than \(b\).}{}{}}
\LANGpl{
\Question{
Jedna trzecia liczby \( a \) jest o \( 75\% \) większa niż jedna siódma liczby \( b \). Wybierz prawdziwe stwierdzenie.}
{Liczba \(a\) jest o \(25\%\) mniejsza niż \(b\).}{Liczba \(a\) jest o \(75\%\) mniejsza niż \(b\).}{Liczba \(a\) jest o \(400\%\) większa niż \(b\).}{Liczba \(a\) jest o \(75\%\) większa niż \(b\).}{}{}}
\LANGcs{
\Question{
Třetina čísla \( a \) je o \( 75\,\% \) větší než sedmina čísla \( b \). Určete pravdivý výrok.}
{Číslo \(a\) je o \(25\,\%\) menší než \(b\).}{Číslo \(a\) je o \(75\,\%\) menší než \(b\).}{Číslo \(a\) je o \(400\,\%\) větší než \(b\).}{Číslo \(a\) je o \(75\,\%\) větší než \(b\).}{}{}}
\LANGsk{
\Question{
Tretina čísla \( a \) je o \( 75\% \) väčšia ako sedmina čísla \( b \). Určte pravdivý výrok.}
{Číslo \(a\) je o \(25\%\) menšie ako \(b\).}{Číslo \(a\) je o \(75\%\) menšie ako \(b\).}{Číslo \(a\) je o \(400\%\) väčšie ako \(b\).}{Číslo \(a\) je o \(75\%\) väčšie ako \(b\).}{}{}}
\LANGes{
\Question{
Un tercio de un número \(a \) es \(75 \% \) mayor que un séptimo de un número \(b \). Elige la proposición verdadera.}
{El número \(a \) es un \(25 \% \) menor que \(b \).}{El número \(a\) es un \(75\%\) menor que  \(b\).}{El número  \(a\) es un \(400\%\) mayor que  \(b\).}{El número  \(a\) es un \(75\%\) mayor que \(b\).}{}{}}
\End

\Data\ID{2010006403}
\Updated{2022-05-16 11:05:29}
\Level{C}
\SubArea{Percent problems}
\LANGen{
\Question{
Find the number \( x \) if \( 40\% \) of its half is \( 70\% \) of the number \( x -3\).}
{\( 4\frac15 \)}{\( 5\frac14 \)}{\( 5 \)}{\( 5\frac25 \)}{}{}}
\LANGpl{
\Question{
Znajdź liczbę \( x \) jeśli \( 40\% \) jej połowy to \( 70\% \) liczby \( x -3\).}
{\( 4\frac15 \)}{\( 5\frac14 \)}{\( 5 \)}{\( 5\frac25 \)}{}{}}
\LANGcs{
\Question{
Najdi číslo \( x \), pro které platí, že \( 40\,\% \) jeho poloviny se rovná \( 70\,\% \) čísla \( x -3\).}
{\( 4\frac15 \)}{\( 5\frac14 \)}{\( 5 \)}{\( 5\frac25 \)}{}{}}
\LANGsk{
\Question{
Určte číslo \( x \), pre ktoré platí, že \( 40\% \) jeho polovice sa rovná \( 70\% \) čísla  \( x -3\).}
{\( 4\frac15 \)}{\( 5\frac14 \)}{\( 5 \)}{\( 5\frac25 \)}{}{}}
\LANGes{
\Question{
Determina el número \(x \) si el \(40 \% \) de su mitad es \(70 \% \) del número \(x -3 \).}
{\( 4\frac15 \)}{\( 5\frac14 \)}{\( 5 \)}{\( 5\frac25 \)}{}{}}
\End
